% =====================================================================================================================
% IEEE two-column journal paper: lesion-faithful bone suppression and its use for focal lesion detection and classification.
% Sections in sections/*.tex (appendix in sections/appendix.tex), references in sections/references.tex,
% supplementary material in supplementary.tex (compiled on its own).
% tables/ and figures/ hold the generated result tables and figures; never retype a number.
% Compile with pdfLaTeX (Overleaf). Red [TODO: ...] markers are the only places that need the authors' input.
% =====================================================================================================================
\documentclass[journal]{IEEEtran}

\usepackage[T1]{fontenc}
\usepackage[utf8]{inputenc}
\usepackage{amsmath,amssymb,amsfonts}
\usepackage{graphicx}
\usepackage{booktabs}
\usepackage{multirow}
\usepackage{tabularx}
\usepackage{array}
\usepackage{algorithm}
\usepackage{algorithmic}
\usepackage{cite}
\usepackage{url}
\usepackage{xcolor}
\usepackage{siunitx}
\usepackage{adjustbox}
\usepackage{placeins}
\usepackage[hidelinks]{hyperref}

\def\UrlBreaks{\do\/\do\.\do\_\do\-}
\setlength{\emergencystretch}{2em}
% Many floats: let more of them share a page so they stay close to the text that cites them.
\renewcommand{\topfraction}{0.9}\renewcommand{\dbltopfraction}{0.9}\renewcommand{\textfraction}{0.07}
\renewcommand{\floatpagefraction}{0.7}\renewcommand{\dblfloatpagefraction}{0.7}
\setcounter{topnumber}{4}\setcounter{dbltopnumber}{4}\setcounter{totalnumber}{6}

% Red placeholder for text only the authors can supply.
\newcommand{\todo}[1]{\textcolor{red}{[TODO: #1]}}
% Figure that shows a red box instead of stopping the build if the file is missing.
\newcommand{\safefig}[2][\linewidth]{%
  \IfFileExists{#2}%
    {\includegraphics[width=#1]{#2}}%
    {\fbox{\parbox[c][2.4cm][c]{0.92\linewidth}{\centering\todo{missing figure #2}}}}}
\providecommand{\doi}[1]{\url{https://doi.org/#1}}
% Generated result table (tables/<name>.tex), shrunk to fit: \restable{name} in a table, \restable[\textwidth]{name} in a table*.
\newcommand{\restable}[2][\columnwidth]{\begin{adjustbox}{max width=#1}\footnotesize\input{tables/#2.tex}\end{adjustbox}}
% Sub-heading inside a table that stacks two generated tables.
\newcommand{\subtab}[1]{\par\smallskip\noindent{\footnotesize\textit{#1}}\par\smallskip}

\begin{document}

\title{Lesion-Faithful Bone Suppression from Chest CT Projections: Counterfactual Auditing and Gains for Focal Lesion Detection
and Classification in Chest Radiographs}

\author{Tanvir~Hossin\todo{remaining authors}%
\thanks{T.~Hossin is with the Department of Computer Science and Engineering, Daffodil International University, Savar, Dhaka,
Bangladesh. \todo{affiliations of the other authors}}%
\thanks{Corresponding author: \todo{name and e-mail address}.}%
\thanks{This retrospective study used public data sets under their licences, VinDr-RibCXR under its data use agreement, and the
paediatric BDCXR-3257 radiographs under ethical approval REC-FSIT-2025/No:~11928 of Daffodil International University and a data use
agreement with Dr.~M~R~Khan Shishu Hospital \& Institute of Child Health, Dhaka. No new patient data were collected.}}

\markboth{\todo{journal name}}{Hossin \MakeLowercase{\textit{et al.}}: Lesion-Faithful Bone Suppression from Chest CT Projections}

\maketitle

\begin{abstract}
\textbf{Background:} Bone suppression can help readers and detectors find lesions that ribs obscure, but a suppressor trained on CT-derived
images may also erase lesions, and fidelity metrics such as PSNR do not show it.
\textbf{Methods:} We train bone suppressors on public CT projections with an exactly known bone image and add a twin lesion-invariance
regulariser. We measure lesion absorption $A$, the fraction of a counterfactual nodule that the suppressor moves into its bone image, on
real radiographs. We then test whether the suppressed images improve nodule and multi-finding detection and classification against
capacity-matched controls, and examine uncertainty, a paediatric external set and a site calibration of the fidelity metric.
\textbf{Results:} On VinDr-RibCXR the twin ensemble absorbed 0.051 of an added lesion, against 0.149 for a released CT-supervised model
and 0.178 for a matched control, at non-inferior bone-region PSNR (21.46 vs.\ 21.51\,dB); the effect held across three encoders and on
paediatric radiographs ($A=0.084$ vs.\ 0.217 and 0.251). Suppressed images added as extra channels raised Node21 detection (CPM $+0.028$,
95\,\% CI 0.009--0.047) and NIH macro-AUROC ($+0.006$, 0.003--0.008). On NIH the soft-tissue image raised the AUROC of focal findings
($+0.0089$, Holm $p=0.0024$) and lowered that of diffuse findings, and its Grad-CAM++ maps located focal lesions better than those of the
radiograph (interaction $+0.145$, $p=0.0002$); a three-stream network gave no overall gain.
\textbf{Conclusion:} A twin regulariser makes CT-supervised bone suppression lesion-faithful, and suppressed images help focal findings
when added as extra channels, not diffuse disease or overall accuracy as separate streams.
\end{abstract}

\begin{IEEEkeywords}
Bone suppression, chest radiography, counterfactual evaluation, lesion fidelity, nodule detection, multi-channel fusion, uncertainty
quantification.
\end{IEEEkeywords}

\section{Introduction}
\label{sec:intro}
\IEEEPARstart{R}{ibs} and clavicles overlap the lung fields on a chest radiograph and can hide small lesions. Dual-energy subtraction
(DES) separates bone from soft tissue, but it needs dedicated hardware and a second exposure \cite{kuhlman2006des,sun2026bsldm}.
Software bone suppression avoids both: radiologists find more nodules with suppressed images \cite{freedman2011lung}, and suppressed
images have been compared with computer-aided detection \cite{schalekamp2014cad} and used as classifier inputs
\cite{rajaraman2021bone}. Learning-based suppressors need paired supervision. It can come from real DES pairs
\cite{chen2024bsdiff,sun2026bsldm,sun2025gllcm,sun2026debonedit}, which are scarce, or from chest CT, whose bone and soft-tissue
components can be projected separately and are therefore known exactly \cite{angaitkar2026anatomy}. Public projections of CT-RATE
volumes \cite{hamamci2026ctrate} make the second route practical.

A suppressor judged by image fidelity can nevertheless erase lesions. A network rewarded for predicting bone can move lesion signal into
the bone image, so that the lesion disappears from the suppressed radiograph, and PSNR or SSIM computed over the whole image are dominated
by normal anatomy. Because a model that does nothing never absorbs a lesion, absorption has to be read next to suppression strength.
Simulated nodules have been used to measure contrast after hand-built clavicle suppression and in phantoms
\cite{hogeweg2013suppression,takagi2020dose}, and learned image-to-image models are known to remove or invent content in other
modalities \cite{cohen2018distribution,antun2020instabilities,bhadra2021hallucinations,tivnan2024hallucination}. In our search (PubMed,
arXiv and Crossref) we found no audit of learned bone suppression against an exactly known, additive lesion signal, and no training
objective aimed at it.

Whether suppression helps prediction is a separate question. The suppressed images are computed from the radiograph and add no
information that a classifier could not extract itself, so a gain must come from the prior that the suppressor learned from CT or from the
way a network uses the extra images. A network that reads three images spends three times the computation, so a gain can be credited to
suppression only against a control of the same capacity, and seed variation and analytic flexibility can easily produce spurious
differences \cite{varoquaux2022failures,simmons2011falsepositive}. Every comparison here is therefore paired, repeated over seeds and
fixed in advance.

We ask whether CT-supervised bone suppression can preserve pulmonary lesions and improve nodule detection, focal disease classification
and localisation. Our contributions are:
\begin{itemize}
\item A counterfactual lesion audit that measures the fraction $A$ of an exactly known, additive lesion signal moved into the bone image,
 at sites shared by all models and reported beside suppression strength.
\item A twin lesion-invariance regulariser, trained on public CT projections in a three-model ensemble. On VinDr-RibCXR the ensemble
 absorbed 0.051 of an added lesion, against 0.149 for a released CT-supervised model and 0.178 for a matched control, at non-inferior
 bone-region PSNR (21.46 vs.\ 21.51\,dB). The effect replicated across three encoders, on JSRT and on paediatric radiographs
 ($A=0.084$ vs.\ 0.217 and 0.251).
\item Evidence that suppressed images help prediction when added as extra input channels: Node21 detection (CPM $+0.028$,
 95\,\% CI 0.009--0.047), VinBigData detection ($+0.008$, 0.002--0.014) and NIH macro-AUROC ($+0.006$, 0.003--0.008), whereas the
 soft-tissue image alone gave no gain.
\item A fusion ablation of five input configurations (175 runs, three data sets, a capacity-matched control). A three-stream network gave
 no overall gain; the benefit was concentrated in focal findings (NIH soft-tissue focal AUROC $+0.0089$, Holm $p=0.0024$; nodule
 $+0.0129$, Holm $p=0.00121$) and in the localisation of focal lesions (pointing-game interaction $+0.145$, $p=0.0002$), at a cost for
 diffuse findings.
\item Deployment analyses: ensemble uncertainty that needs recalibration on real radiographs, a paediatric external validation and an
 exploratory site calibration of the fidelity metric ($+2.12$\,dB in the bone region without raising lesion absorption).
\end{itemize}
Fig.~\ref{fig:workflow} summarises the two studies. Section~\ref{sec:related} reviews related work, Section~\ref{sec:methods} describes
the data, the audit, the suppressors and the downstream studies, Section~\ref{sec:results} reports the results, and
Section~\ref{sec:discussion} discusses them together with their limitations.

\begin{figure*}[!t]
\centering
\safefig[0.9\textwidth]{figures/fig_tmi_workflow.pdf}
\caption{Overview. Study~I trains CT-supervised bone suppressors with a twin lesion-invariance regulariser, audits their lesion absorption
and tests the suppressed images on detection and classification; Study~II compares five input configurations of a multi-stream
classifier, including a capacity-matched control.}
\label{fig:workflow}
\end{figure*}

\section{Related Work and Motivation}
\label{sec:related}

\subsection{Bone Suppression and CT-Derived Supervision}
Early learning-based rib suppression used massive-training artificial neural networks \cite{suzuki2006mtann}. Later work used
convolutional networks, for example a cascade of multi-scale networks in the gradient domain \cite{yang2017cascade}. Recent methods train
diffusion models on paired radiographs and DES soft-tissue images
\cite{chen2024bsdiff,sun2026bsldm,sun2025gllcm,sun2026debonedit}; the SZCH-X-Rays data set was compiled for BS-LDM \cite{sun2026bsldm}
and made public with DeBoneDiT \cite{sun2026debonedit}. Paired DES images are scarce and come from a few scanners. CT removes this
bottleneck: Angaitkar et al. \cite{angaitkar2026anatomy} build digitally reconstructed radiographs from chest CT, project bone, non-lung
soft tissue and lung separately, and train models that predict one component. Their projections are released and are built from CT-RATE
volumes \cite{hamamci2026ctrate}. We train on these projections and use the released weights as the main comparator. Our models differ in
encoder and resolution (ConvNeXt-V2-T \cite{woo2023convnextv2} at 512~px against SwinV2-B \cite{liu2022swinv2} at 1024~px) and in
sample budget, and they add the twin regulariser, the ensemble and the audit.

\subsection{Lesion Fidelity and Counterfactual Auditing}
Learned image-to-image models can remove or invent content: distribution matching in translation networks can drop or hallucinate lesions
\cite{cohen2018distribution}, deep reconstruction can be unstable to small structures \cite{antun2020instabilities}, and hallucinations
of learned reconstruction have been defined and measured \cite{bhadra2021hallucinations,tivnan2024hallucination}. Inserting known
features or signals into images is an established way to test learned MRI reconstruction \cite{cheng2020addressing} and the effect of
learned denoising on signal detection \cite{li2021assessing}. For bone suppression, simulated nodules were inserted to measure contrast
after a hand-built clavicle suppression method \cite{hogeweg2013suppression}, and a phantom with a simulated nodule at many positions was
used to measure contrast after bone suppression \cite{takagi2020dose}. These studies measure lesion contrast. Contrast alone cannot
separate a suppressor that keeps a lesion from one that suppresses weakly, and a score computed over the whole image is dominated by
normal anatomy. Our audit therefore measures the fraction of a known additive lesion signal that a learned suppressor moves into its bone
image, on real radiographs, at sites shared by every model, and reports it next to suppression strength.

\subsection{Suppressed Images as Inputs, Explanations and Uncertainty}
Suppressed images have been used to classify tuberculosis-consistent findings \cite{rajaraman2021bone}, and readers and computer-aided
detection have been compared with and without them \cite{freedman2011lung,schalekamp2014cad}. A suppressed image can be given to a
classifier alone, stacked with the radiograph as channels, or as a separate input stream; the last two add computation, which a fair
comparison has to match with a control of equal capacity. For multi-label classification we use a class-query decoder over spatial tokens
\cite{liu2021query2label,ridnik2023mldecoder}, so that the attention of each class on each input stream can be read out, and attribution
maps from Grad-CAM++ \cite{chattopadhay2018gradcampp} and Integrated Gradients \cite{sundararajan2017axiomatic} scored with the pointing
game \cite{zhang2018topdown}. For uncertainty we use deep ensembles \cite{lakshminarayanan2017ensembles}, split-conformal prediction
\cite{lei2018conformal,angelopoulos2023conformal}, conformal risk control of the false-negative rate \cite{angelopoulos2022crc},
class-conditional (Mondrian) sets \cite{vovk2005algorithmic} and selective classification \cite{geifman2017selective,geifman2019bias}.

Table~\ref{tab:literature} compares the studies closest to ours. Three gaps motivate the design: lesion erasure by learned suppressors
has not been measured against an exactly known signal; gains from suppressed inputs have not been compared against a capacity-matched
control; and the uncertainty and localisation of CT-supervised suppression have not been examined together with its downstream benefit.

\begin{table*}[!t]
\centering
\caption{Comparison of studies on bone suppression and suppressed images as model inputs with the present work. Check marks indicate
whether a study reports each element: a counterfactual or simulated-lesion audit (Lesion), a capacity-matched control (Cap.),
uncertainty quantification or calibration (UQ), validation on external data (Ext.) and explanation maps (XAI). \checkmark: reported;
(\checkmark): qualitative only; --: not reported; n/a: not applicable to the study design (no model trained by the authors, no input added
to a model, or no clinical task). Rows marked $^{\dagger}$ were assessed from the abstract because the full text is not openly accessible;
all other rows from the full text.}
\label{tab:literature}
\footnotesize
\setlength{\tabcolsep}{3pt}
\begin{tabularx}{\textwidth}{@{}>{\raggedright\arraybackslash}p{2.6cm}>{\raggedright\arraybackslash}p{3.4cm}%
>{\raggedright\arraybackslash}p{3.3cm}>{\raggedright\arraybackslash}X*{5}{c}@{}}
\toprule
Study & Data & Method & Key result & Lesion & Cap. & UQ & Ext. & XAI\\
\midrule
Suzuki et al.$^{\dagger}$ \cite{suzuki2006mtann} & DES bone images for training; 118 validation and 136 independent test radiographs (14
 institutions) & Multiresolution massive-training ANN (MTANN) & Ribs and clavicles suppressed while nodule and vessel visibility was kept
 & -- & n/a & -- & \checkmark & n/a\\
Yang et al.$^{\dagger}$ \cite{yang2017cascade} & 504 two-exposure DES radiographs (404 training, 100 test) & Cascade of multi-scale CNNs
 in the gradient domain (CamsNet) & Soft-tissue PSNR 38.7\,dB, bone suppression ratio 83.8\,\%; other X-ray machines shown visually & --
 & n/a & -- & (\checkmark) & n/a\\
Chen et al. \cite{chen2024bsdiff} & Paired radiographs and DES soft tissue of 120 patients, one hospital (7:2:1 split) & Conditional
 diffusion model (BS-Diff) & Diffusion-based suppression with a clinical image-quality evaluation & -- & n/a & -- & -- & n/a\\
Sun et al. \cite{sun2026bsldm} & SZCH-X-Rays and JSRT (separate splits); Shenzhen set for a downstream test & Conditional latent
 diffusion (BS-LDM) & High-resolution suppression; clinical and downstream lesion-detection evaluation & -- & n/a & -- & \checkmark
 & n/a\\
Sun et al. \cite{sun2025gllcm} & SZCH-X-Rays (592/74/75) and JSRT (192/24/25) & Global-local latent consistency model (GL-LCM) &
 Suppression in 50 sampling steps, 8.54\,s against 84.62\,s for BS-LDM & -- & n/a & -- & -- & n/a\\
Sun et al. \cite{sun2026debonedit} & SZCH-X-Rays (released here) and JSRT; Asraf set for a downstream test & Depth-driven conditional
 bridge diffusion transformer (DeBoneDiT) & Public paired DES data; downstream classification on an external set & -- & n/a & -- &
 \checkmark & n/a\\
Angaitkar et al. \cite{angaitkar2026anatomy} & \num{21887} CT-RATE projections; TBX11K, Node21, VinDr-CXR & Anatomy-decomposed CT
 projections, SwinV2-B at 1024~px & Detection gains and higher real-nodule contrast; no lesion-insertion audit & -- & n/a & -- &
 \checkmark & --\\
Hogeweg et al.$^{\dagger}$ \cite{hogeweg2013suppression} & 20 simulated radiographs; 253 radiographs; 100 with CT nodules simulated
 near the clavicles & Model-based suppression of elongated structures & Simulated nodules kept more contrast than the clavicles after
 suppression (1.35 vs.\ 0.55) & \checkmark & n/a & -- & n/a & n/a\\
Takagi et al.$^{\dagger}$ \cite{takagi2020dose} & Anthropomorphic phantom, 12\,mm simulated nodule at 66 positions & Commercial
 bone-suppression processing at low tube voltage & Nodule contrast and CNR rose away from the ribs; lower dose possible & \checkmark &
 n/a & -- & n/a & n/a\\
Freedman et al.$^{\dagger}$ \cite{freedman2011lung} & Reader study, 368 patients, 15 radiologists & Rib and clavicle suppression software
 & Reader LROC AUC rose from 0.460 to 0.558 & -- & n/a & n/a & n/a & n/a\\
Schalekamp et al.$^{\dagger}$ \cite{schalekamp2014cad} & Reader study, 300 radiographs from four institutions, 8 readers & CAD with
 bone-suppressed images available & CAD raised reader AUC from 0.812 to 0.841 beyond suppressed images & -- & n/a & n/a & n/a & n/a\\
Rajaraman et al. \cite{rajaraman2021bone} & JSRT and NIH-CC DES (suppression); Shenzhen and Montgomery TB sets & Suppressed image
 replacing the radiograph as classifier input (VGG-16) & TB AUC 0.899 to 0.954 (Shenzhen) with suppressed inputs; cross-institutional
 suppression test & -- & n/a & -- & \checkmark & \checkmark\\
\midrule
This work & Qure.ai projections; SZCH-X-Rays, VinDr-RibCXR, JSRT, Node21, TBX11K, VinBigData, NIH ChestX-ray14; BDCXR-3257 & Twin
 lesion-invariance regulariser, $K=3$ ensemble; suppressed images as extra channels; multi-stream classifier with a capacity-matched
 control & $A=0.051$ vs.\ 0.149 (released) and 0.178 (control); Node21 CPM $+0.028$ with extra channels
 & \checkmark & \checkmark & \checkmark & \checkmark & \checkmark\\
\bottomrule
\end{tabularx}
\end{table*}

\section{Materials and Methods}
\label{sec:methods}
This section defines the exact image decomposition, the twin lesion-invariance objective and the counterfactual audit
(Section~\ref{sec:m-formulation}) and describes the data, their preprocessing and splits (Section~\ref{sec:m-data}). It then describes
the suppressors of Study~I and their evaluation (Section~\ref{sec:m-study1}), the downstream models (Section~\ref{sec:m-downstream}), the
fusion study of Study~II (Section~\ref{sec:m-study2}), the external validation (Section~\ref{sec:m-bdcxr}) and the statistical analysis
(Section~\ref{sec:m-stats}). Algorithm~\ref{alg:pipeline} summarises the pipeline and Fig.~\ref{fig:architecture} shows the models; the
rendering of the first-generation (v1) training images from LIDC-IDRI CTs is described in Appendix~\ref{app:v1}. Table~\ref{tab:notation}
lists the notation.

\begin{table}[!t]
\centering
\caption{Notation.}
\label{tab:notation}
\footnotesize
\renewcommand{\arraystretch}{1.08}
\setlength{\tabcolsep}{3pt}
\begin{tabular}{@{}>{\raggedright\arraybackslash}p{0.25\linewidth}>{\raggedright\arraybackslash}p{0.71\linewidth}@{}}
\toprule
Symbol & Meaning \\
\midrule
$X$, $S$, $B$ & radiograph, soft-tissue and bone image of a training sample, $S=X-B$ \\
$x$, $b$ & suppressor input and its exact bone target \\
$f$, $f_k$ & suppressor (predicts the bone image); ensemble member $k$ \\
$\bar f$, $\sigma$ & ensemble mean and per-pixel standard deviation ($K=3$) \\
$\hat s$, $\hat B$ & soft-tissue estimate $x-\bar f(x)$; bone image $X-S$ given to the classifier \\
$n$, $\alpha$ & inserted lesion signal and its amplitude \\
$x^{+}$, $M$ & twin $\min(x+n,1)$; lesion support dilated by 3\,px \\
$w$ & twin weight ($w=16$ in the final model) \\
$A$ & lesion absorption, Eq.~\eqref{eq:absorption} \\
$q$, $\mathcal{C}_p$ & conformal scale and interval at pixel $p$ \\
$s$, PSNR$_{\mathrm{cal}}$ & dual-energy soft-tissue reference; calibrated PSNR \\
$R_{\mathrm{rib}}$, $R_{\mathrm{non}}$ & gradient-energy reduction on rib-boundary and other lung pixels \\
$\Delta_G$, $\rho_{c,s}$ & mean AUROC difference of a finding group $G$; routing of class $c$ to stream $s$ \\
$I$ & pointing-game interaction, focal minus other targets \\
\bottomrule
\end{tabular}
\end{table}

\begin{figure*}[!t]
\centering
\safefig[0.9\textwidth]{figures/fig_tmi_architecture.pdf}
\caption{Model architectures. (a) Study~I: bone suppressor trained with the twin lesion-invariance regulariser
(Section~\ref{sec:m-twin}); the twin $x^{+}=\min(x+n,1)$ shares the weights and the bone target of $x$. (b) Study~II: multi-stream
classifier with a shared encoder and a class-query decoder (Section~\ref{sec:m-study2}), and the five input configurations.}
\label{fig:architecture}
\end{figure*}

\begin{algorithm}[!t]
\caption{Training, counterfactual audit and downstream evaluation}
\label{alg:pipeline}
\footnotesize
\begin{algorithmic}[1]
\REQUIRE Qure.ai projections $\mathcal{D}_{\mathrm{qp}}$ with exact $(X,B)$ (Eq.~\eqref{eq:qure}); lesion libraries
 $\mathcal{L}_{\mathrm{train}}$, $\mathcal{L}_{\mathrm{test}}$ from LIDC-IDRI training and test CTs; real test sets
\ENSURE Suppressor ensemble $\bar f$ with intervals $\mathcal{C}_p$; downstream and fusion classifiers; all estimates with 95\,\% intervals
\FOR[\textit{Study~I}]{seed $k\in\{0,1,2\}$}
 \FOR{step $=1$ \TO \num{37500}}
  \STATE Compose a batch: re-weight $B$, shared affine, insertion, tone curve $\gamma$, $B'=X'-S'$ (Eq.~\eqref{eq:tone})
  \STATE Build twins $x^{+}=\min(x+n,1)$, $n$ from $\mathcal{L}_{\mathrm{train}}$, for half of the batch
  \STATE Update $f_k$ with AdamW on $\mathcal{L}$ of Eq.~\eqref{eq:twin}; update EMA weights
 \ENDFOR
\ENDFOR
\STATE $\bar f,\sigma\gets$ Eq.~\eqref{eq:ensemble}; calibrate $q$ on Qure.ai validation pixels (Eq.~\eqref{eq:conformal})
\FOR[\textit{audit}]{each real test image $x$}
 \STATE Draw six lung sites with a fixed seed and lesions from $\mathcal{L}_{\mathrm{test}}$
 \STATE For every model, compute $A$ (Eq.~\eqref{eq:absorption}) and bone overlap from independent masks
\ENDFOR
\STATE Score fidelity, rib benchmark, $\Delta$CNR, detection and classification; test H1$'$--H5$'$ and E1--E3
\FOR[\textit{Study~II}]{config $\in\{$orig, soft, early, hybrid, ctrl$\}$, data set, fold, seed $0..4$}
 \STATE Train the class-query classifier on the $(X,S,\hat B)$ streams (Eqs.~\eqref{eq:tokens}--\eqref{eq:bce})
\ENDFOR
\STATE Seed ensembles (Eq.~\eqref{eq:seedens}); cluster-bootstrap contrasts with Holm; test H6--H10; pointing game, conformal sets, AURC
\RETURN all estimates with 95\,\% intervals, whatever their direction
\end{algorithmic}
\end{algorithm}


% ---------------------------------------------------------------------------------------------------------------------
\subsection{Problem Formulation and Counterfactual Audit Definitions}
\label{sec:m-formulation}

\subsubsection{Training images with an exactly known bone image}
\label{sec:m-qure}
A suppressor $f$ predicts the bone image of a radiograph; the suppressed (soft-tissue) image is the radiograph minus that prediction.
Supervision needs training images whose bone image is known exactly. The final (v2) suppressors are trained on public projections of chest
CT \cite{angaitkar2026anatomy}: each row of the Qure.ai data holds a composite radiograph $T$ and the translated bone, soft-tissue and lung
projections $B_t$, $S_t$ and $L$ (8-bit, $512\times512$, one sample per CT volume at 100~kVp). We verified that
\begin{equation}
\begin{gathered}
T=a\,S_t+\beta\,B_t+c\,L+d,\\
X=\frac{T}{255},\quad B=\frac{\beta\,B_t}{255},\quad S=X-B
\end{gathered}
\label{eq:qure}
\end{equation}
holds up to rounding, with a per-sample bone weight $\beta$ fitted by least squares, so $B$ is the exactly additive bone image of the
normalised input $X$. Over all \num{21887} rows the least-squares fit leaves a mean absolute residual of 0.36 grey levels per projection
(largest 3.2 of 255).

A training sample is composed on the GPU in this order. The bone image is re-weighted ($\tilde B=w_bB$, $w_b$ log-uniform in $[0.6,1.7]$,
so $\tilde X=S+\tilde B$ stays exact); one random affine transform is applied to all maps; in the arms with counterfactual insertion
(Table~\ref{tab:arms}), up to three lesions are inserted into $S$ with the bone target untouched; and a monotone gamma tone curve $\gamma$
gives
\begin{equation}
X'=\gamma(\tilde X),\qquad S'=\gamma(S),\qquad B'=X'-S',
\label{eq:tone}
\end{equation}
which keeps the decomposition exact under any monotone mapping. Linear blur is applied to $X'$ and $B'$, and noise and scatter to $X'$
only. The first-generation (v1) suppressors were trained on radiographs rendered from LIDC-IDRI CTs with polychromatic attenuation physics
(Appendix~\ref{app:v1}).

\subsubsection{Suppression loss}
For a predicted bone image $\hat b=f(x)$ and the target $b$, the suppression loss is
\begin{equation}
\begin{split}
\mathcal{L}_{\mathrm{sup}}(\hat b,b)={}&0.5\,\lVert\hat b-b\rVert_1+0.5\,\lVert\hat b-b\rVert_2^2\\
&+0.5\,\big(1-\mathrm{SSIM}(\hat b,b)\big)\\
&+2.5\,\big(\lVert\hat b-b\rVert_1+\lVert\hat b-b\rVert_2^2\big)_{\mathrm{ROI}},
\end{split}
\label{eq:loss}
\end{equation}
with norms taken as means over pixels and the region of interest (ROI) taken from the lung component. The suppressed image is
$x-\hat b$.

\subsubsection{Twin lesion-invariance objective}
\label{sec:m-twin}
Let $f$ be the suppressor, $x$ a training image and $b$ its exact bone image. A \emph{twin} of $x$ is $x^{+}=\min(x+n,1)$, where $n\ge0$ is
the signal of a few lesion patches taken from the lesion library of \emph{training} CTs and pasted at random positions in the lung fields,
in composite units and after all augmentation. The bone target of $x^{+}$ is the same $b$. Let $M$ be the lesion support dilated by
$3$~px. The loss is
\begin{equation}
\begin{gathered}
\mathcal{L}=\tfrac12\big[\mathcal{L}_{\mathrm{sup}}(f(x),b)+\mathcal{L}_{\mathrm{sup}}(f(x^{+}),b)\big]+w\,\mathcal{L}_{\mathrm{inv}},\\
\mathcal{L}_{\mathrm{inv}}=\frac{\sum_p M_p\,\big|f(x^{+})_p-f(x)_p\big|}{\max\big(1,\sum_p M_p\big)},
\end{gathered}
\label{eq:twin}
\end{equation}
where $w$ is the twin weight. Half of every batch consists of twins, so the cost per step does not change. The supervised term already asks
for the unchanged bone image on $x^{+}$; $\mathcal{L}_{\mathrm{inv}}$ ties the two predictions directly on the lesion support. Each twin
carries one or two pasted lesions (uniform), each with a random scale and amplitude; the supervised term is averaged over all base images
and twins of the batch, and $\mathcal{L}_{\mathrm{inv}}$ is normalised per image by its support and then averaged over the twins.
Fig.~\ref{fig:insertion} shows one pasted lesion, its support $M$ and the unchanged bone target.

\begin{figure}[!t]
\centering
\safefig[\linewidth]{figures/fig_lesion_insertion.pdf}
\caption{Counterfactual lesion insertion, illustrated on a rendered LIDC-IDRI test CT; in training, the twins of the regulariser are built
in the same way from Qure.ai projections with lesions of the training library. (a) Lesion signal $n$ of a 15~mm LIDC-IDRI nodule; (b) the
composite with the lesion, $x+n$, and the dilated support $M$; (c) the bone target under the site, which the insertion leaves unchanged.}
\label{fig:insertion}
\end{figure}

\subsubsection{Counterfactual lesion absorption}
\label{sec:m-absorption}
For a suppressor $f$ that outputs the bone image, an image $x$ and a lesion signal $n\ge0$ in composite-image units, the absorption at
amplitude $\alpha$ is
\begin{equation}
A=\frac{\big\langle f(x+\alpha n)-f(x),\;\alpha n\big\rangle}{\big\langle \alpha n,\;\alpha n\big\rangle},
\label{eq:absorption}
\end{equation}
with both inner products taken over the lesion support (the pixels where the inserted signal exceeds 2\,\% of its peak). $A=0$ means that
the lesion stays in the soft-tissue image and $A=1$ that it was removed with the bone; $A$ is a regression slope and can fall outside
$[0,1]$. A model that does nothing has $A=0$, so $A$ is reported together with a suppression gain and never alone. Two suppression
measures are paired with $A$: the bone-region PSNR on the held-out SZCH pairs and, on VinDr-RibCXR, the reduction of gradient energy on
rib-boundary pixels (the rib benchmark, Eq.~\eqref{eq:rib}).

% ---------------------------------------------------------------------------------------------------------------------
\subsection{Data Sets, Preprocessing and Splits}
\label{sec:m-data}
Ten data sets serve three roles in this work: sources for training the suppressors, real radiographs for evaluating them and the
downstream models, and an external paediatric set for testing transfer to an unseen population. Table~\ref{tab:data} lists each set with
its role, splits and licence or access condition, Fig.~\ref{fig:data_splits} shows the size of every split, and
Fig.~\ref{fig:data_pipeline} with Table~\ref{tab:preproc} follows each image source from its raw form to the study that uses it
(Section~\ref{sec:m-preproc}).

\subsubsection{Training sources}
The LIDC-IDRI collection \cite{armato2011lidc,armato2015lidctcia,clark2013tcia}, with some series read from the LUNA16 copy
\cite{setio2017luna16}, supplies the v1 renders, the lesion library used to paste lesions into training images, and the synthetic test CTs.
We used 1\,018 chest CT series from 1\,010 patients with the radiologists' nodule contours. Series were read from DICOM or from the LUNA16
MetaImage copy (voxel-identical on all 12 series we compared, including nodule slice locations); 1\,000 series could be rendered (18 were
skipped because of non-axial orientation or irregular slice spacing). A render was \emph{usable} if the lungs were not cut by the scan, at
most 50 of the 512 frame rows lay beyond the scan, the lung volume was between 1.5 and 9~L and the bone volume exceeded 300~cm$^3$ (844
usable CTs). Patients were assigned to train/validation/test (70/10/20\,\%) by a hash of the patient identifier (588/98/158 CTs;
1\,057/172/298 consensus nodules) before any model was trained.

The v2 models are trained on the public Qure.ai \texttt{ct2xr-projections} data \cite{angaitkar2026anatomy}, which hold, for each CT
volume, a composite image and translated bone, soft-tissue and lung projections (Section~\ref{sec:m-qure}). The release holds \num{21887}
projections, one per CT volume; the volume names are those of CT-RATE \cite{hamamci2026ctrate}. We split them by patient (\num{14374}
patients) into \num{21051} training, 406 validation and 430 test projections.

\subsubsection{Real radiographs}
\emph{SZCH-X-Rays} holds paired radiographs and dual-energy soft-tissue images from one partner hospital \cite{sun2026bsldm,sun2026debonedit}.
We use the 75-pair test subset; its 75 images are split into a development part (24 images) and a held-out part (51 images) by a hash of
the file name. The 75 images are the published test split of the 741-pair version of the data set (592/74/75), acquired with dual-exposure
DES on a GE Discovery XR656 \cite{sun2026debonedit}; an earlier version of 818 pairs was used for BS-LDM \cite{sun2026bsldm}, and the
papers do not explain the difference.

\emph{VinDr-RibCXR} holds 245 radiographs with per-rib masks \cite{nguyen2021ribcxr}. The data are licensed and are processed only on the
licence holder's own computer and private cloud account; only metrics leave them.

\emph{Nodule and detection sets.} JSRT \cite{shiraishi2000jsrt} (154 nodules), Node21 \cite{sogancioglu2024node21} (1\,476 nodule boxes, of
which two narrower than 4\,px at 512\,px have no defined contrast, so the contrast analysis uses 1\,474), TBX11K \cite{liu2020tbx11k} and
VinBigData (the Kaggle competition data, related to VinDr-CXR \cite{nguyen2022vindrcxr}) are used for real-lesion contrast and for
downstream detection. The detection sets are split by a fixed hash of the image identifier: Node21 \num{3440}/\num{512}/\num{930}, TBX11K
\num{6200}/\num{688}/\num{2088} (its official validation set serves as test set) and VinBigData \num{10433}/\num{1469}/\num{2943} images
for training, validation and test. VinBigData boxes are kept where at least two radiologists marked the same finding; images with only
single-reader findings are left out. No image of any set was unreadable. VinBigData is the Kaggle competition built on the VinDr-CXR images
\cite{vinbigdata2020kaggle}; its rules allow non-commercial research use.

\emph{NIH ChestX-ray14} \cite{wang2017chestxray} is used for classification with its official test list (\num{25596} images) and a
patient-level split of the remaining images into \num{77890} training and \num{8634} validation images. The downstream models of Study~I
and the fusion study of Study~II use the same NIH and VinBigData splits; on VinBigData an image is positive for one of the 14 findings when
at least two of its readers marked it. The 984 expert boxes on 880 NIH test images (eight classes) are used only to score attribution maps.

\subsubsection{External validation set and near-duplicate screening}
\emph{BDCXR-3257} holds 3,257 paediatric chest radiographs (JPEG) from the Radiology and Imaging Department of Dr.~M~R~Khan Shishu Hospital
\& Institute of Child Health, Dhaka, collected in routine care and anonymised, with one label each: normal (880), bacterial pneumonia
(1,477) or viral pneumonia (900) \cite{durjoy2026bdcxr}, a task studied before on other paediatric sets \cite{kermany2018identifying}. Most
children are under five years old (1,132 images under one year). The data set is public on figshare under CC BY 4.0, but we obtained it
under the ethical approval REC-FSIT-2025/No:~11928 (Daffodil International University) and a data use agreement between the hospital and
the authors' research group, whose terms we follow: no radiograph of this set is shown, and no image or model derived from it is
redistributed by us. There are no patient identifiers, so near-duplicate images are grouped by their correlation at $64$~px
($r\ge0.93$; 48 pairs, the rest single images), and a group never spans two folds.

\subsubsection{Development and locked sets}
The locked test sets were fixed before any locked-set evaluation. After the ablation of the v1 suppressors the SZCH pairs had been looked
at, so SZCH became a development set for decisions about the final recipe, which was fixed before any held-out or locked evaluation of a
final model. Decisions used the 24 development images, the synthetic validation data and the real-image audit on all 75 SZCH images; the
audit on SZCH is therefore development evidence only, and the confirmatory audit runs on VinDr-RibCXR (not opened by any model run before
the final protocol was written) and JSRT. Paired fidelity claims use the 51 held-out SZCH pairs. While the recipe was being chosen,
screening tools printed held-out PSNR and SSIM next to the development numbers for v1-family and screening runs; no rule, table or decision
used them, and the tool was changed to print development numbers only.

\subsubsection{Preprocessing steps}
\label{sec:m-preproc}
Every image passes through a fixed chain before any model sees it (Fig.~\ref{fig:data_pipeline}, Table~\ref{tab:preproc}):
\begin{enumerate}
\item \emph{Decomposition of CT into exact components.} LIDC-IDRI CTs are split into bone, non-lung soft tissue and lung by Hounsfield
 rules, projected along the anteroposterior axis and rendered with polychromatic physics (Appendix~\ref{app:v1}); the Qure.ai projections
 are converted into an exactly additive radiograph and bone image (Eq.~\eqref{eq:qure}).
\item \emph{Augmentation that keeps the decomposition exact.} Bone re-weighting, one shared affine transform, lesion insertion and a
 monotone tone curve are applied so that $B'=X'-S'$ holds for every training sample (Eq.~\eqref{eq:tone}).
\item \emph{Real radiographs.} Each radiograph is resized to 1024\,px by area averaging and min--max scaled; the suppressors read it at
 512\,px, the classifiers at 256\,px and the detector at 512\,px.
\item \emph{Lung masks and audit sites.} Lung fields come from a published segmentation network, eroded by 49\,px, and six lesion sites
 per image are drawn once with a fixed seed, so that every model is audited at the same places.
\item \emph{Integrity screening and splits.} Unusable renders are removed by fixed rules, near-duplicate paediatric images are grouped, and
 every split is made by patient, image hash or duplicate group before any model sees the corresponding test data.
\end{enumerate}

\begin{figure*}[!t]
\centering
\safefig[0.88\textwidth]{figures/fig_data_pipeline_tmi.pdf}
\caption{Data preprocessing. Each row follows one image source from its raw form (left) through preprocessing and splitting (train /
validation / test; dev: development pairs, test: held-out pairs) to the studies that use it (right; I, II). Thumbnails are real images
where the licence allows: a LIDC-IDRI CT slice \cite{armato2015lidctcia} (CC~BY~3.0), a SZCH-X-Rays radiograph with its dual-energy
soft-tissue image \cite{sun2026bsldm} (CC~BY-NC~4.0) and an NIH ChestX-ray14 radiograph \cite{wang2017chestxray}. $\dagger$: no image
derived from CT-RATE is shown, so the Qure.ai row shows the composite radiograph $X$ and the exact bone image $B$ of a rendered LIDC-IDRI
CT. $\ddagger$: dummy radiographs, variants (contrast, noise, mirroring, framing) of the same rendered LIDC-IDRI CT, stand in for data sets
whose images are not shown (data use agreements or restrictive licences: VinDr-RibCXR, JSRT, VinBigData, BDCXR-3257, Node21; TBX11K is
not shown either); they are not images of these data sets.}
\label{fig:data_pipeline}
\end{figure*}

\begin{figure}[!t]
\centering
\safefig[\columnwidth]{figures/fig_data_splits_tmi.pdf}
\caption{Size of every data split on a logarithmic scale. SZCH-X-Rays has only development and held-out pairs, VinDr-RibCXR and JSRT are
test sets only, and BDCXR-3257 is used in five-fold cross-validation over near-duplicate groups.}
\label{fig:data_splits}
\end{figure}

\begin{table*}[!t]
\centering
\caption{Data sets, their role, splits (train / validation / test unless stated) and licence or access condition, as stated on the
providers' pages (read on 3 November 2026).}
\label{tab:data}
\footnotesize
\renewcommand{\arraystretch}{1.08}
\begin{tabularx}{\textwidth}{@{}>{\raggedright\arraybackslash}p{3.1cm}>{\raggedright\arraybackslash}p{4.4cm}%
>{\raggedright\arraybackslash}p{4.2cm}>{\raggedright\arraybackslash}X@{}}
\toprule
Data set & Role & Split & Licence or access \\
\midrule
LIDC-IDRI (TCIA) & v1 renders, lesion library, synthetic test CTs & 588 / 98 / 158 CTs (by patient) & CC BY 3.0; attribution with the data DOI \\
Qure.ai ct2xr-projections & v2 training source (one projection per CT-RATE volume) & \num{21051} / 406 / 430 projections (by patient) & CC BY-NC-SA 4.0 \\
SZCH-X-Rays & paired DES evaluation, site calibration & 24 development / 51 held-out pairs & CC BY-NC 4.0 \\
VinDr-RibCXR & locked audit, rib benchmark & 245 test images & licensed; local use, metrics only \\
JSRT & audit, nodule contrast & 247 test images (154 nodules) & free after registration; no redistribution or commercial use \\
Node21 & nodule contrast, detection & \num{3440} / 512 / 930 images & CC BY-NC-ND 4.0 \\
TBX11K & detection & \num{6200} / 688 / \num{2088} images & CC BY 4.0 (data) \\
VinBigData / VinDr-CXR & detection; fusion study & \num{10433} / \num{1469} / \num{2943} images & competition rules: non-commercial research use \\
NIH ChestX-ray14 & downstream classification; fusion study & \num{77890} / \num{8634} / \num{25596} images (test: official list) & unrestricted use with attribution \\
BDCXR-3257 & external validation; fusion study & five folds over near-duplicate groups (\num{3257} images) & CC BY 4.0; used under the providers' data use agreement \\
\bottomrule
\end{tabularx}
\end{table*}

\begin{table*}[!t]
\centering
\caption{Preprocessing of each image source. All splits are by patient, by a fixed hash of the image identifier or, for BDCXR-3257, by
near-duplicate group, and were fixed before any model saw the corresponding test data.}
\label{tab:preproc}
\footnotesize
\renewcommand{\arraystretch}{1.1}
\setlength{\tabcolsep}{3pt}
\begin{tabularx}{\textwidth}{@{}>{\raggedright\arraybackslash}p{3.0cm}>{\raggedright\arraybackslash}p{3.0cm}%
>{\raggedright\arraybackslash}X>{\raggedright\arraybackslash}p{3.4cm}@{}}
\toprule
Source & Used by & Preprocessing & Integrity checks \\
\midrule
LIDC-IDRI CT & v1 suppressors, lesion library & Bone, non-lung soft tissue and lung masks; $\rho=\mathrm{clip}(1+\mathrm{HU}/1000,0,3)$; AP
 projection of each component; body crop, $512\times512$; physics of Eq.~\eqref{eq:physics} at 80--130~kVp (training) or 100~kVp
 (evaluation); composer of Eq.~\eqref{eq:compose} & 18 of 1\,018 series not renderable; usability rules leave 844 CTs; DICOM and LUNA16
 copies voxel-identical on 12 series\\
Qure.ai projections & v2 suppressors & 8-bit, $512\times512$; $X=T/255$, exact $B$ (Eq.~\eqref{eq:qure}); bone weight log-uniform in
 $[0.6,1.7]$, shared affine, insertion, gamma tone curve, blur on $X'$ and $B'$, noise and scatter on $X'$ & Mean absolute residual of the
 additive fit 0.36 grey levels (largest 3.2 of 255)\\
Real radiographs (suppressor input) & Studies I and II & Suppressor run at 512\,px; lung masks from the TorchXRayVision PSPNet at
 1024\,px, eroded by 49\,px for audit sites & Masks exported once and shared by all models\\
Classifier inputs & Study~II, downstream classification, BDCXR-3257 & $256\times256$\,px; per-image standardisation and a learned affine
 map per stream; rotation $\le7^\circ$, scale 0.9--1.1, translation $\le3$\,\%, brightness and contrast, identical across channels &
 VinBigData labels need two readers\\
Detector inputs & Downstream detection & $512$\,px; boxes of findings marked by two or more readers & Node21: two boxes narrower than 4\,px
 excluded from contrast\\
BDCXR-3257 & External validation & JPEG; five class-stratified folds over near-duplicate groups (fold seed 0) & Correlation at 64\,px,
 $r\ge0.93$: 48 pairs grouped\\
\bottomrule
\end{tabularx}
\end{table*}

% ---------------------------------------------------------------------------------------------------------------------
\subsection{Study~I: Lesion-Faithful Bone Suppression and Deep Ensemble Uncertainty}
\label{sec:m-study1}
Study~I asks whether a CT-supervised suppressor can remove bone without removing lesions, and whether a twin regulariser can teach it to.
This subsection describes the suppressors, their training arms and comparators, the ensemble and its intervals, the audit, and the
measures used to score them on real radiographs.

\subsubsection{Suppressor architecture and training}
The suppressor (Fig.~\ref{fig:architecture}a) is a nested (UNet++) decoder \cite{zhou2018unetpp} over the four stages of a pretrained
encoder, with CBAM attention \cite{woo2018cbam}, a plain U-Net tail with two shallow convolutional skips to full resolution and a
$1\times1$ head that outputs the bone image directly. The v1 suppressors use a SwinV2-Tiny encoder \cite{liu2022swinv2} (ImageNet-1k,
window 8, $512^2$ input); relative to the reference this is a smaller backbone at half the resolution, forced by compute. They were trained
with AdamW (learning rate $2\times10^{-4}$, weight decay 0.01, 500-step warm-up, cosine decay), batch 8, mixed precision, gradient clipping
at 1, EMA weights and \num{200000} samples per arm. The dose-response and ablation arms of v2 use the same SwinV2-T encoder; our final model
uses a ConvNeXt-V2-T encoder \cite{woo2023convnextv2} (FCMAE pre-training, ImageNet-22k then ImageNet-1k fine-tuning
\cite{deng2009imagenet}; \texttt{timm} \cite{wightman2019timm} model \texttt{convnextv2\_tiny.fcmae\_ft\_in22k\_in1k}) with a learning rate
$0.25\times$ that of the decoder. Every v2 run is trained at $512$~px with AdamW \cite{loshchilov2019adamw} (learning rate
$2\times10^{-4}$, batch 8), the schedule and loss above, and \num{300000} samples (\num{37500} steps); evaluation uses the EMA weights. The
encoder screen uses \num{100000} samples. Every arm changes exactly one factor of its parent (Table~\ref{tab:arms}); all other settings,
including the seed set, are identical. Table~\ref{tab:settings} collects the training settings of every model in this work.

\begin{table*}[!t]
\centering
\caption{Training settings of all models. Learning rates are those of AdamW unless stated; EMA: exponential moving average of the
weights.}
\label{tab:settings}
\footnotesize
\renewcommand{\arraystretch}{1.08}
\setlength{\tabcolsep}{3pt}
\begin{tabularx}{\textwidth}{@{}>{\raggedright\arraybackslash}p{2.6cm}>{\raggedright\arraybackslash}X%
>{\raggedright\arraybackslash}p{4.4cm}>{\raggedright\arraybackslash}p{2.5cm}>{\raggedright\arraybackslash}p{2.7cm}@{}}
\toprule
Model & Architecture and input & Optimisation & Budget and seeds & Selection \\
\midrule
v1 suppressors & SwinV2-T encoder, UNet++/CBAM decoder; LIDC-IDRI renders at 512\,px & $2\times10^{-4}$, weight decay 0.01, 500-step
 warm-up, cosine; batch 8; mixed precision, clipping at 1; EMA & \num{200000} samples; 3 seeds per arm & validation CTs \\
Final suppressors (ours) & ConvNeXt-V2-T encoder (learning rate $0.25\times$), UNet++/CBAM decoder; Qure.ai projections at 512\,px; twin
 weight $w=16$ & $2\times10^{-4}$, batch 8, same schedule; EMA & \num{300000} samples (\num{37500} steps); seeds 0--2, seed 3 reserve &
 development data, rule fixed in advance \\
Dose-response, ablation & SwinV2-T encoder, otherwise as ours & as ours & \num{300000} samples; 2--4 seeds per arm & -- \\
Encoder screen & ConvNeXt V2, DINOv3, RAD-DINO (learning rate $0.25\times$) & as ours & \num{100000} samples & screening rule \\
Downstream classifier & DenseNet-121, 256\,px & $3\times10^{-4}$, weight decay $10^{-4}$, batch 64, cosine & 10 epochs; 3 seeds & best
 validation macro-AUROC \\
Downstream detector & Faster R-CNN, ResNet-50 FPN, 512\,px & SGD 0.01, momentum 0.9, weight decay $10^{-4}$, batch 8 & 12 epochs; 3
 seeds & best validation CPM \\
Fusion classifier & shared DenseNet-121, class-query decoder (2 layers, 8 heads, $D=256$), 256\,px & $3\times10^{-4}$, weight decay
 $10^{-4}$, batch 32, 200 warm-up steps, cosine; mixed precision & 20 epochs; seeds 0--4 (each BDCXR-3257 fold) & best validation
 macro-AUROC \\
BDCXR-3257 classifier (E1) & as the downstream classifier & as the downstream classifier & 10 epochs; 3 seeds per fold & 10\,\% of each
 training part, by group \\
\bottomrule
\end{tabularx}
\end{table*}

\subsubsection{Training arms}
Table~\ref{tab:arms} lists the arms. The dose-response varies the twin weight over $w\in\{0,4,8,16,32,64\}$ in the SwinV2-T family with 3,
2, 2, 4, 2 and 2 seeds ($w=1$ appears only in the screening). The ablation changes one part of the regulariser at a time: a stop-gradient
on $f(x)$; a support dilated further with a $31\times31$ window, or the whole image, instead of $M$; up to four instead of up to two pasted
lesions; and lesion insertion into the base image without the tie term. The insertion-only arm and the last four rows of
Table~\ref{tab:arms} form this ablation (two seeds each, SwinV2-T, $w=16$, development data).

\begin{table}[!t]
\centering
\caption{Training arms. ``Lesions'' is the maximum number of lesions inserted into the base image (before the physics for the LIDC
renders, before the tone curve for the Qure.ai projections); the twin arms paste 1--2 further lesions per twin (up to 4 in
\texttt{qp\_twin16\_n4}) into the finished composite. ``Twin'' is the twin weight $w$, varied over $\{4,8,16,32,64\}$ in the
\texttt{qp\_twin}$w$ arms ($w=1$ in the screening only). The parent of an arm is \texttt{recipe} for the LIDC arms, \texttt{qp} for the
first Qure.ai arms and \texttt{qp\_twin16} for the last four rows.}
\label{tab:arms}
\footnotesize
\setlength{\tabcolsep}{3pt}
\begin{tabular}{@{}lccc>{\raggedright\arraybackslash}p{0.36\linewidth}@{}}
\toprule
Arm & Source & Lesions & Twin & Change \\
\midrule
\texttt{recipe} & LIDC & 0 & none & v1 baseline \\
\texttt{nodule\_loss} & LIDC & 0 & none & lesion-weighted loss \\
\texttt{clia} & LIDC & 3 & none & counterfactual insertion \\
\texttt{clia\_nodule} & LIDC & 3 & none & insertion and lesion-weighted loss \\
\texttt{qp} & Qure & 0 & none & Qure.ai projections (no-twin control) \\
\texttt{qp\_clia} & Qure & 3 & none & insertion only, no tie term \\
\texttt{qp\_twin}$w$ & Qure & 3 & varied & twin regulariser \\
\texttt{qp\_twin16\_sg} & Qure & 3 & 16 & stop-gradient on the un-pasted branch \\
\texttt{qp\_twin16\_wide} & Qure & 3 & 16 & support dilated with a $31\times31$ window \\
\texttt{qp\_twin16\_all} & Qure & 3 & 16 & whole-image support \\
\texttt{qp\_twin16\_n4} & Qure & 3 & 16 & four lesions per twin \\
\bottomrule
\end{tabular}
\end{table}

\subsubsection{Encoder screen and model selection}
\label{sec:m-encoders}
To check whether a newer encoder helps under the same recipe, ConvNeXt V2 \cite{woo2023convnextv2}, DINOv3 \cite{simeoni2025dinov3} and
RAD-DINO \cite{perezgarcia2025raddino} were trained at a reduced budget, and an encoder moved to the full budget only if it met a rule
fixed in advance: at \num{100000} samples, absorption $A\le0.070$ on the development audit with bone-region PSNR$_{\mathrm{cal}}\ge21.60$~dB
on the development pairs, or PSNR$_{\mathrm{cal}}\ge21.95$~dB with $A\le0.100$ (Supplementary Table~S7). The thresholds come from the
SwinV2-T range at that budget and were written after one SwinV2-B run had been seen. Pretrained encoders other than SwinV2 are trained with
a learning rate $0.25\times$ that of the decoder; ConvNeXt-V2-T diverged at the full rate (fp16 overflow) and was rerun at $0.25\times$, an
extension of the rule written for vision transformers. DINOv3 ViT-S/16 did not improve on SwinV2-T, RAD-DINO ViT-B/14 missed the rule on
PSNR by 0.05~dB but was continued at the full budget as an exploratory encoder, and ConvNeXt-V2-T met it.

\emph{Ours} is the ensemble of three ConvNeXt-V2-T UNet++/CBAM suppressors (seeds 0--2) trained on the Qure.ai projections with twin weight
$w=16$. It replaced the planned SwinV2-T configuration under a replacement rule declared before the final runs; the screening rule only
decided which encoders were trained at the full budget, and the replacement rule applies to every full-budget candidate. At two seeds
(development data) RAD-DINO reached $A=0.046$ and bone-region PSNR$_{\mathrm{cal}}=21.92$\,dB and ConvNeXt-V2-T $A=0.050$ and $22.01$\,dB;
both met the replacement rule, and the tie-break declared in advance (mean absorption within 0.005, then the higher bone-region PSNR) chose
ConvNeXt-V2-T, the only encoder that passed both rules, so excluding RAD-DINO would not have changed the choice. With the third seeds the
absorption gap is 0.0055, just outside the band; the choice was not reopened, and both encoders are reported. The SwinV2-T and RAD-DINO
pairs (twin and no twin, three seeds each) are secondary.

\subsubsection{Comparators}
Comparators are the \emph{no-twin control} (same encoder, data and recipe, $w=0$, three seeds); the v1 model trained on LIDC-IDRI only
(\texttt{clia}, the planned pick of the v1 study); the publicly released weights of Angaitkar et al. \cite{angaitkar2026anatomy} (SwinV2-B
at $1024$~px; research licence), called the \emph{reference model}; and three controls that need no training: identity, CLAHE
\cite{zuiderveld1994clahe} and a homomorphic filter \cite{oppenheim1968homomorphic}. All methods are scored with the same code on the same
images. BS-LDM \cite{sun2026bsldm}, trained on dual-energy pairs of the SZCH hospital, is shown for orientation only.

\subsubsection{Deep ensemble, split-conformal intervals and recalibration}
\label{sec:m-ensemble}
The reported model is a deep ensemble of $K=3$ suppressors trained with seeds $0$, $1$ and $2$ \cite{lakshminarayanan2017ensembles}; seed
$3$ is a reserve. The bone image is the ensemble mean, the soft image is $x$ minus that mean, and the per-pixel spread is the ensemble
standard deviation:
\begin{equation}
\begin{gathered}
\bar f_p=\frac1K\sum_{k=1}^{K}f_k(x)_p,\qquad \hat s=x-\bar f,\\
\sigma_p=\Big(\frac{1}{K-1}\sum_{k=1}^{K}\big(f_k(x)_p-\bar f_p\big)^2\Big)^{1/2}.
\end{gathered}
\label{eq:ensemble}
\end{equation}
With the score $r_p=|b_p-\bar f_p|/(\sigma_p+\varepsilon)$, $\varepsilon=10^{-3}$, a split-conformal scale $q$ is the
$\lceil (n+1)(1-\alpha)\rceil$-th smallest of the $n$ scores pooled over calibration pixels of the Qure.ai validation CTs, and the
interval
\begin{equation}
\mathcal{C}_p=\big[\bar f_p-q\,(\sigma_p+\varepsilon),\;\bar f_p+q\,(\sigma_p+\varepsilon)\big]
\label{eq:conformal}
\end{equation}
has nominal coverage $1-\alpha=0.9$ \cite{lei2018conformal,angelopoulos2023conformal}. Calibration treats pixels as exchangeable, so the
guarantee is marginal over pixels, not over images.

\emph{Recalibration on real pairs (H9, post hoc).} After the intervals were found not to transfer to real radiographs
(Section~\ref{sec:res-uq}), an exploratory recalibration on the 24 SZCH development pairs was added and tested on the 51 held-out pairs
(pseudo-truth bone $b=x-s_{\mathrm{DES}}$ after a gain and offset fit). Each pixel is assigned a stratum
\begin{equation}
\begin{gathered}
k(p)=3\cdot\mathbf{1}\big[\bar f_p>0.25\,p_{99}(\bar f)\big]+t(\sigma_p)\in\{0,\dots,5\},\\
q_k=r^{(k)}_{(\lceil (n_k+1)(1-\alpha)\rceil)},
\end{gathered}
\label{eq:mondrian}
\end{equation}
where $t(\sigma_p)\in\{0,1,2\}$ is the spread tercile (edges from the calibration pixels) and $r^{(k)}_{(j)}$ the $j$-th smallest of the
$n_k$ calibration scores in stratum $k$; the interval of Eq.~\eqref{eq:conformal} then uses $q_{k(p)}$ (Mondrian calibration). The global
variant uses one stratum, and a per-image variant replaces the pooled pixel scores by the per-image $(1-\alpha)$ quantiles, calibrated at
level 0.1. The sparsification curve is the mean absolute error of the pixels that remain after removing the fraction $f\in[0,0.95]$ with
the largest spread; AUSE is the area between this curve and the oracle curve that removes pixels by their error (0 is perfect).

\subsubsection{Audit sites and lesions}
\label{sec:m-audit}
Lesions come from the CTs of the LIDC-IDRI test split only, so no audited lesion was seen in training. On each real radiograph six sites are
drawn at random positions inside the lung fields at $\alpha=1$; the sites depend only on the image, its lung mask, the lesion library and a
fixed seed, and are shared by every model. Lung masks come from the PSPNet anatomical segmentation model of TorchXRayVision
\cite{cohen2022torchxrayvision}, eroded by 49~px at $1024$~px so that a site lies inside the lung field; the model comes from the ChestX-Det
project \cite{lian2021chestxdet}, whose paper does not fully document its training data. The library holds 296 nodules (4.0 to 37.6~mm,
median 7.1~mm); six sites per image give 1,470 sites on VinDr-RibCXR, 1,482 on JSRT and 450 on SZCH. Because the PSPNet masks differ
slightly between GPUs, the masks were exported once on one machine and every evaluation reads them, so every model is scored on identical
sites.

The bone overlap of a site is measured with independent information, never with a model output: the radiologist rib masks of VinDr-RibCXR
\cite{nguyen2021ribcxr} for the confirmatory audit, and the real dual-energy bone signal on SZCH. A site \emph{overlaps bone} when at least
15\,\% of its support lies on an annotated rib. On JSRT no bone reference exists; the protocol named the consensus of the bone maps of ours
and of the reference model, but the audit took the overlap from the reference model's bone map alone (a secondary analysis;
Supplementary Section~S1).

\emph{Why the audit is not circular.} The twin regulariser pastes lesions to teach invariance, and the audit pastes lesions to measure it,
so a model could pass the audit by learning the insertion artefact. Four design choices limit this. The audit lesions come from CTs that
were never used for training; the audit images are real radiographs, whereas the training composites are synthetic; the confirmatory audit
uses rib masks that no model has seen; and the claim that the regulariser helps is also tested with evidence that does not use insertion:
contrast of real nodules on JSRT and Node21, detection of real lesions, and detection recall grouped by bone overlap (H4$'$). These choices
limit the circularity but cannot remove it: a model that recognises pasted signal in general could still score well, which is why the
evidence without insertion is reported beside the audit (Section~\ref{sec:limitations}).

\emph{Audit variants.} On the Qure.ai test split the lesion is added to the soft component, the true bone image is known, and the bone
error inside the lesion support can be measured against it in addition to $A$ (exact-truth audit, Supplementary Table~S5). The SZCH audit
ran on all 75 SZCH images (450 sites) during development and is reported as development evidence only; the SZCH contrasts in the result
tables use its 51 held-out images (306 sites). In an amplitude sweep, $A$ is recomputed for $\alpha\in\{0.25,0.5,1,2,4\}$ on the 450 SZCH
sites (development data, Supplementary Table~S6). For $\alpha\le1$ the absorption is flat (v1 model 0.301--0.306, no-twin model
0.215--0.218, reference model 0.134--0.136), so the leak is linear in the lesion amplitude; at $\alpha=2$ and 4 it falls for every model,
where the added signal starts to saturate at the top of the intensity range.

\subsubsection{Fidelity and suppression measures on real images}
\label{sec:m-fidelity}
\emph{Paired dual-energy data.} Tone differs between a radiograph and its DES image, so metrics on the SZCH pairs are reported as is
(``raw'') and after a per-image least-squares gain and offset of the estimate $\hat s$ onto the reference $s$ (``calibrated'', primary):
\begin{equation}
\begin{gathered}
(\hat a,\hat c)=\arg\min_{a,c}\sum_p\big(a\,\hat s_p+c-s_p\big)^2,\\
\mathrm{PSNR}_{\mathrm{cal}}=10\log_{10}\frac{1}{\frac{1}{|\Omega|}\sum_{p\in\Omega}\big(\mathrm{clip}_{[0,1]}(\hat a\,\hat s_p+\hat c)-s_p\big)^2},
\end{gathered}
\label{eq:psnrcal}
\end{equation}
with intensities in $[0,1]$ and $\Omega$ the evaluated region; SSIM$_{\mathrm{cal}}$ uses the same calibrated estimate. Calibrated PSNR and
SSIM are also reported inside the region where DES removed the most signal (top 10\,\% of calibrated radiograph minus DES, the bone region)
and outside it. Paired fidelity also uses PSNR and SSIM \cite{wang2004ssim} on the whole image, band-limited SSIM, MS-SSIM
\cite{wang2003msssim}, LPIPS \cite{zhang2018lpips} and a high-frequency energy ratio; the same smoothing is applied to every method and never
tuned.

\emph{Rib benchmark.} On VinDr-RibCXR (245 images; evaluated locally, never used for training) we compare the gradient magnitude
$G(I)=\lvert\nabla(g_{1}*I)\rvert$ (Sobel operator after a Gaussian blur of 1~px) of the suppressed image $\hat s$ with that of the
radiograph $x$:
\begin{equation}
\begin{gathered}
R_{\mathrm{rib}}=1-\frac{\sum_{p\in\mathcal{R}}G(\hat s)_p}{\sum_{p\in\mathcal{R}}G(x)_p},\\
R_{\mathrm{non}}=1-\frac{\sum_{p\in\mathcal{N}}G(\hat s)_p}{\sum_{p\in\mathcal{N}}G(x)_p},
\end{gathered}
\label{eq:rib}
\end{equation}
where $\mathcal{R}$ is the rib-boundary band inside the lungs and $\mathcal{N}$ the lung pixels away from ribs. We report both, their
difference $R_{\mathrm{rib}}-R_{\mathrm{non}}$ (selectivity) and the ROC-AUC of the gradient-reduction map for separating the two sets of
pixels.

\emph{Real-nodule contrast.} For a nodule box $\mathcal{B}$ on JSRT or Node21 and the ring $\mathcal{U}$ around it (the box grown by half
its size on each side),
\begin{equation}
\begin{gathered}
\mathrm{CNR}(I)=\frac{\mathrm{mean}_{\mathcal{B}}(I)-\mathrm{mean}_{\mathcal{U}}(I)}{\mathrm{sd}_{\mathcal{U}}(I)},\\
\Delta\mathrm{CNR}=\mathrm{CNR}(\hat s)-\mathrm{CNR}(x),
\end{gathered}
\label{eq:cnr}
\end{equation}
compared between suppressors as paired differences per nodule.

\subsubsection{Site-appearance calibration (exploratory)}
\label{sec:m-sitecal}
BS-LDM, trained on dual-energy pairs of the SZCH hospital, scores far higher than any CT-supervised model on the held-out pairs
(Section~\ref{sec:res-fidelity}). After the H2$'$ results were known, diagnostics on the 24 development pairs only asked why. Subtracting
more of the bone estimate lowered PSNR$_{\mathrm{cal}}$ rather than raising it; a bone signal known exactly but predicted at 512\,px would
reach 41.7\,dB, so resolution does not limit it; and most of the error was broader than 64\,px and present in the unprocessed radiograph
as well. The dual-energy images of this site therefore differ from a radiograph minus its bone image by a smooth appearance field that the
per-image gain and offset of the metric cannot absorb (Fig.~\ref{fig:sitecal_diag}). A protocol written before any held-out pair was scored
with it tests a smooth correction of this field:
\begin{equation}
\begin{gathered}
\hat s'=\hat s+U(F w),\\
w=\arg\min_{w}\sum_{i\in\mathcal{D}}\big\lVert \Pi\,G_{16}(s_i-\hat s_i)-F_i w\big\rVert_2^2+\lambda N\lVert w\rVert_2^2,
\end{gathered}
\label{eq:sitecal}
\end{equation}
where $\mathcal{D}$ are the 24 development pairs with dual-energy soft-tissue images $s_i$ and estimates $\hat s_i$, $G_{16}$ is a Gaussian
blur of 16\,px at 1024\,px, $\Pi$ samples a $64\times64$ grid, $F_i$ holds on that grid the features $[1,u,v,u^2,v^2,uv,I,I^2,B,Iv,Iu]$
(position $u,v$; blurred radiograph $I$ and bone estimate $B=x-\hat s$), $\lambda=10^{-3}$, $N$ is the number of grid samples and $U$ is
bicubic upsampling. One $w$ is fitted per method in the same way (identity, ours, no-twin control, v1 model, reference model) and applied
unchanged to the 51 held-out pairs; BS-LDM, trained on this site's images, is not corrected. The fitted weight on $B$ is about $+0.35$ for
every CT-supervised model, so the correction adds back about a third of the smoothed bone estimate. The correction adds only a smooth field
and so cannot restore an erased lesion. Its hypotheses H16--H18 are stated in Section~\ref{sec:m-stats}. The analysis is exploratory: it
re-uses held-out pairs already used for H2$'$, H2$'$ stays as reported, and calibrated numbers are always shown beside the uncorrected ones.

\begin{figure*}[!t]
\centering
\safefig[0.9\textwidth]{figures/fig_sitecal_diag.pdf}
\caption{Development-pair diagnostics behind the site calibration (24 SZCH development pairs; the held-out pairs were not used). (a)
Calibrated PSNR of the whole image and of the bone region when the bone estimate is subtracted with gain $k$ ($k=1$ as trained):
subtracting more lowers PSNR for both CT-supervised models. (b) Whole-image PSNR of ours as is, after the leave-one-out site correction of
Eq.~\eqref{eq:sitecal}, if every error component broader than 64 or 16\,px were removed, and for an oracle that knows the dual-energy bone
signal exactly but only at 512\,px.}
\label{fig:sitecal_diag}
\end{figure*}

% ---------------------------------------------------------------------------------------------------------------------
\subsection{Downstream Detection and Classification}
\label{sec:m-downstream}
The downstream classifier is a DenseNet-121 \cite{huang2017densenet} and the detector a Faster R-CNN with a ResNet-50 FPN backbone
\cite{ren2017fasterrcnn}, with identical settings for every input variant. The classifier sees $256$~px images, is trained for 10 epochs
with AdamW \cite{loshchilov2019adamw} (learning rate $3\times10^{-4}$, weight decay $10^{-4}$, batch 64, cosine schedule) and keeps the
epoch with the best validation macro-AUROC. The detector sees $512$~px images, is trained for 12 epochs with SGD (learning rate 0.01,
momentum 0.9, weight decay $10^{-4}$, batch 8) and keeps the epoch with the best validation CPM. Each input is trained with three seeds:
the original radiograph, the soft-tissue image of ours or of the reference model, CLAHE, and a multi-channel input that stacks images as
channels (for the classifier the original and our soft image, for the detector the original, our soft and our bone image).
Classification uses NIH ChestX-ray14 (macro-AUROC). Detection uses Node21, TBX11K and VinBigData and is scored by FROC sensitivity at fixed false-positive rates, averaged as
the competition performance metric (CPM) \cite{vanginneken2010anode09,niemeijer2011cpm}, and by mAP@50; detection metrics treat tied scores
as one threshold, so they do not depend on the order in which a sort returns ties.

% ---------------------------------------------------------------------------------------------------------------------
\subsection{Study~II: Multi-Stream Fusion with Capacity-Matched Controls}
\label{sec:m-study2}

\subsubsection{Rationale}
With the frozen suppressor ensemble $\bar f$ of Study~I, a classifier can receive the radiograph $X$, the soft-tissue image
$S=X-\bar f(X)$ and the bone image $\hat B=\mathrm{clip}_{[0,1]}(X-S)$ on the stored 8-bit scale. Because $(S,\hat B)$ is a deterministic
function of $X$, the inputs carry no new information about a label $Y$:
\begin{equation}
I\big(Y;\,X,S,\hat B\big)=I\big(Y;\,X\big).
\label{eq:info}
\end{equation}
A gain can therefore come only from the prior the suppressor learned from CT, or from the inductive bias of the network that reads the
extra images; a network that reads three images also spends three times the encoder computation. The study separates these sources with a
control of the same capacity, and it analyses the effect per finding group, because suppression could help focal lesions behind ribs while
removing cues for diffuse or pleural findings.

\subsubsection{Input configurations and network}
Five configurations were trained with identical settings (Figs.~\ref{fig:architecture}b and~\ref{fig:cls_design}):
\emph{orig} ($X$), \emph{soft} ($S$), \emph{early} ($X$, $S$ and $\hat B$ as the three input channels of one stream), \emph{hybrid} (three
streams $X$, $S$, $\hat B$) and \emph{ctrl} (the hybrid network with $X$ in all three streams). The soft-tissue image comes from our
suppressor ensemble run at its input size of 512\,px; classifiers see all images at $256\times256$\,px. One DenseNet-121 encoder $E$
\cite{huang2017densenet} with shared weights, pretrained on ImageNet \cite{deng2009imagenet,wightman2019timm}, processes each stream $u_s$.
Each stream is standardised per image and passed through its own learned affine map, because the bone image spans a narrow intensity range;
the $8\times8$ encoder output is projected to $D=256$-dimensional tokens with a stream embedding $e_s$ and a two-dimensional sine position
code $\pi_n$:
\begin{equation}
\begin{gathered}
\tilde u_s=\gamma_s\,\frac{u_s-\mathrm{mean}(u_s)}{\mathrm{sd}(u_s)}+\beta_s,\\
z_{s,n}=W\,E(\tilde u_s)_n+e_s+\pi_n\in\mathbb{R}^{D}.
\end{gathered}
\label{eq:tokens}
\end{equation}
A two-layer decoder with eight heads and one learned query $q_c$ per class attends to the tokens of all streams, as in query-based
multi-label heads \cite{liu2021query2label,ridnik2023mldecoder}, and a linear map per class gives the logit $\ell_c$ from the decoder
output $o_c$. With keys $k_{s,n}$, the head-averaged cross-attention of the last layer defines the routing of class $c$ to stream $s$:
\begin{equation}
\begin{gathered}
a_{c,(s,n)}=\operatorname*{softmax}_{(s,n)}\Big(\frac{q_c^{\top}k_{s,n}}{\sqrt{D}}\Big),\\
\rho_{c,s}=\sum_{n}a_{c,(s,n)},\qquad
\ell_c=w_c^{\top}o_c+b_c.
\end{gathered}
\label{eq:routing}
\end{equation}
Single-stream configurations use the same encoder and decoder over 64 tokens. The shared encoder gives the hybrid the parameters of a
single-stream model plus the decoder, at three times the encoder computation; \emph{ctrl} has exactly the hybrid's parameters and
computation but no suppression.

\begin{figure*}[!t]
\centering
\safefig[0.9\textwidth]{figures/fig_study_design.pdf}
\caption{Fusion study design. (a) The three inputs for an NIH ChestX-ray14 test image: the radiograph $X$, the soft-tissue image $S$
predicted by the suppressor and the bone image $\hat B=X-S$ (box: expert annotation of a focal finding). (b) The five configurations, all
with the same DenseNet-121 encoder, class-query decoder and training settings. In the hybrid and the control one encoder with shared
weights processes each stream and the decoder attends to the tokens of all three; the control receives $X$ three times. Bottom: the
contrasts tested.}
\label{fig:cls_design}
\end{figure*}

\subsubsection{Training}
The loss is the binary cross-entropy over the annotated labels $\mathcal{V}$ of a batch (one-hot targets for the three BDCXR-3257 classes):
\begin{equation}
\mathcal{L}_{\mathrm{cls}}=-\frac{1}{|\mathcal{V}|}\sum_{(i,c)\in\mathcal{V}}\Big[y_{ic}\log\sigma(\ell_{ic})+(1-y_{ic})\log\big(1-\sigma(\ell_{ic})\big)\Big].
\label{eq:bce}
\end{equation}
Training used AdamW \cite{loshchilov2019adamw} (learning rate $3\times10^{-4}$, weight decay $10^{-4}$), batch size 32, 20 epochs with a
cosine schedule after 200 warm-up steps, and mixed precision. Augmentation: random rotation by up to $7^\circ$, scaling by 0.9--1.1,
translation by up to 3\,\% of the image size, and brightness and contrast changes, identical across the channels of an image. The epoch with
the best validation macro-AUROC was kept. Every configuration was trained with seeds 0--4 on NIH ChestX-ray14 and VinBigData and with seeds
0--4 in each of the five BDCXR-3257 folds of Section~\ref{sec:m-bdcxr}, giving 175 runs. A configuration's prediction is the seed ensemble
\begin{equation}
\bar p_c(x)=\frac15\sum_{k=0}^{4}\sigma\big(\ell^{(k)}_c(x)\big),
\label{eq:seedens}
\end{equation}
and all tests compare seed ensembles. Runs used cloud T4 GPUs; a run longer than one session resumed from an end-of-epoch checkpoint, which
gave the same predictions as uninterrupted training in a test.

\subsubsection{Encoder selection}
We screened four encoders on the NIH validation split (configuration \emph{orig}, seed 0, 10 epochs): DenseNet-121, ConvNeXt-V2-T
\cite{woo2023convnextv2}, RAD-DINO \cite{perezgarcia2025raddino} and SwinV2-B \cite{liu2022swinv2}. The best encoder would go forward with
DenseNet-121, and an encoder within 0.005 of the best would yield to a cheaper one. All four reached validation macro-AUROC between 0.8384
and 0.8430, and DenseNet-121 was the cheapest, so it went forward alone. ConvNeXt-V2-T failed to train at the planned learning rate;
following the rule of the suppressor study, the three encoders without batch normalisation were screened again with the encoder learning
rate reduced fourfold, a change decided before the other screen results were seen (Supplementary Section~S1).

\subsubsection{Finding groups}
Before any run, the classes were grouped as \emph{focal} (NIH: nodule, mass; VinBigData: nodule/mass), \emph{diffuse} (effusion,
infiltration, consolidation, oedema, pneumonia; VinBigData: pleural effusion, infiltration, consolidation, lung opacity, interstitial lung
disease) and \emph{boundary} (cardiomegaly, atelectasis, pneumothorax, pleural thickening, emphysema, fibrosis, hernia; VinBigData:
cardiomegaly, atelectasis, pneumothorax, pleural thickening, aortic enlargement, pulmonary fibrosis, calcification). VinBigData's ``other
lesion'' belongs to no group, and BDCXR-3257 (normal, bacterial or viral pneumonia) enters only the overall tests. For a group $G$ of classes
and two configurations $m$ and $m'$,
\begin{equation}
\begin{gathered}
\mathrm{AUROC}_{\mathrm{macro}}(m)=\frac1C\sum_{c=1}^{C}\mathrm{AUROC}_c(m),\\
\Delta_G(m,m')=\frac{1}{|G|}\sum_{c\in G}\big(\mathrm{AUROC}_c(m)-\mathrm{AUROC}_c(m')\big).
\end{gathered}
\label{eq:group}
\end{equation}

\subsubsection{Attribution maps and the pointing game}
\label{sec:m-xai}
For each hybrid model, Grad-CAM++ \cite{chattopadhay2018gradcampp} was computed per stream on the last encoder block. Integrated Gradients
\cite{sundararajan2017axiomatic} was computed on the three inputs jointly against a blurred baseline ($\sigma=8$\,px, 32 steps), because a
black image has no contrast after standardisation; the decoder attention was read from its last layer. A target $t$ is a pair of test image
and box class; a map $M_m$ hits when its maximum lies inside a box of that class, with a tolerance of 15\,px at 256\,px fixed in advance
\cite{zhang2018topdown}. With $h^{(k)}_m(t)\in\{0,1\}$ the hit of map $m$ of the seed-$k$ hybrid model, focal targets $\mathcal{F}$
(masses and nodules) and the other targets $\mathcal{O}$ (six box classes),
\begin{equation}
\begin{gathered}
d(t)=\frac15\sum_{k=0}^{4}\big(h^{(k)}_{S}(t)-h^{(k)}_{X}(t)\big),\\
I=\frac{1}{|\mathcal{F}|}\sum_{t\in\mathcal{F}}d(t)-\frac{1}{|\mathcal{O}|}\sum_{t\in\mathcal{O}}d(t),
\end{gathered}
\label{eq:pointing}
\end{equation}
and H10 tests $I>0$ with an image-cluster bootstrap (\num{10000} resamples; Eq.~\eqref{eq:bootp}).

\subsubsection{Classification uncertainty}
\label{sec:m-clsuq}
\emph{NIH: false-negative risk control.} An image's label set holds every class with ensemble probability at least $1-\lambda$; the threshold
is calibrated on the validation split (\num{8634} images) by conformal risk control \cite{angelopoulos2022crc} and the false-negative rate is
measured on the test split:
\begin{equation}
\begin{gathered}
\mathcal{C}_\lambda(x)=\{c:\bar p_c(x)\ge1-\lambda\},\\
\hat R_n(\lambda)=\frac1n\sum_{i=1}^{n}\Big(1-\frac{|\mathcal{C}_\lambda(x_i)\cap Y_i|}{|Y_i|}\Big),\\
\hat\lambda=\min\Big\{\lambda:\frac{n\,\hat R_n(\lambda)+1}{n+1}\le\alpha\Big\},
\end{gathered}
\label{eq:crc}
\end{equation}
over the $n$ calibration images with at least one positive label $Y_i$, at $\alpha=0.10$.

\emph{BDCXR-3257: class-conditional sets.} With probabilities normalised to sum to one, $\hat p_y=\bar p_y/\sum_{y'}\bar p_{y'}$, the score
$1-\hat p_y$ is calibrated separately for each true class \cite{vovk2005algorithmic}, on the other folds' out-of-fold predictions for each
test fold:
\begin{equation}
\begin{gathered}
\mathcal{C}(x)=\{y:1-\hat p_y(x)\le q_y\},\\
q_y=\text{the }j_y\text{-th smallest of }\{1-\hat p_{y_i}(x_i):y_i=y\},\\
j_y=\lceil (n_y+1)(1-\alpha)\rceil.
\end{gathered}
\label{eq:mondriansets}
\end{equation}

\emph{Selective classification.} Images are ranked by a confidence score and the most confident fraction is kept
\cite{geifman2017selective,geifman2019bias}. Score A is one minus the entropy $H$ of the seed ensemble (mean binary entropy over the 14
labels for NIH, Shannon entropy of the normalised probabilities for BDCXR-3257); score B adds the suppressor's own uncertainty,
$-(z_H+z_\sigma)$, where $\sigma$ is the spread of Eq.~\eqref{eq:ensemble} averaged over the lung field and $z$ standardises over the
evaluated images of a data set. With the losses $\ell_{(1)},\dots,\ell_{(N)}$ sorted by decreasing confidence (Brier score for NIH, top-1
error for BDCXR-3257),
\begin{equation}
\mathrm{AURC}=\frac1N\sum_{k=1}^{N}\frac1k\sum_{j=1}^{k}\ell_{(j)}.
\label{eq:aurc}
\end{equation}

% ---------------------------------------------------------------------------------------------------------------------
\subsection{External Validation on Paediatric Radiographs}
\label{sec:m-bdcxr}
The protocol of this validation was fixed before any model of this study saw a BDCXR-3257 image; its one later change, the duplicate
grouping, was made from the images alone before any model ran (Supplementary Section~S1). All suppressors are the frozen models; only the
downstream classifiers are trained. (E1, primary) Pneumonia classification with three one-vs-rest outputs, inputs as in H4$'$b (original,
ours, reference model, CLAHE, original plus ours), the classifier settings of the downstream protocol, five folds stratified by class over
the near-duplicate groups (fold seed 0) with 10\,\% of each training part, by whole groups, for epoch selection, and three seeds per fold and
input (75 runs). Out-of-fold probabilities, averaged over seeds, give one prediction per image. The primary measure is the pooled
macro-AUROC; ours against the original image is non-inferior if the lower bound of the 95\,\% interval of the difference (group bootstrap)
exceeds $-0.01$, and superior if it exceeds 0; per class, paired DeLong tests with Holm over the three classes. (E2, secondary) The lesion
audit on a fixed class-stratified subset of 300 images (seed 7), six sites per image, with bone overlap from the reference model's bone map;
images whose lung mask covers less than 2\,\% of the image are skipped and counted (none was). (E3, descriptive) The per-image ensemble
spread over the lung field, compared with the held-out SZCH pairs and correlated with age.

% ---------------------------------------------------------------------------------------------------------------------
\subsection{Statistical Analysis}
\label{sec:m-stats}

\subsubsection{General framework}
Hypotheses, margins, test sets and tests were written down before the data they govern were evaluated, and they carry the labels under
which they were specified (H1$'$--H5$'$ for the final suppressors, H1 for the v1 models, H6--H10 for the fusion study, H16--H18 for the site
calibration and E1--E3 for the external validation). Analyses added later are labelled exploratory or post hoc, every change made after the
analyses started is listed in Supplementary Section~S1, and all primary and secondary results are reported whatever their direction.

All comparisons are paired. Confidence intervals are percentile bootstrap intervals \cite{efron1993bootstrap} (\num{10000} resamples unless
stated) over patients when several units come from one patient; paired differences are tested with the Wilcoxon signed-rank test
\cite{wilcoxon1945individual}; AUCs are compared with the paired DeLong test \cite{delong1988auc}. Seeds are averaged within unit before
testing and their spread is reported. Multiplicity is handled by the Holm procedure \cite{holm1979simple} within each family: with the $m$
$p$-values of a family ordered as $p_{(1)}\le\dots\le p_{(m)}$,
\begin{equation}
\tilde p_{(j)}=\max_{i\le j}\,\min\big\{1,\,(m-i+1)\,p_{(i)}\big\}.
\label{eq:holm}
\end{equation}
Confidence intervals are unadjusted 95\,\% intervals; effect sizes are always reported with them, and a hypothesis that fails is reported as
failed. Non-inferiority uses the paired mean difference $\bar\Delta=\frac1n\sum_i(m_i^{\mathrm{ours}}-m_i^{\mathrm{ref}})$ and the
one-sided 95\,\% lower bound $L_{95}$ of its image-level bootstrap distribution; it holds if
\begin{equation}
L_{95}(\bar\Delta)>-\delta
\label{eq:ni}
\end{equation}
for a margin $\delta$ fixed in advance.

\subsubsection{Study~I hypotheses (H1$'$--H5$'$)}
\begin{description}
\item[H1$'$, dose-response.] On VinDr-RibCXR (245 images, six sites each, lesions from LIDC test CTs, $\alpha=1$), the per-site absorption
 $A$ averaged over seeds falls with the twin weight. With $\bar A_{s,w}$ the absorption at site $s$ averaged over the seeds trained with
 twin weight $w$, the slope $\beta_1$ of
 \begin{equation}
 \bar A_{s,w}=\beta_0+\beta_1\log_2(1+w)+e_{s,w}
 \label{eq:dose}
 \end{equation}
 has an image-cluster bootstrap interval, and H1$'$ passes if its upper bound is below zero. The same analysis on the Qure.ai test split
 (exact truth) is secondary.
\item[H2$'$, fidelity.] On the 51 held-out SZCH pairs, bone-region PSNR of ours against the reference model is non-inferior
 (Eq.~\eqref{eq:ni}) with margin $\delta=0.5$~dB; SSIM has a margin of 0.01 and is reported either way; ours against the no-twin control
 uses a margin of 0.25~dB (the cost of the regulariser). Metrics are computed calibrated (primary) and raw (Eq.~\eqref{eq:psnrcal}).
\item[H3$'$, lesion absorption.] On VinDr-RibCXR sites, $A$(ours)$<A$(reference model) and $A$(ours)$<A$(no-twin control), as paired
 site-level differences with image-cluster bootstrap intervals, Holm over the two contrasts, and an interaction between sites that overlap
 an annotated rib and sites that do not.
\item[H4$'$, evidence without insertion.] (a) $\Delta$CNR of real nodules on JSRT and Node21 (Eq.~\eqref{eq:cnr}), ours against the
 reference model and against the no-twin control, for all lesions and for lesions with at least 15\,\% bone overlap. (b) Downstream: NIH
 ChestX-ray14 macro-AUROC per input (original, ours, reference model, CLAHE, original plus ours; three seeds); TBX11K, Node21 and
 VinBigData detection by FROC and mAP@50; recall at two false positives per image for lesions with and without bone overlap, as a
 difference in differences against the original input.
\item[H5$'$, reliability.] Spearman correlation between ensemble spread and absolute bone error on the SZCH held-out pairs (pseudo-truth
 bone from the least-squares-aligned DES image); split-conformal coverage calibrated on Qure.ai validation CTs and measured on the Qure.ai
 test split (nominal 90\,\%) and, as measured under shift, on SZCH.
\end{description}
Site-level differences between models are paired, with image-cluster bootstrap intervals. Bootstrap intervals use \num{10000} resamples for
mean absorption and contrasts, \num{2000} for the overlap interaction and 500 for the dose-response slope, always resampling whole images
with replacement. The hypothesis H1 of the v1 models (insertion lowers absorption on synthetic LIDC test CTs at non-inferior bone error) is
tested in the same way on exact truth.

\emph{Site calibration (exploratory, Section~\ref{sec:m-sitecal}).} H16 (primary) is a gain in the bone-region PSNR$_{\mathrm{cal}}$ of ours
from the calibration (one-sided 95\,\% lower bound of the paired difference above 0; image bootstrap, \num{10000} resamples) and H16b the
same for whole-image PSNR$_{\mathrm{cal}}$. H17 is the non-inferiority of the calibrated ours to the calibrated reference model in
bone-region PSNR$_{\mathrm{cal}}$ (margin 0.5\,dB, as H2$'$). H18 is a lower absorption of the calibrated ours than of the calibrated
reference model on VinDr-RibCXR, from one audit run with all four models so that they share every site (image-cluster bootstrap, two-sided
95\,\% upper bound below 0). The rib benchmark of the calibrated models is descriptive. H16 is the only confirmatory test of this analysis;
the others are reported without multiplicity correction.

\subsubsection{Study~II hypotheses (H6--H10)}
H9 concerns the suppressor's intervals (Section~\ref{sec:m-ensemble}); the others concern the classifiers.
\begin{description}
\item[H6 (primary).] The hybrid has a higher macro-AUROC than \emph{orig}; one test per data set, Holm correction over the three.
\item[H7 (secondary).] Hybrid against \emph{ctrl} (does a gain come from the suppression or from the extra computation?) and against
 \emph{early} (do separate streams beat channel stacking?); one Holm family over the six tests.
\item[H8 (secondary).] $\Delta_G$ of Eq.~\eqref{eq:group} for hybrid and for \emph{soft} against \emph{orig} in each finding group; Holm
 within data set over the six tests. Per class, hybrid against \emph{orig} and against \emph{early} with the paired DeLong test
 \cite{delong1988auc}, Holm within data set over classes and both contrasts.
\item[H10 (secondary).] In the hybrid models, the pointing-game hit rate of the soft-tissue stream minus that of the radiograph stream is
 larger for focal targets than for the other targets (Eq.~\eqref{eq:pointing}).
\end{description}
Every contrast $\Delta$ of two seed ensembles is resampled with a paired cluster bootstrap: replicate $b$ draws clusters with replacement
(patients for NIH, images for VinBigData, near-duplicate groups for BDCXR-3257), recomputes $\Delta^{*}_b$ with the same resamples for every
contrast of a data set, and with $B=\num{10000}$ replicates gives the percentile interval $[\Delta^{*}_{(0.025)},\Delta^{*}_{(0.975)}]$ and
\begin{equation}
\begin{gathered}
p=\min\big\{1,\;2\min(p_{-},p_{+})\big\},\\
p_{-}=\frac{1+\#\{b:\Delta^{*}_b\le0\}}{B+1},\quad p_{+}=\frac{1+\#\{b:\Delta^{*}_b\ge0\}}{B+1}.
\end{gathered}
\label{eq:bootp}
\end{equation}
AUPRC differences are reported with bootstrap intervals only. The hypotheses, the code and the split files of this study were fixed before
the encoder screen and before any classifier of this study saw a test split.

\section{Experimental Results}
\label{sec:results}
Of the Study~I hypotheses, H1$'$ (dose-response) passes; H2$'$ passes for bone-region PSNR against the reference model and against the
no-twin control and fails for SSIM; H3$'$ passes for both contrasts on VinDr-RibCXR; H4$'$a is mixed (ours above the reference model on
JSRT and below it on Node21, above the no-twin control on both); H4$'$b shows no downstream benefit of the soft-tissue image used alone but
gains when the suppressed images are added as extra channels; and H5$'$ holds on synthetic data with exact truth and fails under the shift
to real radiographs. The v1 hypothesis H1 passes, and the exploratory site calibration passes H16--H18. In Study~II the primary hypothesis
H6 fails on all three data sets and H7 finds no difference from the capacity control, whereas H8 finds a gain for focal and a loss for
diffuse findings on NIH and H10 passes; the post hoc recalibration of the bone intervals (H9) brings their coverage close to nominal.
Table~\ref{tab:hypotheses} lists every hypothesis with its main estimate and outcome; the subsections below give the details.

\begin{table*}[!t]
\centering
\caption{Hypotheses, their main estimate with 95\,\% interval (or one-sided bound) and outcome. Study~I: final suppressors (H1$'$--H5$'$),
v1 models (H1), site calibration (H16--H18, exploratory) and external validation (E1--E3); Study~II: fusion study (H6--H10; H9 concerns
the suppressor's intervals and is post hoc). $\Delta$: ours (or the first-named model) minus the comparator.}
\label{tab:hypotheses}
\footnotesize
\renewcommand{\arraystretch}{1.08}
\setlength{\tabcolsep}{3pt}
\begin{tabularx}{\textwidth}{@{}>{\raggedright\arraybackslash}p{4.3cm}>{\raggedright\arraybackslash}X%
>{\raggedright\arraybackslash}p{4.6cm}>{\raggedright\arraybackslash}p{2.3cm}@{}}
\toprule
Hypothesis & Data and test & Estimate & Outcome \\
\midrule
H1$'$ absorption falls with the twin weight & VinDr-RibCXR; slope of $A$ on $\log_2(1+w)$, image-cluster bootstrap & $-0.0233$
 $[-0.0249,-0.0217]$ & passed \\
H2$'$ bone-region PSNR$_{\mathrm{cal}}$, ours vs.\ reference & 51 held-out SZCH pairs; non-inferiority, margin 0.5\,dB & $-0.06$\,dB
 (lower bound $-0.15$\,dB) & passed \\
H2$'$ bone-region PSNR$_{\mathrm{cal}}$, ours vs.\ no-twin & same; margin 0.25\,dB & $+0.01$\,dB (lower bound $-0.02$\,dB) & passed \\
H2$'$ SSIM$_{\mathrm{cal}}$, ours vs.\ reference & same; margin 0.01 & $-0.021$ (lower bound $-0.023$) & failed \\
H3$'$ $A$, ours vs.\ reference / no-twin & VinDr-RibCXR, 1,470 sites; paired, Holm & $-0.098$ $[-0.108,-0.089]$ / $-0.127$
 $[-0.134,-0.119]$ & passed \\
H4$'$a $\Delta$CNR of real nodules, ours vs.\ reference & JSRT (154) / Node21 (1,474) & $+0.023$ $[0.008,0.041]$ / $-0.025$
 $[-0.033,-0.018]$ & mixed \\
H4$'$b downstream, our soft image alone & NIH macro-AUROC & $-0.001$ $[-0.004,0.003]$ & no gain \\
H4$'$b downstream, extra channels & NIH macro-AUROC / Node21 CPM & $+0.006$ $[0.003,0.008]$ / $+0.028$ $[0.009,0.047]$ & gain \\
H5$'$ conformal coverage (nominal 0.90) & Qure.ai test split / 51 held-out SZCH pairs & 0.900 $[0.897,0.903]$ / 0.139
 $[0.134,0.145]$ & held / failed under shift \\
H1 (v1) CLIA vs.\ recipe & synthetic LIDC-IDRI test CTs, exact truth & $-0.022$ $[-0.032,-0.013]$ & passed \\
H16 / H16b site-calibration gain of ours & 51 held-out SZCH pairs; one-sided bound & $+2.12$\,dB (bound 1.99) / $+2.74$\,dB (bound 2.61)
 & passed \\
H17 calibrated ours vs.\ calibrated reference & same; margin 0.5\,dB & $-0.35$\,dB (lower bound $-0.44$\,dB) & passed \\
H18 $A$, calibrated ours vs.\ calibrated reference & VinDr-RibCXR, 1,470 sites & $-0.096$ $[-0.105,-0.087]$ & passed \\
E1 pneumonia macro-AUROC, ours vs.\ original & BDCXR-3257, five folds; margin 0.01 & $+0.002$ $[-0.005,0.010]$ & non-inferior \\
E2 $A$, ours vs.\ reference & BDCXR-3257, 300 images & $-0.134$ $[-0.144,-0.124]$ & lower \\
E3 ensemble spread, children vs.\ adults & BDCXR-3257 vs.\ SZCH held-out & 2.7-fold (0.0076 vs.\ 0.0028) & descriptive \\
\midrule
H6 (primary) macro-AUROC, hybrid vs.\ original & NIH / VinBigData / BDCXR-3257; Holm & $+0.0001$ $[-0.0022,0.0024]$ / $+0.0024$
 $[-0.0004,0.0053]$ / $-0.0042$ $[-0.0108,0.0021]$ & failed \\
H7 hybrid vs.\ control / vs.\ early & NIH; Holm over six tests & $+0.0007$ $[-0.0018,0.0032]$ / $-0.0031$ $[-0.0055,-0.0007]$ & no
 difference \\
H8 focal findings, soft vs.\ original & NIH; Holm & $+0.0089$ $[0.0040,0.0140]$, Holm $p=0.0024$ & gain \\
H8 diffuse findings, hybrid vs.\ original & NIH; Holm & $-0.0052$ $[-0.0077,-0.0027]$, Holm $p=0.0024$ & loss \\
H9 recalibrated bone intervals (post hoc) & 24 development / 51 held-out SZCH pairs & coverage 0.886 (global), 0.895 (Mondrian) & near
 nominal \\
H10 pointing game, soft minus radiograph stream, focal vs.\ other & NIH, 984 expert boxes & $+0.145$ $[0.092,0.199]$, $p=0.0002$ &
 passed \\
\bottomrule
\end{tabularx}
\end{table*}

% ---------------------------------------------------------------------------------------------------------------------
\subsection{Fidelity, Dose-Response and Ablation (Study~I: H1\texorpdfstring{$'$}{'}, H2\texorpdfstring{$'$}{'})}
\label{sec:res-fidelity}

\subsubsection{Fidelity on held-out paired radiographs (H2$'$)}
On the 51 held-out SZCH-X-Rays pairs (Table~\ref{tab:fidelity}) the bone-region PSNR$_{\mathrm{cal}}$ of ours is 21.46\,dB and that of the
reference model 21.51\,dB. The paired difference is $-0.06$\,dB with a one-sided 95\,\% lower bound of $-0.15$\,dB, inside the
non-inferiority margin of $-0.5$\,dB. Against the no-twin control the difference is $+0.01$\,dB (lower bound $-0.02$\,dB, margin
$-0.25$\,dB): the twin regulariser costs no measurable bone-region fidelity. SSIM$_{\mathrm{cal}}$ fails its margin: ours is lower by 0.021
(lower bound $-0.023$, margin $-0.01$). The gap sits in fine texture. With a band limit of $\sigma=4$ the SSIM of ours and of the reference
model are 0.946 and 0.948, and the high-frequency energy of the non-bone region, relative to the dual-energy soft-tissue image, is 1.09 for
ours and 0.99 for the reference model; LPIPS is 0.234 and 0.166. BS-LDM, which was trained on dual-energy pairs from the same hospital,
reaches 28.57\,dB; it is shown for orientation and is not one of the planned comparisons. Fig.~\ref{fig:pareto} places fidelity against
lesion absorption, and Fig.~\ref{fig:qualitative} shows three held-out cases.

\begin{table*}[!t]
\centering
\caption{Held-out SZCH-X-Rays pairs (51): calibrated fidelity of ours and the comparators (top; HF ratio: high-frequency energy of the
non-bone region relative to the dual-energy soft-tissue image, 1 = equal) and the H2$'$ non-inferiority tests (bottom; paired differences,
image bootstrap, one-sided 95\,\% lower bound). Every method and seed, raw and calibrated, is in Supplementary Table~S1.}
\label{tab:fidelity}
\restable[\textwidth]{tmi_fidelity}
\subtab{H2$'$ non-inferiority}
\restable[\textwidth]{tmi_h2}
\end{table*}

\begin{figure}[!t]
\centering
\safefig[\columnwidth]{figures/pareto.pdf}
\caption{Lesion absorption on VinDr-RibCXR against bone-region PSNR$_{\mathrm{cal}}$ on the 51 held-out SZCH-X-Rays pairs. Line: the
SwinV2-T family at twin weights $w=0\ldots64$ (single seeds light, seed means dark). Markers: ours and the no-twin control
(ConvNeXt-V2-T, $K=3$), the reference model, the v1 model and the identity. Lower absorption at equal fidelity is better.}
\label{fig:pareto}
\end{figure}

\begin{figure*}[!t]
\centering
\safefig[0.8\textwidth]{figures/qualitative_szch.png}
\caption{Held-out SZCH-X-Rays cases (CC BY-NC 4.0 \cite{sun2026bsldm}): radiograph, dual-energy soft-tissue reference, the soft-tissue
estimates of ours and of the reference model, and the bone estimate of ours. The three cases are not hand-picked: they are the held-out
images at the 25th, 50th and 75th percentile of the bone-region PSNR$_{\mathrm{cal}}$ of ours.}
\label{fig:qualitative}
\end{figure*}

\subsubsection{Dose-response of the twin regulariser (H1$'$)}
In the SwinV2-T family (3, 2, 2, 4, 2 and 2 seeds at $w=0,4,8,16,32,64$; Table~\ref{tab:twin}, Fig.~\ref{fig:dose_response}) the mean
absorption on VinDr-RibCXR falls from 0.183 at $w=0$ to 0.043 at $w=64$. The slope of $A$ on $\log_2(1+w)$ is $-0.0233$
$[-0.0249,-0.0217]$; its upper bound is below zero, so H1$'$ passes. On the Qure.ai test split with exact truth the slope is $-0.0174$
$[-0.0190,-0.0160]$, with $A$ levelling off above $w=16$ (0.118, 0.122 and 0.112 at $w=16,32,64$). Bone-region PSNR$_{\mathrm{cal}}$ on the
held-out pairs stays between 21.22 and 21.28\,dB up to $w=16$ and drops to 20.92 and 21.01\,dB at $w=32$ and 64, which is why $w=16$ is the
working point (Fig.~\ref{fig:pareto}).

\begin{table*}[!t]
\centering
\caption{Twin regulariser. Top: dose-response of the lesion absorption on the twin weight (SwinV2-T family): seeds per weight, mean $A$ on
VinDr-RibCXR and on the Qure.ai test split, and the slope with its image-cluster bootstrap interval; each weight averages all of its seeds
(four at $w=16$, seeds 0--3), whereas the encoder comparisons use seeds 0--2. Bottom: ablation at $w=16$ (SwinV2-T, 300k samples, mean over
seeds), one factor changed per row; development evidence (audit on all 75 SZCH images, bone PSNR$_{\mathrm{cal}}$ on the 24 development
pairs).}
\label{tab:twin}
\restable[\textwidth]{dose_response}
\subtab{Ablation at $w=16$}
\restable[\textwidth]{tmi_ablation}
\end{table*}

\begin{figure}[!t]
\centering
\safefig[\columnwidth]{figures/dose_response.pdf}
\caption{Dose-response of the lesion absorption $A$ on VinDr-RibCXR on the twin weight $w$ (SwinV2-T), one marker per seed, with the slope
on $\log_2(1+w)$.}
\label{fig:dose_response}
\end{figure}

\subsubsection{Ablation of the regulariser}
The ablation (Table~\ref{tab:twin}, bottom) is development evidence: it uses the audit on all 75 SZCH images and the 24 development pairs.
At $w=16$ the full regulariser gives $A=0.067$ against 0.190 without it. Three variants lose the effect: a stop-gradient on the branch
without the pasted lesion (0.185), a consistency term over the whole image instead of the lesion support (0.197), and lesion insertion
without the twin term (CLIA, 0.191). Dilating the support by 31\,px keeps part of it (0.137); up to four lesions per twin instead of up to
two keep all of it (0.074). The effect therefore needs gradients through both branches and a consistency term confined to where the lesion
was pasted; showing the model lesions during training is not enough.

\subsubsection{The v1 hypothesis H1}
The v1 hypothesis H1 was tested on synthetic LIDC test CTs with exact truth, three seeds per arm (Supplementary Table~S8), and passes: CLIA
absorbs less than the recipe arm ($-0.022$ $[-0.032,-0.013]$, Holm $p<0.001$) with non-inferior bone MAE. The lesion-weighted loss also
lowers absorption ($-0.019$) and is non-inferior on the test CTs, although it missed the margin on the validation CTs that decided the
selection; combining both lowers absorption most ($-0.066$) but misses the margin. The 300k-sample finals of CLIA and recipe replicate the
contrast ($-0.024$ $[-0.034,-0.014]$). On real radiographs the v1 models remain far above the v2 models ($A=0.289$ and 0.308 on
VinDr-RibCXR; Table~\ref{tab:audit_means}).

% ---------------------------------------------------------------------------------------------------------------------
\subsection{Counterfactual Audit, Rib Benchmark and Real-Nodule Contrast (Study~I: H3\texorpdfstring{$'$}{'},
H4\texorpdfstring{$'$}{'}a)}
\label{sec:res-audit}

\subsubsection{Audit on VinDr-RibCXR (H3$'$)}
On 1,470 sites in 245 VinDr-RibCXR images (Tables~\ref{tab:audit_means} and~\ref{tab:audit}) ours absorbs $A=0.051$ $[0.045,0.057]$ of an
inserted lesion, the reference model 0.149 $[0.138,0.160]$ and the no-twin control 0.178 $[0.167,0.188]$. The paired differences are
$-0.098$ $[-0.108,-0.089]$ against the reference model and $-0.127$ $[-0.134,-0.119]$ against the no-twin control (Holm-adjusted $p<0.001$
for both), so H3$'$ passes. The advantage does not depend on rib overlap here: the interaction between sites with at least 15\,\% overlap
with an annotated rib and the other sites is $0.003$ $[-0.017,0.027]$ against the reference model. On JSRT (secondary) the pattern repeats
($A$ 0.048, 0.118 and 0.121; differences $-0.070$ and $-0.073$) with a larger advantage at bone-overlapping sites (interaction $-0.060$
$[-0.074,-0.046]$); there the overlap was defined by the reference model's bone map. The SZCH audit (development evidence) agrees ($A$ 0.054,
0.130 and 0.195). Fig.~\ref{fig:audit_maps} shows one site.

\begin{table}[!t]
\centering
\caption{Counterfactual lesion audit on VinDr-RibCXR (1,470 sites): mean absorption of each model with image-cluster bootstrap intervals
(lower is better; ensembles of seeds 0--2). Every single seed is in Supplementary Table~S2.}
\label{tab:audit_means}
\restable{tmi_audit_means}
\end{table}

\begin{table*}[!t]
\centering
\caption{H3$'$: paired site-level contrasts on the three audit sets, Holm over the two contrasts, and the interaction with bone overlap
($\geq$15\,\% vs.\ $<$15\,\%; top). SZCH rows use the 51 held-out images (306 sites) and remain development evidence because the audit on
all 75 images was used during development; the JSRT overlap comes from the reference model's bone map. Bottom: rib benchmark on
VinDr-RibCXR (metrics only); every method is in Supplementary Table~S3.}
\label{tab:audit}
\restable[\textwidth]{tmi_audit}
\subtab{Rib benchmark: gradient-energy reduction on rib-boundary pixels and on other lung pixels, selectivity and AUC}
\restable[\textwidth]{tmi_rib}
\end{table*}

\begin{figure*}[!t]
\centering
\safefig[0.95\textwidth]{figures/audit_maps.pdf}
\caption{The audit on one held-out SZCH-X-Rays image (CC BY-NC 4.0): a 7\,mm nodule from a LIDC-IDRI test CT is added to the radiograph;
the change of each bone estimate (red = lesion signal moved into the bone image) and the resulting soft-tissue image. Not hand-picked: eight
placements over the bone-overlapping lung field were drawn with a fixed seed and the one with the median $A(\text{ours})-A(\text{reference})$
is shown.}
\label{fig:audit_maps}
\end{figure*}

\subsubsection{Rib benchmark}
Absorption must be read beside suppression strength, because doing nothing gives $A=0$ (Table~\ref{tab:audit}, bottom). The reference model
removes more edge energy, both on the rib boundaries (48\,\% against 30\,\% for ours) and on the other lung pixels (22\,\% against 8\,\%);
its selectivity is 0.264 against 0.220, and the AUC separating rib from non-rib pixels is the same (0.671 and 0.672). Ours suppresses ribs
less strongly, but its absorption is a third of the reference model's while its rib-edge removal is about 60\,\% of it, and its bone-region
fidelity is non-inferior (H2$'$).

\subsubsection{Encoders (exploratory)}
The twin effect replicates across three encoders trained under the same recipe (Table~\ref{tab:encoders}). On VinDr-RibCXR the twin lowers
$A$ by 0.127 for ConvNeXt-V2-T, 0.114 for SwinV2-T and 0.112 for RAD-DINO, and on JSRT by 0.073, 0.061 and 0.064. ConvNeXt-V2-T was chosen
as ``ours'' on development data before the final evaluation (Section~\ref{sec:m-encoders}); RAD-DINO reaches the same absorption ($A=0.053$
on VinDr-RibCXR).

\begin{table}[!t]
\centering
\caption{Twin regulariser per encoder (exploratory): absorption difference to the matched no-twin control and its interaction with bone
overlap; three seeds per model (seeds 0--2), SZCH on its 51 held-out images. The single seeds of the final models are in
Supplementary Table~S4.}
\label{tab:encoders}
\restable{tmi_s_encoders}
\end{table}

\subsubsection{Real nodules without lesion insertion (H4$'$a)}
Real nodules give a mixed answer (Table~\ref{tab:cnr}). On JSRT (154 nodules) the contrast-to-noise ratio rises by 0.041 with ours and by
0.018 with the reference model; the difference is $+0.023$ $[0.008,0.041]$, and $+0.032$ $[0.014,0.052]$ for the 123 nodules with at least
15\,\% bone overlap. On Node21 (1,474 boxes) the reference model raises contrast more: $+0.103$ against $+0.078$, a difference of $-0.025$
$[-0.033,-0.018]$, and $-0.017$ $[-0.026,-0.009]$ for the 1,109 bone-overlapping boxes. Against the no-twin control ours raises real-nodule
contrast in all four comparisons ($+0.007$ to $+0.027$, all intervals above zero; Holm within data set $p<0.001$).

\begin{table*}[!t]
\centering
\caption{Contrast-to-noise change of real nodules on JSRT and Node21 (soft-tissue image minus original), all lesions and lesions with at
least 15\,\% bone overlap; image-cluster bootstrap intervals.}
\label{tab:cnr}
\restable[0.85\textwidth]{tmi_cnr}
\end{table*}

% ---------------------------------------------------------------------------------------------------------------------
\subsection{Downstream Detection and Classification (Study~I: H4\texorpdfstring{$'$}{'}b)}
\label{sec:res-downstream}
Our soft-tissue image used alone does not help downstream models (Table~\ref{tab:downstream}, Fig.~\ref{fig:downstream_forest}). On
NIH ChestX-ray14 the macro-AUROC of the three-seed ensemble is 0.818 with the original image and 0.817 with ours ($-0.001$
$[-0.004,0.003]$, patient-cluster bootstrap); per class, ours is lower for effusion, infiltration and pleural thickening ($-0.004$,
$-0.007$, $-0.013$) and higher for nodule ($+0.011$; Holm-adjusted DeLong $p<0.05$ for these four). The reference model's soft image alone
gives $+0.004$ $[0.002,0.007]$. In detection the change in CPM with ours alone is $+0.001$ $[-0.023,0.023]$ on Node21, $+0.006$
$[-0.008,0.020]$ on TBX11K and $-0.009$ $[-0.016,-0.002]$ on VinBigData.

Added as extra channels, our suppressed images do help. The original plus our soft image raises NIH macro-AUROC by $+0.006$
$[0.003,0.008]$ (0.818 to 0.824). The three-channel input (original, soft and bone image) raises the CPM on Node21 by $+0.028$
$[0.009,0.047]$ (0.743 to 0.771) and on VinBigData by $+0.008$ $[0.002,0.014]$ (0.726 to 0.734). The gain is established for CPM only:
mAP@50 changes by $+0.030$ $[-0.002,0.062]$ on Node21 and $-0.001$ $[-0.017,0.012]$ on VinBigData. On TBX11K the three-channel input of ours
gives $+0.008$ $[-0.000,0.019]$, whereas the reference model's soft image alone gains $+0.029$ $[0.016,0.043]$. No difference in differences
of the recall at two false positives per image between bone-overlapping and other lesions excludes zero (Node21 $-0.005$ $[-0.071,0.061]$,
VinBigData $+0.009$ $[-0.034,0.053]$ for our soft image); on TBX11K it cannot be estimated, because every test box overlaps bone by at least
20\,\%.

\begin{table*}[!t]
\centering
\caption{Downstream models. Top: classification on NIH ChestX-ray14, macro-AUROC of the three-seed ensemble, per-seed mean $\pm$ SD and
the difference to the original image (patient-cluster bootstrap). Bottom: detection on Node21, TBX11K and VinBigData, CPM and mAP@50 (mean
$\pm$ SD over three seeds), differences to the original image and the difference in differences (DiD) of recall at two false positives per
image between bone-overlapping and other lesions (image bootstrap). ``Original + ours + bone'': the original, soft-tissue and bone images
as three channels.}
\label{tab:downstream}
\restable[\textwidth]{tmi_down_cls}
\subtab{Detection}
\restable[\textwidth]{tmi_down_det}
\end{table*}

\begin{figure}[!t]
\centering
\safefig[\columnwidth]{figures/downstream_forest.pdf}
\caption{Downstream effects of each input against the original image with 95\,\% intervals: macro-AUROC on NIH ChestX-ray14, CPM on the
detection sets.}
\label{fig:downstream_forest}
\end{figure}

% ---------------------------------------------------------------------------------------------------------------------
\subsection{Fusion Study: Overall Performance, Focal Findings and Localisation (Study~II: H6--H8, H10)}
\label{sec:res-fusion}

\subsubsection{Overall performance (H6, H7)}
All 175 runs finished; Table~\ref{tab:fusion} gives the test metrics of the seed ensembles. The hybrid did not improve on the radiograph
alone on any data set (Fig.~\ref{fig:cls_forest}). On NIH the macro-AUROC difference was $+0.0001$ $[-0.0022,0.0024]$ (Holm $p=0.967$), so
a gain larger than 0.0024 is not compatible with the data; on VinBigData it was $+0.0024$ $[-0.0004,0.0053]$ (Holm $p=0.28$) and on
BDCXR-3257 $-0.0042$ $[-0.0108,0.0021]$ (Holm $p=0.39$). The AUPRC differences agreed in direction and none excluded zero. The hybrid did not
differ from the capacity control on any data set (Holm $p=1$). Against channel stacking it was nominally worse on NIH ($-0.0031$
$[-0.0055,-0.0007]$, Holm $p=0.066$; AUPRC $-0.0055$ $[-0.0100,-0.0012]$) and no different elsewhere. Early fusion had the highest NIH
macro-AUROC of the five configurations (0.8273 against 0.8241 for the radiograph alone); this contrast was not one of the tested
hypotheses and is descriptive. Three-stream fusion therefore did not pay for its three times larger encoder computation, and lightweight
channel stacking did at least as well.

\begin{table*}[!t]
\centering
\caption{Fusion study. Top: test macro-AUROC and macro-AUPRC of the seed ensembles (five seeds; BDCXR-3257: out-of-fold over five folds);
Original: radiograph; Soft: soft-tissue image; Early: three images as channels; Hybrid: three streams; Control: radiograph in three
streams. Bottom: H6 (primary) and H7, differences of the seed ensembles with 95\,\% cluster-bootstrap intervals (\num{10000} resamples);
Holm over the three data sets (H6) and over the six tests of H7.}
\label{tab:fusion}
\restable[\textwidth]{tmi_macro}
\subtab{H6 and H7}
\restable[\textwidth]{tmi_primary}
\end{table*}

\begin{figure*}[!t]
\centering
\safefig[0.62\textwidth]{figures/fig_hypothesis_forest_densenet121.pdf}
\caption{Contrasts of the fusion study with 95\,\% cluster-bootstrap intervals. (a) H6 and H7: macro-AUROC differences per data set.
(b) H8: differences in the mean AUROC of each finding group's classes. Filled markers: Holm $p<0.05$; right: Holm-adjusted $p$.}
\label{fig:cls_forest}
\end{figure*}

\subsubsection{Focal and diffuse findings (H8)}
The effect of the inputs depended on the finding group (Table~\ref{tab:groups}, Fig.~\ref{fig:cls_forest}b). On NIH the soft-tissue image
raised the mean AUROC of the focal findings by $+0.0089$ $[0.0040,0.0140]$ (Holm $p=0.0024$) and the hybrid by $+0.0063$ $[0.0005,0.0117]$
(Holm $p=0.101$). Diffuse findings lost with both inputs (hybrid $-0.0052$ $[-0.0077,-0.0027]$, Holm $p=0.0024$; soft $-0.0032$
$[-0.0057,-0.0008]$, Holm $p=0.039$), and boundary findings did not change. On VinBigData the focal estimates had the same direction and
size (hybrid $+0.0085$ $[-0.0023,0.0196]$; soft $+0.0083$ $[-0.0044,0.0202]$) but wide intervals, the diffuse estimates were close to zero
(hybrid $+0.0011$, soft $+0.0003$), and no group difference survived the correction.

\begin{table}[!t]
\centering
\caption{H8: difference in the mean per-class AUROC of each finding group, with 95\,\% cluster-bootstrap intervals; Holm within data set
over groups and contrasts.}
\label{tab:groups}
\restable{tmi_groups}
\end{table}

Per class (Supplementary Table~S9, Fig.~\ref{fig:cls_per_class}), two NIH classes differed between the hybrid and the radiograph after
correction: nodule gained ($+0.0129$, Holm $p=0.00121$) and infiltration lost ($-0.0160$, Holm $p<0.001$; also $-0.0128$ against channel
stacking). Mass did not change ($-0.0004$). No class differed on VinBigData or BDCXR-3257.

\begin{figure*}[!t]
\centering
\safefig[0.7\textwidth]{figures/fig_auroc_per_class_densenet121.pdf}
\caption{NIH ChestX-ray14: test AUROC per class and per finding group for each configuration, and each configuration minus the
radiograph, with 95\,\% patient-cluster bootstrap bands from 1,000 resamples (display only; the bootstrap tests of
Tables~\ref{tab:fusion} and~\ref{tab:groups} use \num{10000}).}
\label{fig:cls_per_class}
\end{figure*}

\subsubsection{Where the streams point (H10)}
Across 984 targets (164 focal) on 880 test images, Grad-CAM++ maps of the radiograph stream hit focal targets less often than other targets
(0.213 against 0.374), whereas soft-tissue stream maps hit both about equally (0.385 against 0.401) (Table~\ref{tab:pointing};
per class in Supplementary Fig.~S7). The soft-minus-radiograph difference was therefore larger for focal targets ($+0.172$) than for the others
($+0.027$): interaction $+0.145$ $[0.092,0.199]$, $p=0.0002$, so H10 passes. Integrated Gradients ranked the two streams the same way
(focal: 0.580 soft-tissue against 0.505 radiograph; other: 0.621 against 0.645; maps in Supplementary Fig.~S6). Bone-stream maps rarely hit
any target (0.049 and 0.167). Fig.~\ref{fig:cls_maps} shows examples.

\begin{table}[!t]
\centering
\caption{H10: pointing-game hit rates on the NIH expert boxes, means over targets of the per-target means over the five hybrid models;
tolerance 15\,px at 256\,px. Last row: soft-tissue minus radiograph stream (Grad-CAM++) and the interaction with its 95\,\%
image-cluster bootstrap interval.}
\label{tab:pointing}
\restable{v3_pointing}
\end{table}

\begin{figure*}[!t]
\centering
\safefig[0.87\textwidth]{figures/fig_xai_comparison_densenet121.pdf}
\caption{Attribution maps of one hybrid model (seed 0) on NIH ChestX-ray14 test images: the three inputs, Grad-CAM++ per stream and the
decoder attention, with the expert boxes. Rows: targets (row label: the box class).}
\label{fig:cls_maps}
\end{figure*}

% ---------------------------------------------------------------------------------------------------------------------
\subsection{External Validation on Paediatric Radiographs (BDCXR-3257: E1--E3)}
\label{sec:res-bdcxr}
On 3,257 paediatric radiographs from a hospital in Dhaka (Table~\ref{tab:bdcxr}), five-fold cross-validated pneumonia classification
(normal, bacterial, viral) reaches a pooled out-of-fold macro-AUROC of 0.814 on the original image and 0.816 on the soft-tissue image of
ours. The difference, $+0.002$ $[-0.005,0.010]$, meets the non-inferiority margin of 0.01 and shows no superiority (E1); the reference
model's soft image gives $+0.001$ $[-0.007,0.010]$ and CLAHE $-0.006$ $[-0.012,0.001]$, which misses the margin. In the lesion audit on 300
of these images (E2) ours absorbs $A=0.084$, the reference model 0.217 and the no-twin control 0.251 (differences $-0.134$
$[-0.144,-0.124]$ and $-0.167$ $[-0.175,-0.160]$), with a larger advantage at bone-overlapping sites. The inserted lesions come from adult
CTs, so this is a stress test rather than a clinical finding. The ensemble disagreement on these children (E3) is 2.7 times that on the
adult SZCH-X-Rays pairs (median lung-field spread 0.0076 against 0.0028) and is larger in younger children (Spearman with age $-0.22$).

\begin{table*}[!t]
\centering
\caption{BDCXR-3257. Top (E1): pneumonia classification (normal, bacterial, viral), five folds over near-duplicate groups $\times$ three
seeds; pooled out-of-fold macro-AUROC, per-fold mean $\pm$ SD, the difference to the original image (group bootstrap), per-class AUROC and
balanced accuracy. Bottom (E2): lesion audit on a fixed class-stratified subset of 300 images (1,800 sites).}
\label{tab:bdcxr}
\restable[\textwidth]{tmi_s_bdcxr_cls}
\subtab{Lesion audit (E2)}
\restable[0.4\textwidth]{tmi_s_bdcxr_audit}
\end{table*}

% ---------------------------------------------------------------------------------------------------------------------
\subsection{Uncertainty of the Suppressor and the Classifiers (H5\texorpdfstring{$'$}{'}, H9)}
\label{sec:res-uq}

\subsubsection{Reliability of the bone-image uncertainty (H5$'$)}
On the Qure.ai test split with exact truth the split-conformal intervals of the three-model ensemble cover 0.900 $[0.897,0.903]$ of the
pixels at a nominal 0.90, and the ensemble spread ranks the absolute bone error well (Spearman 0.76; Table~\ref{tab:reliability};
per-image coverage in Supplementary Fig.~S1). On the 51 held-out SZCH-X-Rays pairs, measured against a pseudo-truth bone image from the dual-energy pair,
the same intervals cover only 0.139 $[0.134,0.145]$ and the spread ranks the error weakly (Spearman 0.11 $[0.07,0.15]$). The uncertainty is
calibrated in the training domain and does not transfer to real radiographs without recalibration.

\begin{table*}[!t]
\centering
\caption{Bone-image uncertainty. Top (H5$'$): ensemble spread and split-conformal intervals ($K=3$, nominal coverage 0.90), with
image-bootstrap intervals. Bottom (H9, exploratory and post hoc): intervals recalibrated on the 24 SZCH-X-Rays development pairs and
tested on the 51 held-out pairs at nominal 0.90; Mondrian: calibration per stratum of predicted bone overlap and ensemble spread.}
\label{tab:reliability}
\restable[\textwidth]{reliability}
\subtab{Recalibration on real pairs (H9)}
\restable[\textwidth]{v3_bone_intervals}
\end{table*}

\subsubsection{Recalibration on real pairs (H9, post hoc)}
An exploratory analysis, added after this result was known and therefore post hoc on already-seen pairs, recalibrated the intervals on the
24 SZCH-X-Rays development pairs and tested them on the 51 held-out pairs (Table~\ref{tab:reliability}, bottom; Supplementary Fig.~S2).
Pooled coverage was 0.886 with one global calibration and 0.895 with a separate calibration per stratum of predicted bone overlap and
ensemble spread (Mondrian), whose intervals were less than half as wide (mean width 0.215 against 0.447). A few real pairs thus restore
nominal coverage on average, while the coverage of single images still varies.

\subsubsection{Classification uncertainty}
Conformal risk control held the NIH false-negative rate at 0.060--0.069 on the test split for nominal 0.10, but the label sets held
6.33--6.78 of 14 labels on average (Supplementary Table~S10 and Fig.~S4). On BDCXR-3257 the Mondrian sets covered 0.897--0.900 of
the images at nominal 0.90 with 1.76--1.81 of 3 classes per set. In selective classification (Supplementary Fig.~S5), adding the
suppressor's spread to the confidence score raised the AURC in every configuration: NIH 0.0352--0.0358 with the entropy alone against
0.0372--0.0375 with the spread, both on the \num{25580} test images with a lung-field spread (16 of \num{25596} have an empty lung mask),
and BDCXR-3257 0.1754--0.1938 against 0.2138--0.2260. Every single run is shown next to the seed ensembles in Supplementary Fig.~S3.

% ---------------------------------------------------------------------------------------------------------------------
\subsection{Site-Appearance Calibration (Exploratory, H16--H18)}
\label{sec:res-sitecal}
\begin{table*}[!t]
\centering
\caption{Site-appearance calibration (exploratory). Top: fidelity on the 51 held-out SZCH-X-Rays pairs, every method as is and with the
smooth correction of Eq.~\eqref{eq:sitecal} fitted on the 24 development pairs (BS-LDM, trained on this site's pairs, is not corrected).
Middle: H16--H18, paired differences with image bootstrap (H16, H16b, H17: one-sided 95\,\% lower bounds) or image-cluster bootstrap
intervals (absorption on 1,470 VinDr-RibCXR sites), \num{10000} resamples. Bottom: VinDr-RibCXR (metrics only), mean lesion absorption of
one audit run with all four models (shared sites) and the rib benchmark, as is and site-calibrated.}
\label{tab:sitecal}
\restable[\textwidth]{site_cal_fidelity}
\subtab{H16--H18}
\restable[\textwidth]{site_cal_tests}
\subtab{VinDr-RibCXR: absorption and rib benchmark}
\restable[\textwidth]{site_cal_ribcxr}
\end{table*}

The smooth site correction of Section~\ref{sec:m-sitecal}, fitted on the 24 development pairs, raised the bone-region PSNR$_{\mathrm{cal}}$
of ours on the 51 held-out pairs from 21.45 to 23.58\,dB, a paired gain of 2.12\,dB (one-sided 95\,\% lower bound 1.99\,dB; H16 passes),
and its whole-image PSNR$_{\mathrm{cal}}$ from 25.56 to 28.30\,dB (gain 2.74\,dB, lower bound 2.61\,dB; H16b passes;
Table~\ref{tab:sitecal}). Every CT-supervised model gained about as much (reference model 21.51 to 23.92\,dB, v1 model 21.95 to 24.06\,dB,
no-twin control 21.45 to 23.74\,dB) and the unprocessed radiograph less (20.21 to 21.46\,dB): the correction calibrates the site rather
than improving a suppressor, and the order of the CT-supervised models did not change. After calibration ours remained non-inferior to the
calibrated reference model, by a narrower margin than before ($-0.35$\,dB, lower bound $-0.44$\,dB against the margin of $-0.5$\,dB; H17
passes), and the no-twin control was slightly ahead of ours (23.74 against 23.58\,dB; not tested).

The calibration did not cost lesion fidelity (Table~\ref{tab:sitecal}, bottom): on VinDr-RibCXR the calibrated ours absorbed 0.036 of an
inserted lesion and the calibrated reference model 0.132 (difference $-0.096$ $[-0.105,-0.087]$; H18 passes), and the calibration lowered
the absorption of both models slightly ($-0.015$ and $-0.017$); the uncalibrated models of this run reproduced the audit of
Table~\ref{tab:audit_means}. It weakened suppression a little and made it more selective: the rib-edge reduction of ours fell from 0.298 to
0.282 and the reduction on other lung pixels from 0.078 to 0.050, so the selectivity rose from 0.220 to 0.232. BS-LDM, trained on this
site's dual-energy pairs, remains about 5\,dB ahead in the bone region (28.57 against 23.58\,dB) and 3.65\,dB in the whole image.


\FloatBarrier
\section{Discussion}
\label{sec:discussion}

\subsection{Principal Findings, Lesion Fidelity and the Cost of the Regulariser}
\label{sec:d-principal}
CT-supervised bone suppression can be made lesion-faithful, and its images help prediction in a specific way. Lesion absorption fell
steadily with the twin weight (Table~\ref{tab:twin}); on an audit set that no model had seen during development, the twin ensemble absorbed
a third of what the released CT-supervised model absorbed and less than a third of what the matched control absorbed
(Tables~\ref{tab:audit_means} and~\ref{tab:audit}), at non-inferior bone-region fidelity; and the effect replicated with two other encoders
(Table~\ref{tab:encoders}) and on paediatric radiographs (Section~\ref{sec:res-bdcxr}). Added as extra channels, the suppressed images
raised nodule detection on Node21, detection on VinBigData and classification on NIH (Table~\ref{tab:downstream}), and in the fusion study
they favoured focal findings, both in AUROC and in where the attribution maps pointed (Tables~\ref{tab:groups} and~\ref{tab:pointing}).

Several planned results did not hold. Our soft-tissue image alone did not improve any downstream task (Table~\ref{tab:downstream}); a
multi-stream classifier that read the radiograph, soft-tissue and bone images did not improve on the radiograph alone (H6,
Table~\ref{tab:fusion}); SSIM fell short of the reference model by more than its margin (Table~\ref{tab:fidelity}); real-nodule contrast was
higher than with the reference model on JSRT and lower on Node21 (Table~\ref{tab:cnr}); and conformal intervals calibrated on synthetic
data covered 14\,\% instead of 90\,\% of real bone pixels, 89\,\% after an exploratory recalibration on 24 real pairs
(Table~\ref{tab:reliability}).

\emph{What the audit adds.} Fidelity metrics and lesion absorption rank suppressors differently (Fig.~\ref{fig:pareto}). The v1 model has
the highest bone-region PSNR of the CT-supervised models and the highest absorption; ours and the reference model are within 0.06\,dB of
each other and a factor of three apart in absorption. BS-LDM, trained on dual-energy pairs from the evaluation hospital, is about 7\,dB ahead
in fidelity on that hospital's images, yet absorbed 0.42 of an inserted lesion in the development audit (at unpaired sites, so only
descriptive): fidelity to a dual-energy target does not protect what the target never showed. Part of that lead is the appearance of the
site's dual-energy images rather than suppression: a smooth calibration learned from 24 pairs of the site raised every CT-supervised model
by 2.1 to 2.4\,dB in the bone region without raising lesion absorption, and BS-LDM remained about 5\,dB ahead
(Section~\ref{sec:res-sitecal}; exploratory). Absorption must still be read beside suppression strength. The reference model removes more
rib-edge energy than ours (48\,\% against 30\,\%) and more non-rib energy (22\,\% against 8\,\%), so part of its higher absorption may
reflect stronger suppression in general.

\emph{What the regulariser costs.} On real paired radiographs the twin term costs no measurable bone-region fidelity up to $w=16$ (21.22 to
21.28\,dB across $w=0\ldots16$); at $w=32$ and 64 PSNR drops by about 0.3\,dB, so $w=16$, chosen on development data before the final
evaluation, is the usable operating point. The cost appears where the truth is exact: on the Qure.ai test split the bone error of the twin
ensemble is 41\,\% higher than that of the control (MAE 0.0041 against 0.0029; Supplementary Table~S5). The regulariser thus trades
synthetic bone accuracy for lesion invariance, a trade not visible on real images. The ablation shows that the effect needs gradients
through both twin branches and a consistency term confined to the pasted lesion; showing lesions during training without the twin term
(CLIA) did not lower absorption. The reference model is a SwinV2-B at $1024$~px trained with a far larger budget, so only the matched
no-twin control, which absorbs slightly more than the reference model (0.178 against 0.149), isolates the regulariser. After the
exploratory site calibration the no-twin control was slightly ahead of ours (23.74 against 23.58\,dB), a cost of the regulariser that the
uncalibrated comparison did not show; it was not tested. Diffusion models trained on real dual-energy pairs
\cite{sun2026bsldm,sun2026debonedit} need paired data from the target site, sample slowly (about 2.5 minutes per image for BS-LDM on a T4
GPU) and were not designed for lesion invariance.

\subsection{Why Suppressed Inputs Help Focal Findings}
\label{sec:d-focal}
The suppressed images are deterministic functions of the radiograph (Eq.~\eqref{eq:info}): the suppressor adds a transformation learned
from CT but no information, and a classifier trained on thousands of radiographs may learn whatever bone-invariant features its labels
reward. The results fit a modest, cheap prior rather than new information. Extra channels helped, whereas the soft image alone did not; the
multi-stream hybrid matched the radiograph alone on all three data sets, and on NIH its interval rules out a gain larger than 0.0024 in
macro-AUROC, so the null result is informative. The hybrid also matched the capacity control, and plain channel stacking did at least as
well as separate streams at a third of the encoder computation. Nodules and masses had the largest estimates on both adult data sets
($+0.008$ to $+0.009$ for the soft-tissue input), and the pointing game agreed: maps of the radiograph stream found focal lesions less often
than other findings, whereas maps of the soft-tissue stream found both about equally often. Not every signal favoured ours: on TBX11K the
reference model's soft image helped more, and on Node21 it raised real-nodule contrast more, so our less strongly suppressing model may
trade some contrast for lesion fidelity.

Three cautions apply. The NIH test split had already been used by the downstream classifiers of Study~I, so the NIH result of the fusion
study is a re-test with new models; VinBigData, the independent test, has a single merged nodule/mass class whose interval included zero;
and the detection gains are established for CPM, whereas the mAP@50 intervals include zero. The loss for diffuse findings appeared on NIH
only, mostly for infiltration; with report-derived NIH labels \cite{wang2017chestxray} and no such loss on the radiologist-labelled
VinBigData, these data cannot tell whether suppression removes cues for diffuse disease or interacts with label noise.

\subsection{Audit Circularity, Uncertainty and Clinical Translation}
\label{sec:d-translation}
\emph{Circularity.} The regulariser pastes lesions to teach invariance and the audit pastes lesions to measure it. The design limits the
circularity (Section~\ref{sec:m-audit}): audit lesions come from held-out LIDC-IDRI test CTs, audit images are real radiographs, the
confirmatory audit uses rib masks that no model has seen, and evidence that does not use insertion points the same way, namely the contrast
of real JSRT nodules (ours above the reference model and the no-twin control) and the Node21 detection gain of the three-channel input. These
choices limit the circularity but cannot remove it: a model that recognises pasted signal in general could still score well.

\emph{Uncertainty.} The bone-image intervals held where they were calibrated and failed under the shift to real radiographs; 24 real pairs
from the target site sufficed to restore coverage here, although transfer between sites is untested. For the classifiers, conformal risk
control kept its guarantee for every input, but label sets holding about half of the 14 NIH labels are of little use to a reader, and the
suppressor's own uncertainty made abstention worse in every configuration: uncertainty about the bone image is not uncertainty about the
diagnosis.

\emph{Recommendations.} A bone suppressor should be audited for lesion absorption, reported next to a suppression measure, before it is
used; suppressed images should be shown beside the original rather than instead of it; and the ensemble spread should serve as a warning
flag rather than a calibrated interval unless recalibrated on real pairs (it rose 2.7-fold on paediatric radiographs). Detection and
classification models should receive suppressed images as extra channels next to the radiograph rather than as separate streams, should
not be expected to gain more than small amounts of overall accuracy, and should be tested per finding against a control of the same
capacity. Fidelity comparisons between CT-supervised and dual-energy-trained suppressors should separate the appearance of the test site
from suppression, for example with a calibration on a few development pairs reported beside the uncorrected numbers.

\subsection{Limitations and Future Research}
\label{sec:limitations}
Each limitation is stated with its consequence for how the results must be read.
\begin{enumerate}
\item \emph{Synthetic supervision and one source of real pairs.} The bone image of a real radiograph is never observed. Real-image fidelity
 is measured against dual-energy soft-tissue images from one hospital and one acquisition set-up (GE Discovery XR656), which differ from
 the radiograph in tone and fine texture, so raw and band-limited metrics are both reported and calibrated metrics are not absolute.
\item \emph{The audit is linear and additive, and shares its mechanism with the regulariser.} Lesions are added in image units; beam
 hardening, scatter and tone mapping are not modelled, so $A$ measures the leak of an added signal rather than clinical detectability, and
 it is read next to a suppression gain because $A=0$ is reached by doing nothing. The real-nodule contrast and downstream subgroup analyses
 limit, but do not remove, the circularity of inserting lesions both in training and in the audit.
\item \emph{Development use of SZCH.} All 75 SZCH images were used in development (audit) and 24 for paired decisions, so only the 51
 held-out pairs support fidelity claims; the recipe was revised once after the SZCH pairs had been looked at. The exploratory site
 calibration was designed after H2$'$ was known, re-uses the 51 held-out pairs and is specific to one hospital and scanner, whose own 24
 development pairs it was fitted on.
\item \emph{Resolution, backbone and operating point.} Our suppressors predict at 512\,px with small encoders, the reference model at
 1024\,px with SwinV2-B. The SSIM gap sits in fine texture (with a band limit of $\sigma=4$ it shrinks from 0.021 to 0.002); whether
 prediction at 1024\,px closes it was not tested. Only $w=16$ was evaluated with the full protocol.
\item \emph{Uncertainty.} Conformal calibration treats pixels as exchangeable and was calibrated on synthetic data; the recalibration on
 real pairs was chosen after the shift was seen and used 24 pairs from one hospital. Ensemble members share training data.
\item \emph{Scope and external validity.} Training and paired evaluation use adult frontal radiographs; BDCXR-3257 is a single-centre
 paediatric set with image-level labels and no patient identifiers, and its audit inserts adult nodules into children's radiographs. There
 is no lateral imaging and no reader study. On JSRT the bone overlap of each audit site came from the reference model's bone map.
 Downstream models were not tuned per input, the overlap analysis is not estimable on TBX11K, and the downstream gains are small and were
 measured with one classifier and one detector. The models are for research use only.
\item \emph{Fusion study.} One encoder (DenseNet-121) and one resolution ($256\times256$\,px) were tested; heavier encoders or full
 resolution may use the extra streams differently. The NIH test split served the downstream experiments of Study~I as well, its labels are
 report-derived, and VinBigData has one focal class. The per-class DeLong tests treat NIH images as independent although patients
 contribute several images, and the pointing game depends on the attribution method, the 15\,px tolerance and the eight annotated NIH
 classes.
\end{enumerate}

\emph{Future research.} Future work should test prediction at $1024$~px, recalibration on real pairs from more sites, other lesion types,
detectors that receive the suppressed images as extra channels at full resolution, a reader study and a prospective evaluation.

\section{Conclusion}
\label{sec:conclusion}
Bone suppressors that match a dual-energy reference equally well differed threefold in how much of a lesion they moved into their bone
image. A counterfactual audit makes this loss measurable on real radiographs, and a twin lesion-invariance regulariser trained on public
CT-derived projections lowered it to a third of a released CT-supervised model's, at non-inferior bone-region fidelity, across encoders
and on paediatric radiographs. The suppressed images are worth using in a specific way: added as extra channels they improved nodule
detection on Node21 and classification on NIH, and they favoured focal findings, in AUROC and in the localisation of Grad-CAM++ maps,
whereas separate streams gave no overall gain and diffuse findings lost. Lesion absorption should therefore be audited before
CT-supervised suppression is used on radiographs, the ensemble spread treated as a warning flag unless recalibrated on real pairs, and
gains from suppressed inputs expected, and tested against a control of equal capacity, for specific findings rather than overall.

% Appendix (IEEE: after the conclusion, before the acknowledgment).
\appendices
\section{First-Generation (v1) Training Images Rendered from LIDC-IDRI}
\label{app:v1}
The first (v1) suppressors, the lesion library and the synthetic test CTs come from LIDC-IDRI. This appendix describes how a CT is split
into exact components and projected, how a training sample is composed, and how lesions were used in v1.

\emph{CT decomposition.} Each CT was split into bone, non-lung soft tissue and lung. The body is the largest component of $\mathrm{HU}>-300$
(slice-wise hole filling, then a $1{\times}9{\times}9$ opening to detach the couch). Bone follows the iterative-relaxation recipe of the
reference method \cite[Algorithm~1]{angaitkar2026anatomy}, with a 10~mm support cube, the first pass starting from the reference mask and a
single closing after the ladder: a reference mask at 300~HU, thresholds $\tau\in\{200,150,100\}$~HU admitted only where at least 5\,\% of a
10~mm neighbourhood was already bone, union with the reference mask, a radius-2 closing to fill marrow, and removal of components smaller
than 1\,\% of the largest. Lung is air ($<-320$~HU) inside the body in components at least 20\,\% the size of the largest, closed and
hole-filled slice-wise (so vessels and nodules belong to the lung component) and without bone. Bone voxels were replaced in the
soft-tissue component by a normalised convolution of the surrounding soft tissue (Gaussian, 4~mm, with 12 and 36~mm fall-backs). Density is
water-equivalent, $\rho=\mathrm{clip}(1+\mathrm{HU}/1000,0,3)$~g\,cm$^{-3}$.

\emph{Polychromatic projection.} Each component was integrated along the anteroposterior axis (parallel beam), giving an areal density
$a_c$ (g\,cm$^{-2}$). The frame is a body crop: width 1.08 times the body width, height the lung extent plus 60~mm clipped to 0.75--1.2 of
the width (short scans cannot reach 1.2), area-resized to $512\times512$ and later squashed to a square exactly as real radiographs are.
Padded rows are set to air. The AP view is the horizontal mirror of the PA view. For a component $c$ with areal density $a_c$, a tube
spectrum $\phi(E)$ and an energy-integrating detector, a pixel value is
\begin{equation}
P_c=-\ln\frac{\sum_E \phi(E)\,E\,\exp\!\big(-(\mu/\rho)_c(E)\,a_c\big)}{\sum_E\phi(E)\,E},
\label{eq:physics}
\end{equation}
with ICRU-44 \cite{icru1989report44} cortical bone, soft tissue and lung attenuation (Table~4 of NIST SRD 126 \cite{hubbell2004xraymass})
and W-anode spectra (SpekPy \cite{poludniowski2021spekpy}, 12$^\circ$, 2.5~mm Al) at 70--140~kVp (training draws 80--130~kVp, the
synthetic evaluation uses 100~kVp). Components are stored as areal densities and the spectrum is applied at training time, so one render
serves every tube voltage. The sum of separately projected components differs from the projection of the whole body because of beam
hardening; we report this additivity gap (6.1\,\% at 100~kVp for 20~g\,cm$^{-2}$ of soft tissue and 3~g\,cm$^{-2}$ of bone).
Fig.~\ref{fig:pipeline} shows the decomposition and projection of one CT.

\begin{figure*}[!t]
\centering
\safefig[0.7\textwidth]{figures/fig_preprocessing_pipeline.pdf}
\caption{CT decomposition and polychromatic projection for one LIDC-IDRI test CT. (a) CT slice; (b) bone, lung and non-lung soft-tissue
masks; (c)--(e) areal densities of bone, non-lung soft tissue and lung; (f) composite radiograph $X$ at 100~kVp; (g) exact bone image $B$;
(h) exact soft-tissue image $S=X-B$; (i) a real radiograph (SZCH-X-Rays) for comparison.}
\label{fig:pipeline}
\end{figure*}

\emph{Training samples.} A training sample is composed on the GPU from the stored areal densities in this order: lesion insertion, one
random affine transform shared by all maps (rotation $\pm15^\circ$, scale 0.9--1.15, shift $\pm6\,\%$, mirror), the physics of
Eq.~\eqref{eq:physics}, component weighting and a monotone tone curve $\gamma$:
\begin{equation}
\begin{gathered}
S=\mathrm{norm}\big(P_{\mathrm{nl}}+g\,P_{\mathrm{lung}}\big),\quad B=\frac{w_b\,P_{\mathrm{bone}}}{p_{99}(P_{\mathrm{bone}})},\\
X=\frac{S+B}{1.5},\quad X'=\gamma(X),\quad S'=\gamma(S),\quad B'=X'-S',
\end{gathered}
\label{eq:compose}
\end{equation}
with $g\sim U(1,3)$ and $w_b\sim U(0.2,0.5)$. The tone curve is either a random gamma or histogram matching to the intensity quantiles of
a randomly drawn real radiograph (4\,000 NIH ChestX-ray14 training-list images). Defining $B'=X'-S'$ keeps the decomposition exact under
any monotone mapping. Gaussian blur is linear and applied identically to $X'$ and $B'$; quantum noise and scatter modify only $X'$. The
region of interest of the suppression loss (Eq.~\eqref{eq:loss}) comes from the lung component (top structures:
$P_{\mathrm{lung}}/p_{99.9}>0.3$). Unlike the Qure.ai projections, where lesions are inserted after the affine transform, lesions enter the
LIDC renders before the affine transform and the physics.

\emph{Lesion-faithful supervision in v1.} Each lesion (LIDC consensus of at least two readers, 50\,\% voxel agreement) was projected as its
excess density over the surrounding lung. The \emph{lesion-weighted loss} adds $(\mathrm{L1}+\mathrm{MSE})$ on a dilated lesion mask with
weight 1. \emph{Counterfactual lesion insertion (CLIA)} adds, before the physics, 0--3 lesion patches drawn from \emph{training} CTs at
random lung positions (including over bone) with random scale (0.7--1.5) and an amplitude of 0.6--1.6 times that scale; the bone target is
untouched, so the network is explicitly taught that a lesion on a rib is not bone. The v1 hypothesis H1 compares these arms on synthetic
LIDC test CTs with exact truth (Section~\ref{sec:res-fidelity}).

\section*{Acknowledgment}
We thank the providers of the data used here: Qure.ai for the CT projections and the released bone-suppression weights, the CT-RATE team,
The Cancer Imaging Archive and the LIDC-IDRI investigators, the authors of SZCH-X-Rays, VinDr (VinDr-RibCXR and VinDr-CXR), the Japanese
Society of Radiological Technology (JSRT), the Node21, TBX11K and NIH ChestX-ray14 teams, and Dr.~M~R~Khan Shishu Hospital \& Institute of
Child Health, Dhaka, for BDCXR-3257 (ethical approval REC-FSIT-2025/No:~11928). \todo{funding and colleagues}

\section*{Ethics Statement}
This is a retrospective study of existing images. Public data sets were used under their licences, VinDr-RibCXR under its data use
agreement, and BDCXR-3257 under the ethical approval REC-FSIT-2025/No:~11928 of Daffodil International University and a data use agreement
with the providing hospital. No new patient data were collected, and no image of a licensed or restricted set is shown.

\section*{Data and Code Availability}
The public data sets are available from their providers (Table~\ref{tab:data}); the Qure.ai projections are built from CT-RATE volumes
\cite{hamamci2026ctrate} and are available under CC BY-NC-SA~4.0, and no image derived from them is shown. VinDr-RibCXR was used under its
data use agreement and is not redistributed. BDCXR-3257 is public on figshare under CC BY 4.0 \cite{durjoy2026bdcxr}; the authors obtained
it under a data use agreement with the providing hospital and follow its terms, so no image of it is shown and nothing derived from it is
redistributed. VinBigData was processed only in cloud accounts that accepted the competition rules. Only metrics left these environments,
and no image, mask, map or model derived from VinDr-RibCXR, VinBigData or BDCXR-3257 images is released.

\emph{Compute.} Training used NVIDIA T4 GPUs in cloud notebooks: a run of \num{300000} samples took 6 to 7 GPU-hours, a screening run of
\num{100000} samples about 2 GPU-hours. The locked-set evaluation ran on a consumer GPU and on cloud T4 GPUs with identical code; the audit
sites of every model were pinned to the same lung masks, and model outputs differ between the two GPU types by at most 0.0024\,dB.

\emph{Release.} The code is available under the MIT licence at \url{https://github.com/HTanvir0011/Thesis-full-project}; the three
checkpoints of the ensemble and a model card are released there for research use under CC BY-NC-SA~4.0, the licence of the training
projections, as release v1.0.0 with SHA-256 checksums. The reporting follows the
CLAIM 2024 checklist \cite{tejani2024claim}, which is submitted with the manuscript.

\section*{CRediT Authorship Contribution Statement}
\todo{author contributions in CRediT terms, for example conceptualisation, methodology, software, formal analysis, writing}

\section*{Declaration of Competing Interest}
\todo{competing interests, or: The authors declare that they have no known competing financial interests or personal relationships that
could have appeared to influence the work reported in this paper.}

\section*{Declaration of Generative AI Use}
During the preparation of this work the authors used an AI assistant (Claude, Anthropic) for code, analysis automation and drafting of the
text. After using this tool, the authors reviewed and edited the content as needed and take full responsibility for the content of the
publication. \todo{confirm or adapt this statement to the actual use before submission}


% References: thebibliography frozen from BibTeX (see README.txt to regenerate).
% Reference list frozen from BibTeX (refs.bib, style vancouver-doi.bst), numbered in order of first citation.
% To regenerate after adding citations: in main.tex, replace % Reference list frozen from BibTeX (refs.bib, style vancouver-doi.bst), numbered in order of first citation.
% To regenerate after adding citations: in main.tex, replace % Reference list frozen from BibTeX (refs.bib, style vancouver-doi.bst), numbered in order of first citation.
% To regenerate after adding citations: in main.tex, replace \input{sections/references} by
% \bibliographystyle{vancouver-doi} \bibliography{refs}, compile, and run tools/freeze_refs.py on main.bbl.
\begin{thebibliography}{10}

\bibitem{kuhlman2006des}
Kuhlman JE, Collins J, Brooks GN, Yandow DR, Broderick LS.
\newblock Dual-Energy Subtraction Chest Radiography: What to Look for beyond
  Calcified Nodules.
\newblock RadioGraphics. 2006;26(1):79-92.
\newblock \url{https://doi.org/10.1148/rg.261055034}.

\bibitem{sun2026bsldm}
Sun Y, Chen Z, Zheng H, Deng W, Liu J, Min W, et~al.
\newblock {BS-LDM}: Effective Bone Suppression in High-Resolution Chest {X}-Ray
  Images With Conditional Latent Diffusion Models.
\newblock {IEEE} Journal of Biomedical and Health Informatics.
  2026;30(1):376-87.
\newblock Online 14 July 2025.
\newblock \url{https://doi.org/10.1109/JBHI.2025.3588138}.

\bibitem{freedman2011lung}
Freedman MT, Lo SCB, Seibel JC, Bromley CM.
\newblock Lung Nodules: Improved Detection with Software That Suppresses the
  Rib and Clavicle on Chest Radiographs.
\newblock Radiology. 2011;260(1):265-73.
\newblock \url{https://doi.org/10.1148/radiol.11100153}.

\bibitem{schalekamp2014cad}
Schalekamp S, van Ginneken B, Koedam E, Snoeren MM, Tiehuis AM, Wittenberg R,
  et~al.
\newblock Computer-aided Detection Improves Detection of Pulmonary Nodules in
  Chest Radiographs beyond the Support by Bone-suppressed Images.
\newblock Radiology. 2014;272(1):252-61.
\newblock \url{https://doi.org/10.1148/radiol.14131315}.

\bibitem{rajaraman2021bone}
Rajaraman S, Zamzmi G, Folio L, Alderson P, Antani S.
\newblock Chest {X}-ray Bone Suppression for Improving Classification of
  Tuberculosis-Consistent Findings.
\newblock Diagnostics. 2021;11(5):840.
\newblock \url{https://doi.org/10.3390/diagnostics11050840}.

\bibitem{chen2024bsdiff}
Chen Z, Sun Y, Qin W, Ge R, Pan C, Deng W, et~al. {BS-Diff}: Effective Bone
  Suppression Using Conditional Diffusion Models from Chest {X}-Ray Images;
  2024.
\newblock Accepted at IEEE ISBI 2024.
\newblock arXiv preprint arXiv:2311.15328.
\newblock \url{https://doi.org/10.48550/arXiv.2311.15328}.

\bibitem{sun2025gllcm}
Sun Y, Chen Z, Zheng H, Lu Y, Duan L, Fan F, et~al. {GL-LCM}: Global-Local
  Latent Consistency Models for Fast High-Resolution Bone Suppression in Chest
  {X}-Ray Images; 2025.
\newblock Accepted at MICCAI 2025.
\newblock arXiv preprint arXiv:2508.03357.
\newblock \url{https://doi.org/10.48550/arXiv.2508.03357}.

\bibitem{sun2026debonedit}
Sun Y, Fan F, Jia J, Deng W, Xu H, Wang C, et~al. {DeBoneDiT}: Depth-Driven
  Conditional Bridge Diffusion Transformers for Bone Suppression in Chest
  {X}-Ray Images; 2026.
\newblock arXiv preprint arXiv:2610.19466.
\newblock \url{https://doi.org/10.48550/arXiv.2610.19466}.

\bibitem{angaitkar2026anatomy}
Angaitkar M, Kumar P, Satia A, Rao P, Mittal A, Tadepalli M, et~al.
  Anatomy-Decomposed Chest Computed Tomography ({CT}) Projections as Scalable
  Supervision for Bone Suppression in Chest Radiographs; 2026.
\newblock Preprint.
\newblock arXiv preprint arXiv:2609.24937.
\newblock \url{https://doi.org/10.48550/arXiv.2609.24937}.

\bibitem{hamamci2026ctrate}
Hamamci IE, Er S, Wang C, Almas F, Simsek AG, Esirgun SN, et~al.
\newblock Generalist foundation models from a multimodal dataset for {3D}
  computed tomography.
\newblock Nature Biomedical Engineering. 2026;10(8):1610-28.
\newblock Preprint: arXiv:2403.17834 (2024).
\newblock \url{https://doi.org/10.1038/s41551-025-01599-y}.

\bibitem{hogeweg2013suppression}
Hogeweg L, Sanchez CI, van Ginneken B.
\newblock Suppression of Translucent Elongated Structures: Applications in
  Chest Radiography.
\newblock {IEEE} Transactions on Medical Imaging. 2013;32(11):2099-113.
\newblock \url{https://doi.org/10.1109/TMI.2013.2274212}.

\bibitem{takagi2020dose}
Takagi S, Yaegashi T, Ishikawa M.
\newblock Dose reduction and image quality improvement of chest radiography by
  using bone-suppression technique and low tube voltage: a phantom study.
\newblock European Radiology. 2020;30(1):571-80.
\newblock \url{https://doi.org/10.1007/s00330-019-06375-6}.

\bibitem{cohen2018distribution}
Cohen JP, Luck M, Honari S.
\newblock Distribution Matching Losses Can Hallucinate Features in Medical
  Image Translation.
\newblock In: Medical Image Computing and Computer Assisted Intervention --
  {MICCAI} 2018. Lecture Notes in Computer Science. Springer; 2018. p. 529-36.
\newblock Available from: \url{https://doi.org/10.1007/978-3-030-00928-1_60}.

\bibitem{antun2020instabilities}
Antun V, Renna F, Poon C, Adcock B, Hansen AC.
\newblock On instabilities of deep learning in image reconstruction and the
  potential costs of {AI}.
\newblock Proceedings of the National Academy of Sciences.
  2020;117(48):30088-95.
\newblock \url{https://doi.org/10.1073/pnas.1907377117}.

\bibitem{bhadra2021hallucinations}
Bhadra S, Kelkar VA, Brooks FJ, Anastasio MA.
\newblock On Hallucinations in Tomographic Image Reconstruction.
\newblock {IEEE} Transactions on Medical Imaging. 2021;40(11):3249-60.
\newblock \url{https://doi.org/10.1109/TMI.2021.3077857}.

\bibitem{tivnan2024hallucination}
Tivnan M, Yoon S, Chen Z, Li X, Wu D, Li Q.
\newblock Hallucination Index: An Image Quality Metric for Generative
  Reconstruction Models.
\newblock In: Medical Image Computing and Computer Assisted Intervention --
  {MICCAI} 2024. Lecture Notes in Computer Science. Springer; 2024. p. 449-58.
\newblock Available from: \url{https://doi.org/10.1007/978-3-031-72117-5_42}.

\bibitem{varoquaux2022failures}
Varoquaux G, Cheplygina V.
\newblock Machine learning for medical imaging: methodological failures and
  recommendations for the future.
\newblock npj Digital Medicine. 2022;5:48.
\newblock \url{https://doi.org/10.1038/s41746-022-00592-y}.

\bibitem{simmons2011falsepositive}
Simmons JP, Nelson LD, Simonsohn U.
\newblock False-Positive Psychology: Undisclosed Flexibility in Data Collection
  and Analysis Allows Presenting Anything as Significant.
\newblock Psychological Science. 2011;22(11):1359-66.
\newblock \url{https://doi.org/10.1177/0956797611417632}.

\bibitem{suzuki2006mtann}
Suzuki K, Abe H, MacMahon H, Doi K.
\newblock Image-processing technique for suppressing ribs in chest radiographs
  by means of massive training artificial neural network ({MTANN}).
\newblock {IEEE} Transactions on Medical Imaging. 2006;25(4):406-16.
\newblock \url{https://doi.org/10.1109/TMI.2006.871549}.

\bibitem{yang2017cascade}
Yang W, Chen Y, Liu Y, Zhong L, Qin G, Lu Z, et~al.
\newblock Cascade of multi-scale convolutional neural networks for bone
  suppression of chest radiographs in gradient domain.
\newblock Medical Image Analysis. 2017;35:421-33.
\newblock \url{https://doi.org/10.1016/j.media.2016.08.004}.

\bibitem{woo2023convnextv2}
Woo S, Debnath S, Hu R, Chen X, Liu Z, Kweon IS, et~al.
\newblock {ConvNeXt V2}: Co-designing and Scaling {ConvNets} with Masked
  Autoencoders.
\newblock In: Proceedings of the {IEEE/CVF} Conference on Computer Vision and
  Pattern Recognition ({CVPR}); 2023. p. 16133-42.
\newblock \url{https://doi.org/10.1109/CVPR52729.2023.01548}.

\bibitem{liu2022swinv2}
Liu Z, Hu H, Lin Y, Yao Z, Xie Z, Wei Y, et~al.
\newblock Swin Transformer {V2}: Scaling Up Capacity and Resolution.
\newblock In: Proceedings of the {IEEE/CVF} Conference on Computer Vision and
  Pattern Recognition ({CVPR}); 2022. p. 11999-2009.
\newblock \url{https://doi.org/10.1109/CVPR52688.2022.01170}.

\bibitem{cheng2020addressing}
Cheng K, Caliv{\'a} F, Shah R, Han M, Majumdar S, Pedoia V.
\newblock Addressing The False Negative Problem of Deep Learning {MRI}
  Reconstruction Models by Adversarial Attacks and Robust Training.
\newblock In: Proceedings of the Third Conference on Medical Imaging with Deep
  Learning. vol. 121 of Proceedings of Machine Learning Research. PMLR; 2020.
  p. 121-35.
\newblock Available from:
  \url{https://proceedings.mlr.press/v121/cheng20a.html}.

\bibitem{li2021assessing}
Li K, Zhou W, Li H, Anastasio MA.
\newblock Assessing the Impact of Deep Neural Network-Based Image Denoising on
  Binary Signal Detection Tasks.
\newblock {IEEE} Transactions on Medical Imaging. 2021;40(9):2295-305.
\newblock \url{https://doi.org/10.1109/TMI.2021.3076810}.

\bibitem{liu2021query2label}
Liu S, Zhang L, Yang X, Su H, Zhu J. {Query2Label}: A Simple Transformer Way to
  Multi-Label Classification; 2021.
\newblock arXiv preprint arXiv:2107.10834.
\newblock \url{https://doi.org/10.48550/arXiv.2107.10834}.

\bibitem{ridnik2023mldecoder}
Ridnik T, Sharir G, Ben-Cohen A, Ben-Baruch E, Noy A.
\newblock {ML-Decoder}: Scalable and Versatile Classification Head.
\newblock In: 2023 IEEE/CVF Winter Conference on Applications of Computer
  Vision (WACV); 2023. p. 32-41.
\newblock \url{https://doi.org/10.1109/WACV56688.2023.00012}.

\bibitem{chattopadhay2018gradcampp}
Chattopadhay A, Sarkar A, Howlader P, Balasubramanian VN.
\newblock {Grad-CAM++}: Generalized Gradient-Based Visual Explanations for Deep
  Convolutional Networks.
\newblock In: 2018 IEEE Winter Conference on Applications of Computer Vision
  (WACV); 2018. p. 839-47.
\newblock \url{https://doi.org/10.1109/WACV.2018.00097}.

\bibitem{sundararajan2017axiomatic}
Sundararajan M, Taly A, Yan Q. Axiomatic Attribution for Deep Networks; 2017.
\newblock arXiv preprint arXiv:1703.01365.
\newblock \url{https://doi.org/10.48550/arXiv.1703.01365}.

\bibitem{zhang2018topdown}
Zhang J, Bargal SA, Lin Z, Brandt J, Shen X, Sclaroff S.
\newblock Top-Down Neural Attention by Excitation Backprop.
\newblock International Journal of Computer Vision. 2018;126(10):1084-102.
\newblock \url{https://doi.org/10.1007/s11263-017-1059-x}.

\bibitem{lakshminarayanan2017ensembles}
Lakshminarayanan B, Pritzel A, Blundell C.
\newblock Simple and Scalable Predictive Uncertainty Estimation using Deep
  Ensembles.
\newblock In: Advances in Neural Information Processing Systems. vol.~30; 2017.
  Available from:
  \url{https://proceedings.neurips.cc/paper_files/paper/2017/hash/9ef2ed4b7fd2c810847ffa5fa85bce38-Abstract.html}.

\bibitem{lei2018conformal}
Lei J, G'Sell M, Rinaldo A, Tibshirani RJ, Wasserman L.
\newblock Distribution-Free Predictive Inference for Regression.
\newblock Journal of the American Statistical Association.
  2018;113(523):1094-111.
\newblock \url{https://doi.org/10.1080/01621459.2017.1307116}.

\bibitem{angelopoulos2023conformal}
Angelopoulos AN, Bates S.
\newblock Conformal Prediction: A Gentle Introduction.
\newblock Foundations and Trends in Machine Learning. 2023;16(4):494-591.
\newblock Preprint title: A Gentle Introduction to Conformal Prediction and
  Distribution-Free Uncertainty Quantification.
\newblock \url{https://doi.org/10.1561/2200000101}.

\bibitem{angelopoulos2022crc}
Angelopoulos AN, Bates S, Fisch A, Lei L, Schuster T. Conformal Risk Control;
  2022.
\newblock arXiv preprint arXiv:2208.02814.
\newblock \url{https://doi.org/10.48550/arXiv.2208.02814}.

\bibitem{vovk2005algorithmic}
Vovk V, Gammerman A, Shafer G.
\newblock Algorithmic Learning in a Random World.
\newblock Springer; 2005.
\newblock \url{https://doi.org/10.1007/b106715}.

\bibitem{geifman2017selective}
Geifman Y, El-Yaniv R. Selective Classification for Deep Neural Networks; 2017.
\newblock arXiv preprint arXiv:1705.08500.
\newblock \url{https://doi.org/10.48550/arXiv.1705.08500}.

\bibitem{geifman2019bias}
Geifman Y, Uziel G, El-Yaniv R. Bias-Reduced Uncertainty Estimation for Deep
  Neural Classifiers; 2019.
\newblock International Conference on Learning Representations (ICLR); arXiv
  preprint arXiv:1805.08206.
\newblock \url{https://doi.org/10.48550/arXiv.1805.08206}.

\bibitem{armato2011lidc}
Armato SG III, McLennan G, Bidaut L, McNitt-Gray MF, Meyer CR, Reeves AP,
  et~al.
\newblock The {Lung Image Database Consortium (LIDC)} and {Image Database
  Resource Initiative (IDRI)}: A completed reference database of lung nodules
  on {CT} scans.
\newblock Medical Physics. 2011;38(2):915-31.
\newblock \url{https://doi.org/10.1118/1.3528204}.

\bibitem{armato2015lidctcia}
Armato SG III, et~al. Data From {LIDC-IDRI}; 2015.
\newblock Licence: CC BY 3.0.
\newblock The Cancer Imaging Archive.
\newblock \url{https://doi.org/10.7937/K9/TCIA.2015.LO9QL9SX}.

\bibitem{clark2013tcia}
Clark K, Vendt B, Smith K, Freymann J, Kirby J, Koppel P, et~al.
\newblock The {Cancer Imaging Archive (TCIA)}: Maintaining and Operating a
  Public Information Repository.
\newblock Journal of Digital Imaging. 2013;26(6):1045-57.
\newblock \url{https://doi.org/10.1007/s10278-013-9622-7}.

\bibitem{setio2017luna16}
Setio AAA, Traverso A, de~Bel T, Berens MSN, van~den Bogaard C, Cerello P,
  et~al.
\newblock Validation, comparison, and combination of algorithms for automatic
  detection of pulmonary nodules in computed tomography images: The {LUNA16}
  challenge.
\newblock Medical Image Analysis. 2017;42:1-13.
\newblock \url{https://doi.org/10.1016/j.media.2017.06.015}.

\bibitem{nguyen2021ribcxr}
Nguyen HC, Le TT, Pham HH, Nguyen HQ.
\newblock {VinDr-RibCXR}: A Benchmark Dataset for Automatic Segmentation and
  Labeling of Individual Ribs on Chest {X}-rays.
\newblock In: Medical Imaging with Deep Learning ({MIDL}); 2021. Short paper;
  arXiv:2107.01327.

\bibitem{shiraishi2000jsrt}
Shiraishi J, Katsuragawa S, Ikezoe J, Matsumoto T, Kobayashi T, Komatsu Ki,
  et~al.
\newblock Development of a Digital Image Database for Chest Radiographs With
  and Without a Lung Nodule: Receiver Operating Characteristic Analysis of
  Radiologists' Detection of Pulmonary Nodules.
\newblock American Journal of Roentgenology. 2000;174(1):71-4.
\newblock \url{https://doi.org/10.2214/ajr.174.1.1740071}.

\bibitem{sogancioglu2024node21}
Sogancioglu E, van Ginneken B, Behrendt F, Bengs M, Schlaefer A, Radu M, et~al.
\newblock Nodule Detection and Generation on Chest {X}-Rays: {NODE21}
  Challenge.
\newblock {IEEE} Transactions on Medical Imaging. 2024;43(8):2839-53.
\newblock \url{https://doi.org/10.1109/TMI.2024.3382042}.

\bibitem{liu2020tbx11k}
Liu Y, Wu YH, Ban Y, Wang H, Cheng MM.
\newblock Rethinking Computer-Aided Tuberculosis Diagnosis.
\newblock In: Proceedings of the {IEEE/CVF} Conference on Computer Vision and
  Pattern Recognition ({CVPR}); 2020. p. 2643-52.
\newblock \url{https://doi.org/10.1109/CVPR42600.2020.00272}.

\bibitem{nguyen2022vindrcxr}
Nguyen HQ, Lam K, Le LT, Pham HH, Tran DQ, Nguyen DB, et~al.
\newblock {VinDr-CXR}: An open dataset of chest {X}-rays with radiologist's
  annotations.
\newblock Scientific Data. 2022;9:429.
\newblock \url{https://doi.org/10.1038/s41597-022-01498-w}.

\bibitem{vinbigdata2020kaggle}
Nguyen D, {DungNB}, Nguyen HQ, Elliott J, {NguyenThanhNhan}, Culliton P.
  {VinBigData} Chest {X}-ray Abnormalities Detection; 2020.
\newblock Accessed 2026-11-03.
\newblock Kaggle competition.
\newblock Available from:
  \url{https://www.kaggle.com/competitions/vinbigdata-chest-xray-abnormalities-detection}.

\bibitem{wang2017chestxray}
Wang X, Peng Y, Lu L, Lu Z, Bagheri M, Summers RM.
\newblock {ChestX-ray8}: Hospital-scale Chest {X}-ray Database and Benchmarks
  on Weakly-Supervised Classification and Localization of Common Thorax
  Diseases.
\newblock In: {IEEE} Conference on Computer Vision and Pattern Recognition
  ({CVPR}); 2017. p. 3462-71.
\newblock \url{https://doi.org/10.1109/CVPR.2017.369}.

\bibitem{durjoy2026bdcxr}
Durjoy SH, Shikder ME, Bijoy MHI, Shoib MMH, Rahman A, Hasan MZ, et~al.
  {BDCXR-3257}: A Comprehensive Chest Radiography Dataset with
  Multi-Etiological Pneumonia Labels from a Clinical Setting; 2026.
\newblock Version 1; licence CC BY 4.0.
\newblock figshare, dataset.
\newblock \url{https://doi.org/10.6084/m9.figshare.32086362.v1}.

\bibitem{kermany2018identifying}
Kermany DS, Goldbaum M, Cai W, Valentim CCS, Liang H, Baxter SL, et~al.
\newblock Identifying Medical Diagnoses and Treatable Diseases by Image-Based
  Deep Learning.
\newblock Cell. 2018;172(5):1122-31.e9.
\newblock \url{https://doi.org/10.1016/j.cell.2018.02.010}.

\bibitem{zhou2018unetpp}
Zhou Z, Rahman~Siddiquee MM, Tajbakhsh N, Liang J.
\newblock {UNet++}: A Nested {U-Net} Architecture for Medical Image
  Segmentation.
\newblock In: Deep Learning in Medical Image Analysis and Multimodal Learning
  for Clinical Decision Support ({DLMIA} 2018). vol. 11045 of Lecture Notes in
  Computer Science. Springer; 2018. p. 3-11.
\newblock Available from: \url{https://doi.org/10.1007/978-3-030-00889-5_1}.

\bibitem{woo2018cbam}
Woo S, Park J, Lee JY, Kweon IS.
\newblock {CBAM}: Convolutional Block Attention Module.
\newblock In: Computer Vision -- {ECCV} 2018. Lecture Notes in Computer
  Science. Springer; 2018. p. 3-19.
\newblock Available from: \url{https://doi.org/10.1007/978-3-030-01234-2_1}.

\bibitem{deng2009imagenet}
Deng J, Dong W, Socher R, Li LJ, Li K, Fei-Fei L.
\newblock {ImageNet}: A large-scale hierarchical image database.
\newblock In: {IEEE} Conference on Computer Vision and Pattern Recognition
  ({CVPR}); 2009. p. 248-55.
\newblock \url{https://doi.org/10.1109/CVPR.2009.5206848}.

\bibitem{wightman2019timm}
Wightman R. {PyTorch} Image Models. GitHub; 2019.
\newblock \url{https://github.com/rwightman/pytorch-image-models}.
\newblock \url{https://doi.org/10.5281/zenodo.4414861}.

\bibitem{loshchilov2019adamw}
Loshchilov I, Hutter F.
\newblock Decoupled Weight Decay Regularization.
\newblock In: International Conference on Learning Representations ({ICLR});
  2019. Available from: \url{https://arxiv.org/abs/1711.05101}.

\bibitem{simeoni2025dinov3}
Sim{\'e}oni O, Vo HV, Seitzer M, Baldassarre F, Oquab M, Jose C, et~al.
  {DINOv3}; 2025.
\newblock arXiv preprint arXiv:2508.10104.
\newblock \url{https://doi.org/10.48550/arXiv.2508.10104}.

\bibitem{perezgarcia2025raddino}
P{\'e}rez-Garc{\'\i}a F, Sharma H, Bond-Taylor S, Bouzid K, Salvatelli V, Ilse
  M, et~al.
\newblock Exploring scalable medical image encoders beyond text supervision.
\newblock Nature Machine Intelligence. 2025;7(1):119-30.
\newblock \url{https://doi.org/10.1038/s42256-024-00965-w}.

\bibitem{zuiderveld1994clahe}
Zuiderveld K.
\newblock Contrast Limited Adaptive Histogram Equalization.
\newblock In: Heckbert PS, editor. Graphics Gems {IV}. Boston: AP Professional;
  1994. p. 474-85.
\newblock \url{https://doi.org/10.1016/B978-0-12-336156-1.50061-6}.

\bibitem{oppenheim1968homomorphic}
Oppenheim AV, Schafer RW, Stockham TG.
\newblock Nonlinear filtering of multiplied and convolved signals.
\newblock Proceedings of the {IEEE}. 1968;56(8):1264-91.
\newblock \url{https://doi.org/10.1109/PROC.1968.6570}.

\bibitem{cohen2022torchxrayvision}
Cohen JP, Viviano JD, Bertin P, Morrison P, Torabian P, Guarrera M, et~al.
\newblock {TorchXRayVision}: A library of chest {X}-ray datasets and models.
\newblock In: Konukoglu E, Menze B, Venkataraman A, Baumgartner C, Dou Q,
  Albarqouni S, editors. Proceedings of The 5th International Conference on
  Medical Imaging with Deep Learning. vol. 172 of Proceedings of Machine
  Learning Research. PMLR; 2022. p. 231-49.

\bibitem{lian2021chestxdet}
Lian J, Liu J, Zhang S, Gao K, Liu X, Zhang D, et~al.
\newblock A Structure-Aware Relation Network for Thoracic Diseases Detection
  and Segmentation.
\newblock {IEEE} Transactions on Medical Imaging. 2021;40(8):2042-52.
\newblock \url{https://doi.org/10.1109/TMI.2021.3070847}.

\bibitem{wang2004ssim}
Wang Z, Bovik AC, Sheikh HR, Simoncelli EP.
\newblock Image quality assessment: from error visibility to structural
  similarity.
\newblock {IEEE} Transactions on Image Processing. 2004;13(4):600-12.
\newblock \url{https://doi.org/10.1109/TIP.2003.819861}.

\bibitem{wang2003msssim}
Wang Z, Simoncelli EP, Bovik AC.
\newblock Multiscale structural similarity for image quality assessment.
\newblock In: The Thirty-Seventh Asilomar Conference on Signals, Systems \&
  Computers. vol.~2; 2003. p. 1398-402.
\newblock \url{https://doi.org/10.1109/ACSSC.2003.1292216}.

\bibitem{zhang2018lpips}
Zhang R, Isola P, Efros AA, Shechtman E, Wang O.
\newblock The Unreasonable Effectiveness of Deep Features as a Perceptual
  Metric.
\newblock In: Proceedings of the {IEEE/CVF} Conference on Computer Vision and
  Pattern Recognition ({CVPR}); 2018. p. 586-95.
\newblock \url{https://doi.org/10.1109/CVPR.2018.00068}.

\bibitem{huang2017densenet}
Huang G, Liu Z, van~der Maaten L, Weinberger KQ.
\newblock Densely Connected Convolutional Networks.
\newblock In: {IEEE} Conference on Computer Vision and Pattern Recognition
  ({CVPR}); 2017. p. 2261-9.
\newblock \url{https://doi.org/10.1109/CVPR.2017.243}.

\bibitem{ren2017fasterrcnn}
Ren S, He K, Girshick R, Sun J.
\newblock Faster {R-CNN}: Towards Real-Time Object Detection with Region
  Proposal Networks.
\newblock {IEEE} Transactions on Pattern Analysis and Machine Intelligence.
  2017;39(6):1137-49.
\newblock \url{https://doi.org/10.1109/TPAMI.2016.2577031}.

\bibitem{vanginneken2010anode09}
van Ginneken B, Armato SG, de~Hoop B, van Amelsvoort-van~de Vorst S, Duindam T,
  Niemeijer M, et~al.
\newblock Comparing and combining algorithms for computer-aided detection of
  pulmonary nodules in computed tomography scans: The {ANODE09} study.
\newblock Medical Image Analysis. 2010;14(6):707-22.
\newblock \url{https://doi.org/10.1016/j.media.2010.05.005}.

\bibitem{niemeijer2011cpm}
Niemeijer M, Loog M, Abr{\`a}moff MD, Viergever MA, Prokop M, van Ginneken B.
\newblock On Combining Computer-Aided Detection Systems.
\newblock {IEEE} Transactions on Medical Imaging. 2011;30(2):215-23.
\newblock \url{https://doi.org/10.1109/TMI.2010.2072789}.

\bibitem{efron1993bootstrap}
Efron B, Tibshirani RJ.
\newblock An Introduction to the Bootstrap.
\newblock Boston, MA: Springer US; 1993.
\newblock \url{https://doi.org/10.1007/978-1-4899-4541-9}.

\bibitem{wilcoxon1945individual}
Wilcoxon F.
\newblock Individual Comparisons by Ranking Methods.
\newblock Biometrics Bulletin. 1945;1(6):80-3.
\newblock \url{https://doi.org/10.2307/3001968}.

\bibitem{delong1988auc}
DeLong ER, DeLong DM, Clarke-Pearson DL.
\newblock Comparing the Areas under Two or More Correlated Receiver Operating
  Characteristic Curves: A Nonparametric Approach.
\newblock Biometrics. 1988;44(3):837-45.
\newblock \url{https://doi.org/10.2307/2531595}.

\bibitem{holm1979simple}
Holm S.
\newblock A Simple Sequentially Rejective Multiple Test Procedure.
\newblock Scandinavian Journal of Statistics. 1979;6(2):65-70.
\newblock Available from: \url{https://www.jstor.org/stable/4615733}.

\bibitem{icru1989report44}
{International Commission on Radiation Units and Measurements}.
\newblock Tissue Substitutes in Radiation Dosimetry and Measurement.
\newblock International Commission on Radiation Units and Measurements; 1989.
  44.
\newblock Journal of the ICRU, vol. os23, no. 1.
\newblock \url{https://doi.org/10.1093/jicru/os23.1.Report44}.

\bibitem{hubbell2004xraymass}
Hubbell JH, Seltzer SM. Tables of {X}-Ray Mass Attenuation Coefficients and
  Mass Energy-Absorption Coefficients (version 1.4); 2004.
\newblock NIST Standard Reference Database 126; originally published as NISTIR
  5632 (1995). Accessed 2026-11-03.
\newblock National Institute of Standards and Technology, Gaithersburg, MD.
  \url{http://physics.nist.gov/xaamdi}.
\newblock \url{https://doi.org/10.18434/T4D01F}.

\bibitem{poludniowski2021spekpy}
Poludniowski GG, Omar A, Bujila R, Andreo P.
\newblock Technical Note: {SpekPy} v2.0---a software toolkit for modeling x-ray
  tube spectra.
\newblock Medical Physics. 2021;48(7):3630-7.
\newblock \url{https://doi.org/10.1002/mp.14945}.

\bibitem{tejani2024claim}
Tejani AS, Klontzas ME, Gatti AA, Mongan JT, Moy L, Park SH, et~al.
\newblock Checklist for Artificial Intelligence in Medical Imaging ({CLAIM}):
  2024 Update.
\newblock Radiology: Artificial Intelligence. 2024;6(4):e240300.
\newblock \url{https://doi.org/10.1148/ryai.240300}.

\end{thebibliography}
 by
% \bibliographystyle{vancouver-doi} \bibliography{refs}, compile, and run tools/freeze_refs.py on main.bbl.
\begin{thebibliography}{10}

\bibitem{kuhlman2006des}
Kuhlman JE, Collins J, Brooks GN, Yandow DR, Broderick LS.
\newblock Dual-Energy Subtraction Chest Radiography: What to Look for beyond
  Calcified Nodules.
\newblock RadioGraphics. 2006;26(1):79-92.
\newblock \url{https://doi.org/10.1148/rg.261055034}.

\bibitem{sun2026bsldm}
Sun Y, Chen Z, Zheng H, Deng W, Liu J, Min W, et~al.
\newblock {BS-LDM}: Effective Bone Suppression in High-Resolution Chest {X}-Ray
  Images With Conditional Latent Diffusion Models.
\newblock {IEEE} Journal of Biomedical and Health Informatics.
  2026;30(1):376-87.
\newblock Online 14 July 2025.
\newblock \url{https://doi.org/10.1109/JBHI.2025.3588138}.

\bibitem{freedman2011lung}
Freedman MT, Lo SCB, Seibel JC, Bromley CM.
\newblock Lung Nodules: Improved Detection with Software That Suppresses the
  Rib and Clavicle on Chest Radiographs.
\newblock Radiology. 2011;260(1):265-73.
\newblock \url{https://doi.org/10.1148/radiol.11100153}.

\bibitem{schalekamp2014cad}
Schalekamp S, van Ginneken B, Koedam E, Snoeren MM, Tiehuis AM, Wittenberg R,
  et~al.
\newblock Computer-aided Detection Improves Detection of Pulmonary Nodules in
  Chest Radiographs beyond the Support by Bone-suppressed Images.
\newblock Radiology. 2014;272(1):252-61.
\newblock \url{https://doi.org/10.1148/radiol.14131315}.

\bibitem{rajaraman2021bone}
Rajaraman S, Zamzmi G, Folio L, Alderson P, Antani S.
\newblock Chest {X}-ray Bone Suppression for Improving Classification of
  Tuberculosis-Consistent Findings.
\newblock Diagnostics. 2021;11(5):840.
\newblock \url{https://doi.org/10.3390/diagnostics11050840}.

\bibitem{chen2024bsdiff}
Chen Z, Sun Y, Qin W, Ge R, Pan C, Deng W, et~al. {BS-Diff}: Effective Bone
  Suppression Using Conditional Diffusion Models from Chest {X}-Ray Images;
  2024.
\newblock Accepted at IEEE ISBI 2024.
\newblock arXiv preprint arXiv:2311.15328.
\newblock \url{https://doi.org/10.48550/arXiv.2311.15328}.

\bibitem{sun2025gllcm}
Sun Y, Chen Z, Zheng H, Lu Y, Duan L, Fan F, et~al. {GL-LCM}: Global-Local
  Latent Consistency Models for Fast High-Resolution Bone Suppression in Chest
  {X}-Ray Images; 2025.
\newblock Accepted at MICCAI 2025.
\newblock arXiv preprint arXiv:2508.03357.
\newblock \url{https://doi.org/10.48550/arXiv.2508.03357}.

\bibitem{sun2026debonedit}
Sun Y, Fan F, Jia J, Deng W, Xu H, Wang C, et~al. {DeBoneDiT}: Depth-Driven
  Conditional Bridge Diffusion Transformers for Bone Suppression in Chest
  {X}-Ray Images; 2026.
\newblock arXiv preprint arXiv:2610.19466.
\newblock \url{https://doi.org/10.48550/arXiv.2610.19466}.

\bibitem{angaitkar2026anatomy}
Angaitkar M, Kumar P, Satia A, Rao P, Mittal A, Tadepalli M, et~al.
  Anatomy-Decomposed Chest Computed Tomography ({CT}) Projections as Scalable
  Supervision for Bone Suppression in Chest Radiographs; 2026.
\newblock Preprint.
\newblock arXiv preprint arXiv:2609.24937.
\newblock \url{https://doi.org/10.48550/arXiv.2609.24937}.

\bibitem{hamamci2026ctrate}
Hamamci IE, Er S, Wang C, Almas F, Simsek AG, Esirgun SN, et~al.
\newblock Generalist foundation models from a multimodal dataset for {3D}
  computed tomography.
\newblock Nature Biomedical Engineering. 2026;10(8):1610-28.
\newblock Preprint: arXiv:2403.17834 (2024).
\newblock \url{https://doi.org/10.1038/s41551-025-01599-y}.

\bibitem{hogeweg2013suppression}
Hogeweg L, Sanchez CI, van Ginneken B.
\newblock Suppression of Translucent Elongated Structures: Applications in
  Chest Radiography.
\newblock {IEEE} Transactions on Medical Imaging. 2013;32(11):2099-113.
\newblock \url{https://doi.org/10.1109/TMI.2013.2274212}.

\bibitem{takagi2020dose}
Takagi S, Yaegashi T, Ishikawa M.
\newblock Dose reduction and image quality improvement of chest radiography by
  using bone-suppression technique and low tube voltage: a phantom study.
\newblock European Radiology. 2020;30(1):571-80.
\newblock \url{https://doi.org/10.1007/s00330-019-06375-6}.

\bibitem{cohen2018distribution}
Cohen JP, Luck M, Honari S.
\newblock Distribution Matching Losses Can Hallucinate Features in Medical
  Image Translation.
\newblock In: Medical Image Computing and Computer Assisted Intervention --
  {MICCAI} 2018. Lecture Notes in Computer Science. Springer; 2018. p. 529-36.
\newblock Available from: \url{https://doi.org/10.1007/978-3-030-00928-1_60}.

\bibitem{antun2020instabilities}
Antun V, Renna F, Poon C, Adcock B, Hansen AC.
\newblock On instabilities of deep learning in image reconstruction and the
  potential costs of {AI}.
\newblock Proceedings of the National Academy of Sciences.
  2020;117(48):30088-95.
\newblock \url{https://doi.org/10.1073/pnas.1907377117}.

\bibitem{bhadra2021hallucinations}
Bhadra S, Kelkar VA, Brooks FJ, Anastasio MA.
\newblock On Hallucinations in Tomographic Image Reconstruction.
\newblock {IEEE} Transactions on Medical Imaging. 2021;40(11):3249-60.
\newblock \url{https://doi.org/10.1109/TMI.2021.3077857}.

\bibitem{tivnan2024hallucination}
Tivnan M, Yoon S, Chen Z, Li X, Wu D, Li Q.
\newblock Hallucination Index: An Image Quality Metric for Generative
  Reconstruction Models.
\newblock In: Medical Image Computing and Computer Assisted Intervention --
  {MICCAI} 2024. Lecture Notes in Computer Science. Springer; 2024. p. 449-58.
\newblock Available from: \url{https://doi.org/10.1007/978-3-031-72117-5_42}.

\bibitem{varoquaux2022failures}
Varoquaux G, Cheplygina V.
\newblock Machine learning for medical imaging: methodological failures and
  recommendations for the future.
\newblock npj Digital Medicine. 2022;5:48.
\newblock \url{https://doi.org/10.1038/s41746-022-00592-y}.

\bibitem{simmons2011falsepositive}
Simmons JP, Nelson LD, Simonsohn U.
\newblock False-Positive Psychology: Undisclosed Flexibility in Data Collection
  and Analysis Allows Presenting Anything as Significant.
\newblock Psychological Science. 2011;22(11):1359-66.
\newblock \url{https://doi.org/10.1177/0956797611417632}.

\bibitem{suzuki2006mtann}
Suzuki K, Abe H, MacMahon H, Doi K.
\newblock Image-processing technique for suppressing ribs in chest radiographs
  by means of massive training artificial neural network ({MTANN}).
\newblock {IEEE} Transactions on Medical Imaging. 2006;25(4):406-16.
\newblock \url{https://doi.org/10.1109/TMI.2006.871549}.

\bibitem{yang2017cascade}
Yang W, Chen Y, Liu Y, Zhong L, Qin G, Lu Z, et~al.
\newblock Cascade of multi-scale convolutional neural networks for bone
  suppression of chest radiographs in gradient domain.
\newblock Medical Image Analysis. 2017;35:421-33.
\newblock \url{https://doi.org/10.1016/j.media.2016.08.004}.

\bibitem{woo2023convnextv2}
Woo S, Debnath S, Hu R, Chen X, Liu Z, Kweon IS, et~al.
\newblock {ConvNeXt V2}: Co-designing and Scaling {ConvNets} with Masked
  Autoencoders.
\newblock In: Proceedings of the {IEEE/CVF} Conference on Computer Vision and
  Pattern Recognition ({CVPR}); 2023. p. 16133-42.
\newblock \url{https://doi.org/10.1109/CVPR52729.2023.01548}.

\bibitem{liu2022swinv2}
Liu Z, Hu H, Lin Y, Yao Z, Xie Z, Wei Y, et~al.
\newblock Swin Transformer {V2}: Scaling Up Capacity and Resolution.
\newblock In: Proceedings of the {IEEE/CVF} Conference on Computer Vision and
  Pattern Recognition ({CVPR}); 2022. p. 11999-2009.
\newblock \url{https://doi.org/10.1109/CVPR52688.2022.01170}.

\bibitem{cheng2020addressing}
Cheng K, Caliv{\'a} F, Shah R, Han M, Majumdar S, Pedoia V.
\newblock Addressing The False Negative Problem of Deep Learning {MRI}
  Reconstruction Models by Adversarial Attacks and Robust Training.
\newblock In: Proceedings of the Third Conference on Medical Imaging with Deep
  Learning. vol. 121 of Proceedings of Machine Learning Research. PMLR; 2020.
  p. 121-35.
\newblock Available from:
  \url{https://proceedings.mlr.press/v121/cheng20a.html}.

\bibitem{li2021assessing}
Li K, Zhou W, Li H, Anastasio MA.
\newblock Assessing the Impact of Deep Neural Network-Based Image Denoising on
  Binary Signal Detection Tasks.
\newblock {IEEE} Transactions on Medical Imaging. 2021;40(9):2295-305.
\newblock \url{https://doi.org/10.1109/TMI.2021.3076810}.

\bibitem{liu2021query2label}
Liu S, Zhang L, Yang X, Su H, Zhu J. {Query2Label}: A Simple Transformer Way to
  Multi-Label Classification; 2021.
\newblock arXiv preprint arXiv:2107.10834.
\newblock \url{https://doi.org/10.48550/arXiv.2107.10834}.

\bibitem{ridnik2023mldecoder}
Ridnik T, Sharir G, Ben-Cohen A, Ben-Baruch E, Noy A.
\newblock {ML-Decoder}: Scalable and Versatile Classification Head.
\newblock In: 2023 IEEE/CVF Winter Conference on Applications of Computer
  Vision (WACV); 2023. p. 32-41.
\newblock \url{https://doi.org/10.1109/WACV56688.2023.00012}.

\bibitem{chattopadhay2018gradcampp}
Chattopadhay A, Sarkar A, Howlader P, Balasubramanian VN.
\newblock {Grad-CAM++}: Generalized Gradient-Based Visual Explanations for Deep
  Convolutional Networks.
\newblock In: 2018 IEEE Winter Conference on Applications of Computer Vision
  (WACV); 2018. p. 839-47.
\newblock \url{https://doi.org/10.1109/WACV.2018.00097}.

\bibitem{sundararajan2017axiomatic}
Sundararajan M, Taly A, Yan Q. Axiomatic Attribution for Deep Networks; 2017.
\newblock arXiv preprint arXiv:1703.01365.
\newblock \url{https://doi.org/10.48550/arXiv.1703.01365}.

\bibitem{zhang2018topdown}
Zhang J, Bargal SA, Lin Z, Brandt J, Shen X, Sclaroff S.
\newblock Top-Down Neural Attention by Excitation Backprop.
\newblock International Journal of Computer Vision. 2018;126(10):1084-102.
\newblock \url{https://doi.org/10.1007/s11263-017-1059-x}.

\bibitem{lakshminarayanan2017ensembles}
Lakshminarayanan B, Pritzel A, Blundell C.
\newblock Simple and Scalable Predictive Uncertainty Estimation using Deep
  Ensembles.
\newblock In: Advances in Neural Information Processing Systems. vol.~30; 2017.
  Available from:
  \url{https://proceedings.neurips.cc/paper_files/paper/2017/hash/9ef2ed4b7fd2c810847ffa5fa85bce38-Abstract.html}.

\bibitem{lei2018conformal}
Lei J, G'Sell M, Rinaldo A, Tibshirani RJ, Wasserman L.
\newblock Distribution-Free Predictive Inference for Regression.
\newblock Journal of the American Statistical Association.
  2018;113(523):1094-111.
\newblock \url{https://doi.org/10.1080/01621459.2017.1307116}.

\bibitem{angelopoulos2023conformal}
Angelopoulos AN, Bates S.
\newblock Conformal Prediction: A Gentle Introduction.
\newblock Foundations and Trends in Machine Learning. 2023;16(4):494-591.
\newblock Preprint title: A Gentle Introduction to Conformal Prediction and
  Distribution-Free Uncertainty Quantification.
\newblock \url{https://doi.org/10.1561/2200000101}.

\bibitem{angelopoulos2022crc}
Angelopoulos AN, Bates S, Fisch A, Lei L, Schuster T. Conformal Risk Control;
  2022.
\newblock arXiv preprint arXiv:2208.02814.
\newblock \url{https://doi.org/10.48550/arXiv.2208.02814}.

\bibitem{vovk2005algorithmic}
Vovk V, Gammerman A, Shafer G.
\newblock Algorithmic Learning in a Random World.
\newblock Springer; 2005.
\newblock \url{https://doi.org/10.1007/b106715}.

\bibitem{geifman2017selective}
Geifman Y, El-Yaniv R. Selective Classification for Deep Neural Networks; 2017.
\newblock arXiv preprint arXiv:1705.08500.
\newblock \url{https://doi.org/10.48550/arXiv.1705.08500}.

\bibitem{geifman2019bias}
Geifman Y, Uziel G, El-Yaniv R. Bias-Reduced Uncertainty Estimation for Deep
  Neural Classifiers; 2019.
\newblock International Conference on Learning Representations (ICLR); arXiv
  preprint arXiv:1805.08206.
\newblock \url{https://doi.org/10.48550/arXiv.1805.08206}.

\bibitem{armato2011lidc}
Armato SG III, McLennan G, Bidaut L, McNitt-Gray MF, Meyer CR, Reeves AP,
  et~al.
\newblock The {Lung Image Database Consortium (LIDC)} and {Image Database
  Resource Initiative (IDRI)}: A completed reference database of lung nodules
  on {CT} scans.
\newblock Medical Physics. 2011;38(2):915-31.
\newblock \url{https://doi.org/10.1118/1.3528204}.

\bibitem{armato2015lidctcia}
Armato SG III, et~al. Data From {LIDC-IDRI}; 2015.
\newblock Licence: CC BY 3.0.
\newblock The Cancer Imaging Archive.
\newblock \url{https://doi.org/10.7937/K9/TCIA.2015.LO9QL9SX}.

\bibitem{clark2013tcia}
Clark K, Vendt B, Smith K, Freymann J, Kirby J, Koppel P, et~al.
\newblock The {Cancer Imaging Archive (TCIA)}: Maintaining and Operating a
  Public Information Repository.
\newblock Journal of Digital Imaging. 2013;26(6):1045-57.
\newblock \url{https://doi.org/10.1007/s10278-013-9622-7}.

\bibitem{setio2017luna16}
Setio AAA, Traverso A, de~Bel T, Berens MSN, van~den Bogaard C, Cerello P,
  et~al.
\newblock Validation, comparison, and combination of algorithms for automatic
  detection of pulmonary nodules in computed tomography images: The {LUNA16}
  challenge.
\newblock Medical Image Analysis. 2017;42:1-13.
\newblock \url{https://doi.org/10.1016/j.media.2017.06.015}.

\bibitem{nguyen2021ribcxr}
Nguyen HC, Le TT, Pham HH, Nguyen HQ.
\newblock {VinDr-RibCXR}: A Benchmark Dataset for Automatic Segmentation and
  Labeling of Individual Ribs on Chest {X}-rays.
\newblock In: Medical Imaging with Deep Learning ({MIDL}); 2021. Short paper;
  arXiv:2107.01327.

\bibitem{shiraishi2000jsrt}
Shiraishi J, Katsuragawa S, Ikezoe J, Matsumoto T, Kobayashi T, Komatsu Ki,
  et~al.
\newblock Development of a Digital Image Database for Chest Radiographs With
  and Without a Lung Nodule: Receiver Operating Characteristic Analysis of
  Radiologists' Detection of Pulmonary Nodules.
\newblock American Journal of Roentgenology. 2000;174(1):71-4.
\newblock \url{https://doi.org/10.2214/ajr.174.1.1740071}.

\bibitem{sogancioglu2024node21}
Sogancioglu E, van Ginneken B, Behrendt F, Bengs M, Schlaefer A, Radu M, et~al.
\newblock Nodule Detection and Generation on Chest {X}-Rays: {NODE21}
  Challenge.
\newblock {IEEE} Transactions on Medical Imaging. 2024;43(8):2839-53.
\newblock \url{https://doi.org/10.1109/TMI.2024.3382042}.

\bibitem{liu2020tbx11k}
Liu Y, Wu YH, Ban Y, Wang H, Cheng MM.
\newblock Rethinking Computer-Aided Tuberculosis Diagnosis.
\newblock In: Proceedings of the {IEEE/CVF} Conference on Computer Vision and
  Pattern Recognition ({CVPR}); 2020. p. 2643-52.
\newblock \url{https://doi.org/10.1109/CVPR42600.2020.00272}.

\bibitem{nguyen2022vindrcxr}
Nguyen HQ, Lam K, Le LT, Pham HH, Tran DQ, Nguyen DB, et~al.
\newblock {VinDr-CXR}: An open dataset of chest {X}-rays with radiologist's
  annotations.
\newblock Scientific Data. 2022;9:429.
\newblock \url{https://doi.org/10.1038/s41597-022-01498-w}.

\bibitem{vinbigdata2020kaggle}
Nguyen D, {DungNB}, Nguyen HQ, Elliott J, {NguyenThanhNhan}, Culliton P.
  {VinBigData} Chest {X}-ray Abnormalities Detection; 2020.
\newblock Accessed 2026-11-03.
\newblock Kaggle competition.
\newblock Available from:
  \url{https://www.kaggle.com/competitions/vinbigdata-chest-xray-abnormalities-detection}.

\bibitem{wang2017chestxray}
Wang X, Peng Y, Lu L, Lu Z, Bagheri M, Summers RM.
\newblock {ChestX-ray8}: Hospital-scale Chest {X}-ray Database and Benchmarks
  on Weakly-Supervised Classification and Localization of Common Thorax
  Diseases.
\newblock In: {IEEE} Conference on Computer Vision and Pattern Recognition
  ({CVPR}); 2017. p. 3462-71.
\newblock \url{https://doi.org/10.1109/CVPR.2017.369}.

\bibitem{durjoy2026bdcxr}
Durjoy SH, Shikder ME, Bijoy MHI, Shoib MMH, Rahman A, Hasan MZ, et~al.
  {BDCXR-3257}: A Comprehensive Chest Radiography Dataset with
  Multi-Etiological Pneumonia Labels from a Clinical Setting; 2026.
\newblock Version 1; licence CC BY 4.0.
\newblock figshare, dataset.
\newblock \url{https://doi.org/10.6084/m9.figshare.32086362.v1}.

\bibitem{kermany2018identifying}
Kermany DS, Goldbaum M, Cai W, Valentim CCS, Liang H, Baxter SL, et~al.
\newblock Identifying Medical Diagnoses and Treatable Diseases by Image-Based
  Deep Learning.
\newblock Cell. 2018;172(5):1122-31.e9.
\newblock \url{https://doi.org/10.1016/j.cell.2018.02.010}.

\bibitem{zhou2018unetpp}
Zhou Z, Rahman~Siddiquee MM, Tajbakhsh N, Liang J.
\newblock {UNet++}: A Nested {U-Net} Architecture for Medical Image
  Segmentation.
\newblock In: Deep Learning in Medical Image Analysis and Multimodal Learning
  for Clinical Decision Support ({DLMIA} 2018). vol. 11045 of Lecture Notes in
  Computer Science. Springer; 2018. p. 3-11.
\newblock Available from: \url{https://doi.org/10.1007/978-3-030-00889-5_1}.

\bibitem{woo2018cbam}
Woo S, Park J, Lee JY, Kweon IS.
\newblock {CBAM}: Convolutional Block Attention Module.
\newblock In: Computer Vision -- {ECCV} 2018. Lecture Notes in Computer
  Science. Springer; 2018. p. 3-19.
\newblock Available from: \url{https://doi.org/10.1007/978-3-030-01234-2_1}.

\bibitem{deng2009imagenet}
Deng J, Dong W, Socher R, Li LJ, Li K, Fei-Fei L.
\newblock {ImageNet}: A large-scale hierarchical image database.
\newblock In: {IEEE} Conference on Computer Vision and Pattern Recognition
  ({CVPR}); 2009. p. 248-55.
\newblock \url{https://doi.org/10.1109/CVPR.2009.5206848}.

\bibitem{wightman2019timm}
Wightman R. {PyTorch} Image Models. GitHub; 2019.
\newblock \url{https://github.com/rwightman/pytorch-image-models}.
\newblock \url{https://doi.org/10.5281/zenodo.4414861}.

\bibitem{loshchilov2019adamw}
Loshchilov I, Hutter F.
\newblock Decoupled Weight Decay Regularization.
\newblock In: International Conference on Learning Representations ({ICLR});
  2019. Available from: \url{https://arxiv.org/abs/1711.05101}.

\bibitem{simeoni2025dinov3}
Sim{\'e}oni O, Vo HV, Seitzer M, Baldassarre F, Oquab M, Jose C, et~al.
  {DINOv3}; 2025.
\newblock arXiv preprint arXiv:2508.10104.
\newblock \url{https://doi.org/10.48550/arXiv.2508.10104}.

\bibitem{perezgarcia2025raddino}
P{\'e}rez-Garc{\'\i}a F, Sharma H, Bond-Taylor S, Bouzid K, Salvatelli V, Ilse
  M, et~al.
\newblock Exploring scalable medical image encoders beyond text supervision.
\newblock Nature Machine Intelligence. 2025;7(1):119-30.
\newblock \url{https://doi.org/10.1038/s42256-024-00965-w}.

\bibitem{zuiderveld1994clahe}
Zuiderveld K.
\newblock Contrast Limited Adaptive Histogram Equalization.
\newblock In: Heckbert PS, editor. Graphics Gems {IV}. Boston: AP Professional;
  1994. p. 474-85.
\newblock \url{https://doi.org/10.1016/B978-0-12-336156-1.50061-6}.

\bibitem{oppenheim1968homomorphic}
Oppenheim AV, Schafer RW, Stockham TG.
\newblock Nonlinear filtering of multiplied and convolved signals.
\newblock Proceedings of the {IEEE}. 1968;56(8):1264-91.
\newblock \url{https://doi.org/10.1109/PROC.1968.6570}.

\bibitem{cohen2022torchxrayvision}
Cohen JP, Viviano JD, Bertin P, Morrison P, Torabian P, Guarrera M, et~al.
\newblock {TorchXRayVision}: A library of chest {X}-ray datasets and models.
\newblock In: Konukoglu E, Menze B, Venkataraman A, Baumgartner C, Dou Q,
  Albarqouni S, editors. Proceedings of The 5th International Conference on
  Medical Imaging with Deep Learning. vol. 172 of Proceedings of Machine
  Learning Research. PMLR; 2022. p. 231-49.

\bibitem{lian2021chestxdet}
Lian J, Liu J, Zhang S, Gao K, Liu X, Zhang D, et~al.
\newblock A Structure-Aware Relation Network for Thoracic Diseases Detection
  and Segmentation.
\newblock {IEEE} Transactions on Medical Imaging. 2021;40(8):2042-52.
\newblock \url{https://doi.org/10.1109/TMI.2021.3070847}.

\bibitem{wang2004ssim}
Wang Z, Bovik AC, Sheikh HR, Simoncelli EP.
\newblock Image quality assessment: from error visibility to structural
  similarity.
\newblock {IEEE} Transactions on Image Processing. 2004;13(4):600-12.
\newblock \url{https://doi.org/10.1109/TIP.2003.819861}.

\bibitem{wang2003msssim}
Wang Z, Simoncelli EP, Bovik AC.
\newblock Multiscale structural similarity for image quality assessment.
\newblock In: The Thirty-Seventh Asilomar Conference on Signals, Systems \&
  Computers. vol.~2; 2003. p. 1398-402.
\newblock \url{https://doi.org/10.1109/ACSSC.2003.1292216}.

\bibitem{zhang2018lpips}
Zhang R, Isola P, Efros AA, Shechtman E, Wang O.
\newblock The Unreasonable Effectiveness of Deep Features as a Perceptual
  Metric.
\newblock In: Proceedings of the {IEEE/CVF} Conference on Computer Vision and
  Pattern Recognition ({CVPR}); 2018. p. 586-95.
\newblock \url{https://doi.org/10.1109/CVPR.2018.00068}.

\bibitem{huang2017densenet}
Huang G, Liu Z, van~der Maaten L, Weinberger KQ.
\newblock Densely Connected Convolutional Networks.
\newblock In: {IEEE} Conference on Computer Vision and Pattern Recognition
  ({CVPR}); 2017. p. 2261-9.
\newblock \url{https://doi.org/10.1109/CVPR.2017.243}.

\bibitem{ren2017fasterrcnn}
Ren S, He K, Girshick R, Sun J.
\newblock Faster {R-CNN}: Towards Real-Time Object Detection with Region
  Proposal Networks.
\newblock {IEEE} Transactions on Pattern Analysis and Machine Intelligence.
  2017;39(6):1137-49.
\newblock \url{https://doi.org/10.1109/TPAMI.2016.2577031}.

\bibitem{vanginneken2010anode09}
van Ginneken B, Armato SG, de~Hoop B, van Amelsvoort-van~de Vorst S, Duindam T,
  Niemeijer M, et~al.
\newblock Comparing and combining algorithms for computer-aided detection of
  pulmonary nodules in computed tomography scans: The {ANODE09} study.
\newblock Medical Image Analysis. 2010;14(6):707-22.
\newblock \url{https://doi.org/10.1016/j.media.2010.05.005}.

\bibitem{niemeijer2011cpm}
Niemeijer M, Loog M, Abr{\`a}moff MD, Viergever MA, Prokop M, van Ginneken B.
\newblock On Combining Computer-Aided Detection Systems.
\newblock {IEEE} Transactions on Medical Imaging. 2011;30(2):215-23.
\newblock \url{https://doi.org/10.1109/TMI.2010.2072789}.

\bibitem{efron1993bootstrap}
Efron B, Tibshirani RJ.
\newblock An Introduction to the Bootstrap.
\newblock Boston, MA: Springer US; 1993.
\newblock \url{https://doi.org/10.1007/978-1-4899-4541-9}.

\bibitem{wilcoxon1945individual}
Wilcoxon F.
\newblock Individual Comparisons by Ranking Methods.
\newblock Biometrics Bulletin. 1945;1(6):80-3.
\newblock \url{https://doi.org/10.2307/3001968}.

\bibitem{delong1988auc}
DeLong ER, DeLong DM, Clarke-Pearson DL.
\newblock Comparing the Areas under Two or More Correlated Receiver Operating
  Characteristic Curves: A Nonparametric Approach.
\newblock Biometrics. 1988;44(3):837-45.
\newblock \url{https://doi.org/10.2307/2531595}.

\bibitem{holm1979simple}
Holm S.
\newblock A Simple Sequentially Rejective Multiple Test Procedure.
\newblock Scandinavian Journal of Statistics. 1979;6(2):65-70.
\newblock Available from: \url{https://www.jstor.org/stable/4615733}.

\bibitem{icru1989report44}
{International Commission on Radiation Units and Measurements}.
\newblock Tissue Substitutes in Radiation Dosimetry and Measurement.
\newblock International Commission on Radiation Units and Measurements; 1989.
  44.
\newblock Journal of the ICRU, vol. os23, no. 1.
\newblock \url{https://doi.org/10.1093/jicru/os23.1.Report44}.

\bibitem{hubbell2004xraymass}
Hubbell JH, Seltzer SM. Tables of {X}-Ray Mass Attenuation Coefficients and
  Mass Energy-Absorption Coefficients (version 1.4); 2004.
\newblock NIST Standard Reference Database 126; originally published as NISTIR
  5632 (1995). Accessed 2026-11-03.
\newblock National Institute of Standards and Technology, Gaithersburg, MD.
  \url{http://physics.nist.gov/xaamdi}.
\newblock \url{https://doi.org/10.18434/T4D01F}.

\bibitem{poludniowski2021spekpy}
Poludniowski GG, Omar A, Bujila R, Andreo P.
\newblock Technical Note: {SpekPy} v2.0---a software toolkit for modeling x-ray
  tube spectra.
\newblock Medical Physics. 2021;48(7):3630-7.
\newblock \url{https://doi.org/10.1002/mp.14945}.

\bibitem{tejani2024claim}
Tejani AS, Klontzas ME, Gatti AA, Mongan JT, Moy L, Park SH, et~al.
\newblock Checklist for Artificial Intelligence in Medical Imaging ({CLAIM}):
  2024 Update.
\newblock Radiology: Artificial Intelligence. 2024;6(4):e240300.
\newblock \url{https://doi.org/10.1148/ryai.240300}.

\end{thebibliography}
 by
% \bibliographystyle{vancouver-doi} \bibliography{refs}, compile, and run tools/freeze_refs.py on main.bbl.
\begin{thebibliography}{10}

\bibitem{kuhlman2006des}
Kuhlman JE, Collins J, Brooks GN, Yandow DR, Broderick LS.
\newblock Dual-Energy Subtraction Chest Radiography: What to Look for beyond
  Calcified Nodules.
\newblock RadioGraphics. 2006;26(1):79-92.
\newblock \url{https://doi.org/10.1148/rg.261055034}.

\bibitem{sun2026bsldm}
Sun Y, Chen Z, Zheng H, Deng W, Liu J, Min W, et~al.
\newblock {BS-LDM}: Effective Bone Suppression in High-Resolution Chest {X}-Ray
  Images With Conditional Latent Diffusion Models.
\newblock {IEEE} Journal of Biomedical and Health Informatics.
  2026;30(1):376-87.
\newblock Online 14 July 2025.
\newblock \url{https://doi.org/10.1109/JBHI.2025.3588138}.

\bibitem{freedman2011lung}
Freedman MT, Lo SCB, Seibel JC, Bromley CM.
\newblock Lung Nodules: Improved Detection with Software That Suppresses the
  Rib and Clavicle on Chest Radiographs.
\newblock Radiology. 2011;260(1):265-73.
\newblock \url{https://doi.org/10.1148/radiol.11100153}.

\bibitem{schalekamp2014cad}
Schalekamp S, van Ginneken B, Koedam E, Snoeren MM, Tiehuis AM, Wittenberg R,
  et~al.
\newblock Computer-aided Detection Improves Detection of Pulmonary Nodules in
  Chest Radiographs beyond the Support by Bone-suppressed Images.
\newblock Radiology. 2014;272(1):252-61.
\newblock \url{https://doi.org/10.1148/radiol.14131315}.

\bibitem{rajaraman2021bone}
Rajaraman S, Zamzmi G, Folio L, Alderson P, Antani S.
\newblock Chest {X}-ray Bone Suppression for Improving Classification of
  Tuberculosis-Consistent Findings.
\newblock Diagnostics. 2021;11(5):840.
\newblock \url{https://doi.org/10.3390/diagnostics11050840}.

\bibitem{chen2024bsdiff}
Chen Z, Sun Y, Qin W, Ge R, Pan C, Deng W, et~al. {BS-Diff}: Effective Bone
  Suppression Using Conditional Diffusion Models from Chest {X}-Ray Images;
  2024.
\newblock Accepted at IEEE ISBI 2024.
\newblock arXiv preprint arXiv:2311.15328.
\newblock \url{https://doi.org/10.48550/arXiv.2311.15328}.

\bibitem{sun2025gllcm}
Sun Y, Chen Z, Zheng H, Lu Y, Duan L, Fan F, et~al. {GL-LCM}: Global-Local
  Latent Consistency Models for Fast High-Resolution Bone Suppression in Chest
  {X}-Ray Images; 2025.
\newblock Accepted at MICCAI 2025.
\newblock arXiv preprint arXiv:2508.03357.
\newblock \url{https://doi.org/10.48550/arXiv.2508.03357}.

\bibitem{sun2026debonedit}
Sun Y, Fan F, Jia J, Deng W, Xu H, Wang C, et~al. {DeBoneDiT}: Depth-Driven
  Conditional Bridge Diffusion Transformers for Bone Suppression in Chest
  {X}-Ray Images; 2026.
\newblock arXiv preprint arXiv:2610.19466.
\newblock \url{https://doi.org/10.48550/arXiv.2610.19466}.

\bibitem{angaitkar2026anatomy}
Angaitkar M, Kumar P, Satia A, Rao P, Mittal A, Tadepalli M, et~al.
  Anatomy-Decomposed Chest Computed Tomography ({CT}) Projections as Scalable
  Supervision for Bone Suppression in Chest Radiographs; 2026.
\newblock Preprint.
\newblock arXiv preprint arXiv:2609.24937.
\newblock \url{https://doi.org/10.48550/arXiv.2609.24937}.

\bibitem{hamamci2026ctrate}
Hamamci IE, Er S, Wang C, Almas F, Simsek AG, Esirgun SN, et~al.
\newblock Generalist foundation models from a multimodal dataset for {3D}
  computed tomography.
\newblock Nature Biomedical Engineering. 2026;10(8):1610-28.
\newblock Preprint: arXiv:2403.17834 (2024).
\newblock \url{https://doi.org/10.1038/s41551-025-01599-y}.

\bibitem{hogeweg2013suppression}
Hogeweg L, Sanchez CI, van Ginneken B.
\newblock Suppression of Translucent Elongated Structures: Applications in
  Chest Radiography.
\newblock {IEEE} Transactions on Medical Imaging. 2013;32(11):2099-113.
\newblock \url{https://doi.org/10.1109/TMI.2013.2274212}.

\bibitem{takagi2020dose}
Takagi S, Yaegashi T, Ishikawa M.
\newblock Dose reduction and image quality improvement of chest radiography by
  using bone-suppression technique and low tube voltage: a phantom study.
\newblock European Radiology. 2020;30(1):571-80.
\newblock \url{https://doi.org/10.1007/s00330-019-06375-6}.

\bibitem{cohen2018distribution}
Cohen JP, Luck M, Honari S.
\newblock Distribution Matching Losses Can Hallucinate Features in Medical
  Image Translation.
\newblock In: Medical Image Computing and Computer Assisted Intervention --
  {MICCAI} 2018. Lecture Notes in Computer Science. Springer; 2018. p. 529-36.
\newblock Available from: \url{https://doi.org/10.1007/978-3-030-00928-1_60}.

\bibitem{antun2020instabilities}
Antun V, Renna F, Poon C, Adcock B, Hansen AC.
\newblock On instabilities of deep learning in image reconstruction and the
  potential costs of {AI}.
\newblock Proceedings of the National Academy of Sciences.
  2020;117(48):30088-95.
\newblock \url{https://doi.org/10.1073/pnas.1907377117}.

\bibitem{bhadra2021hallucinations}
Bhadra S, Kelkar VA, Brooks FJ, Anastasio MA.
\newblock On Hallucinations in Tomographic Image Reconstruction.
\newblock {IEEE} Transactions on Medical Imaging. 2021;40(11):3249-60.
\newblock \url{https://doi.org/10.1109/TMI.2021.3077857}.

\bibitem{tivnan2024hallucination}
Tivnan M, Yoon S, Chen Z, Li X, Wu D, Li Q.
\newblock Hallucination Index: An Image Quality Metric for Generative
  Reconstruction Models.
\newblock In: Medical Image Computing and Computer Assisted Intervention --
  {MICCAI} 2024. Lecture Notes in Computer Science. Springer; 2024. p. 449-58.
\newblock Available from: \url{https://doi.org/10.1007/978-3-031-72117-5_42}.

\bibitem{varoquaux2022failures}
Varoquaux G, Cheplygina V.
\newblock Machine learning for medical imaging: methodological failures and
  recommendations for the future.
\newblock npj Digital Medicine. 2022;5:48.
\newblock \url{https://doi.org/10.1038/s41746-022-00592-y}.

\bibitem{simmons2011falsepositive}
Simmons JP, Nelson LD, Simonsohn U.
\newblock False-Positive Psychology: Undisclosed Flexibility in Data Collection
  and Analysis Allows Presenting Anything as Significant.
\newblock Psychological Science. 2011;22(11):1359-66.
\newblock \url{https://doi.org/10.1177/0956797611417632}.

\bibitem{suzuki2006mtann}
Suzuki K, Abe H, MacMahon H, Doi K.
\newblock Image-processing technique for suppressing ribs in chest radiographs
  by means of massive training artificial neural network ({MTANN}).
\newblock {IEEE} Transactions on Medical Imaging. 2006;25(4):406-16.
\newblock \url{https://doi.org/10.1109/TMI.2006.871549}.

\bibitem{yang2017cascade}
Yang W, Chen Y, Liu Y, Zhong L, Qin G, Lu Z, et~al.
\newblock Cascade of multi-scale convolutional neural networks for bone
  suppression of chest radiographs in gradient domain.
\newblock Medical Image Analysis. 2017;35:421-33.
\newblock \url{https://doi.org/10.1016/j.media.2016.08.004}.

\bibitem{woo2023convnextv2}
Woo S, Debnath S, Hu R, Chen X, Liu Z, Kweon IS, et~al.
\newblock {ConvNeXt V2}: Co-designing and Scaling {ConvNets} with Masked
  Autoencoders.
\newblock In: Proceedings of the {IEEE/CVF} Conference on Computer Vision and
  Pattern Recognition ({CVPR}); 2023. p. 16133-42.
\newblock \url{https://doi.org/10.1109/CVPR52729.2023.01548}.

\bibitem{liu2022swinv2}
Liu Z, Hu H, Lin Y, Yao Z, Xie Z, Wei Y, et~al.
\newblock Swin Transformer {V2}: Scaling Up Capacity and Resolution.
\newblock In: Proceedings of the {IEEE/CVF} Conference on Computer Vision and
  Pattern Recognition ({CVPR}); 2022. p. 11999-2009.
\newblock \url{https://doi.org/10.1109/CVPR52688.2022.01170}.

\bibitem{cheng2020addressing}
Cheng K, Caliv{\'a} F, Shah R, Han M, Majumdar S, Pedoia V.
\newblock Addressing The False Negative Problem of Deep Learning {MRI}
  Reconstruction Models by Adversarial Attacks and Robust Training.
\newblock In: Proceedings of the Third Conference on Medical Imaging with Deep
  Learning. vol. 121 of Proceedings of Machine Learning Research. PMLR; 2020.
  p. 121-35.
\newblock Available from:
  \url{https://proceedings.mlr.press/v121/cheng20a.html}.

\bibitem{li2021assessing}
Li K, Zhou W, Li H, Anastasio MA.
\newblock Assessing the Impact of Deep Neural Network-Based Image Denoising on
  Binary Signal Detection Tasks.
\newblock {IEEE} Transactions on Medical Imaging. 2021;40(9):2295-305.
\newblock \url{https://doi.org/10.1109/TMI.2021.3076810}.

\bibitem{liu2021query2label}
Liu S, Zhang L, Yang X, Su H, Zhu J. {Query2Label}: A Simple Transformer Way to
  Multi-Label Classification; 2021.
\newblock arXiv preprint arXiv:2107.10834.
\newblock \url{https://doi.org/10.48550/arXiv.2107.10834}.

\bibitem{ridnik2023mldecoder}
Ridnik T, Sharir G, Ben-Cohen A, Ben-Baruch E, Noy A.
\newblock {ML-Decoder}: Scalable and Versatile Classification Head.
\newblock In: 2023 IEEE/CVF Winter Conference on Applications of Computer
  Vision (WACV); 2023. p. 32-41.
\newblock \url{https://doi.org/10.1109/WACV56688.2023.00012}.

\bibitem{chattopadhay2018gradcampp}
Chattopadhay A, Sarkar A, Howlader P, Balasubramanian VN.
\newblock {Grad-CAM++}: Generalized Gradient-Based Visual Explanations for Deep
  Convolutional Networks.
\newblock In: 2018 IEEE Winter Conference on Applications of Computer Vision
  (WACV); 2018. p. 839-47.
\newblock \url{https://doi.org/10.1109/WACV.2018.00097}.

\bibitem{sundararajan2017axiomatic}
Sundararajan M, Taly A, Yan Q. Axiomatic Attribution for Deep Networks; 2017.
\newblock arXiv preprint arXiv:1703.01365.
\newblock \url{https://doi.org/10.48550/arXiv.1703.01365}.

\bibitem{zhang2018topdown}
Zhang J, Bargal SA, Lin Z, Brandt J, Shen X, Sclaroff S.
\newblock Top-Down Neural Attention by Excitation Backprop.
\newblock International Journal of Computer Vision. 2018;126(10):1084-102.
\newblock \url{https://doi.org/10.1007/s11263-017-1059-x}.

\bibitem{lakshminarayanan2017ensembles}
Lakshminarayanan B, Pritzel A, Blundell C.
\newblock Simple and Scalable Predictive Uncertainty Estimation using Deep
  Ensembles.
\newblock In: Advances in Neural Information Processing Systems. vol.~30; 2017.
  Available from:
  \url{https://proceedings.neurips.cc/paper_files/paper/2017/hash/9ef2ed4b7fd2c810847ffa5fa85bce38-Abstract.html}.

\bibitem{lei2018conformal}
Lei J, G'Sell M, Rinaldo A, Tibshirani RJ, Wasserman L.
\newblock Distribution-Free Predictive Inference for Regression.
\newblock Journal of the American Statistical Association.
  2018;113(523):1094-111.
\newblock \url{https://doi.org/10.1080/01621459.2017.1307116}.

\bibitem{angelopoulos2023conformal}
Angelopoulos AN, Bates S.
\newblock Conformal Prediction: A Gentle Introduction.
\newblock Foundations and Trends in Machine Learning. 2023;16(4):494-591.
\newblock Preprint title: A Gentle Introduction to Conformal Prediction and
  Distribution-Free Uncertainty Quantification.
\newblock \url{https://doi.org/10.1561/2200000101}.

\bibitem{angelopoulos2022crc}
Angelopoulos AN, Bates S, Fisch A, Lei L, Schuster T. Conformal Risk Control;
  2022.
\newblock arXiv preprint arXiv:2208.02814.
\newblock \url{https://doi.org/10.48550/arXiv.2208.02814}.

\bibitem{vovk2005algorithmic}
Vovk V, Gammerman A, Shafer G.
\newblock Algorithmic Learning in a Random World.
\newblock Springer; 2005.
\newblock \url{https://doi.org/10.1007/b106715}.

\bibitem{geifman2017selective}
Geifman Y, El-Yaniv R. Selective Classification for Deep Neural Networks; 2017.
\newblock arXiv preprint arXiv:1705.08500.
\newblock \url{https://doi.org/10.48550/arXiv.1705.08500}.

\bibitem{geifman2019bias}
Geifman Y, Uziel G, El-Yaniv R. Bias-Reduced Uncertainty Estimation for Deep
  Neural Classifiers; 2019.
\newblock International Conference on Learning Representations (ICLR); arXiv
  preprint arXiv:1805.08206.
\newblock \url{https://doi.org/10.48550/arXiv.1805.08206}.

\bibitem{armato2011lidc}
Armato SG III, McLennan G, Bidaut L, McNitt-Gray MF, Meyer CR, Reeves AP,
  et~al.
\newblock The {Lung Image Database Consortium (LIDC)} and {Image Database
  Resource Initiative (IDRI)}: A completed reference database of lung nodules
  on {CT} scans.
\newblock Medical Physics. 2011;38(2):915-31.
\newblock \url{https://doi.org/10.1118/1.3528204}.

\bibitem{armato2015lidctcia}
Armato SG III, et~al. Data From {LIDC-IDRI}; 2015.
\newblock Licence: CC BY 3.0.
\newblock The Cancer Imaging Archive.
\newblock \url{https://doi.org/10.7937/K9/TCIA.2015.LO9QL9SX}.

\bibitem{clark2013tcia}
Clark K, Vendt B, Smith K, Freymann J, Kirby J, Koppel P, et~al.
\newblock The {Cancer Imaging Archive (TCIA)}: Maintaining and Operating a
  Public Information Repository.
\newblock Journal of Digital Imaging. 2013;26(6):1045-57.
\newblock \url{https://doi.org/10.1007/s10278-013-9622-7}.

\bibitem{setio2017luna16}
Setio AAA, Traverso A, de~Bel T, Berens MSN, van~den Bogaard C, Cerello P,
  et~al.
\newblock Validation, comparison, and combination of algorithms for automatic
  detection of pulmonary nodules in computed tomography images: The {LUNA16}
  challenge.
\newblock Medical Image Analysis. 2017;42:1-13.
\newblock \url{https://doi.org/10.1016/j.media.2017.06.015}.

\bibitem{nguyen2021ribcxr}
Nguyen HC, Le TT, Pham HH, Nguyen HQ.
\newblock {VinDr-RibCXR}: A Benchmark Dataset for Automatic Segmentation and
  Labeling of Individual Ribs on Chest {X}-rays.
\newblock In: Medical Imaging with Deep Learning ({MIDL}); 2021. Short paper;
  arXiv:2107.01327.

\bibitem{shiraishi2000jsrt}
Shiraishi J, Katsuragawa S, Ikezoe J, Matsumoto T, Kobayashi T, Komatsu Ki,
  et~al.
\newblock Development of a Digital Image Database for Chest Radiographs With
  and Without a Lung Nodule: Receiver Operating Characteristic Analysis of
  Radiologists' Detection of Pulmonary Nodules.
\newblock American Journal of Roentgenology. 2000;174(1):71-4.
\newblock \url{https://doi.org/10.2214/ajr.174.1.1740071}.

\bibitem{sogancioglu2024node21}
Sogancioglu E, van Ginneken B, Behrendt F, Bengs M, Schlaefer A, Radu M, et~al.
\newblock Nodule Detection and Generation on Chest {X}-Rays: {NODE21}
  Challenge.
\newblock {IEEE} Transactions on Medical Imaging. 2024;43(8):2839-53.
\newblock \url{https://doi.org/10.1109/TMI.2024.3382042}.

\bibitem{liu2020tbx11k}
Liu Y, Wu YH, Ban Y, Wang H, Cheng MM.
\newblock Rethinking Computer-Aided Tuberculosis Diagnosis.
\newblock In: Proceedings of the {IEEE/CVF} Conference on Computer Vision and
  Pattern Recognition ({CVPR}); 2020. p. 2643-52.
\newblock \url{https://doi.org/10.1109/CVPR42600.2020.00272}.

\bibitem{nguyen2022vindrcxr}
Nguyen HQ, Lam K, Le LT, Pham HH, Tran DQ, Nguyen DB, et~al.
\newblock {VinDr-CXR}: An open dataset of chest {X}-rays with radiologist's
  annotations.
\newblock Scientific Data. 2022;9:429.
\newblock \url{https://doi.org/10.1038/s41597-022-01498-w}.

\bibitem{vinbigdata2020kaggle}
Nguyen D, {DungNB}, Nguyen HQ, Elliott J, {NguyenThanhNhan}, Culliton P.
  {VinBigData} Chest {X}-ray Abnormalities Detection; 2020.
\newblock Accessed 2026-11-03.
\newblock Kaggle competition.
\newblock Available from:
  \url{https://www.kaggle.com/competitions/vinbigdata-chest-xray-abnormalities-detection}.

\bibitem{wang2017chestxray}
Wang X, Peng Y, Lu L, Lu Z, Bagheri M, Summers RM.
\newblock {ChestX-ray8}: Hospital-scale Chest {X}-ray Database and Benchmarks
  on Weakly-Supervised Classification and Localization of Common Thorax
  Diseases.
\newblock In: {IEEE} Conference on Computer Vision and Pattern Recognition
  ({CVPR}); 2017. p. 3462-71.
\newblock \url{https://doi.org/10.1109/CVPR.2017.369}.

\bibitem{durjoy2026bdcxr}
Durjoy SH, Shikder ME, Bijoy MHI, Shoib MMH, Rahman A, Hasan MZ, et~al.
  {BDCXR-3257}: A Comprehensive Chest Radiography Dataset with
  Multi-Etiological Pneumonia Labels from a Clinical Setting; 2026.
\newblock Version 1; licence CC BY 4.0.
\newblock figshare, dataset.
\newblock \url{https://doi.org/10.6084/m9.figshare.32086362.v1}.

\bibitem{kermany2018identifying}
Kermany DS, Goldbaum M, Cai W, Valentim CCS, Liang H, Baxter SL, et~al.
\newblock Identifying Medical Diagnoses and Treatable Diseases by Image-Based
  Deep Learning.
\newblock Cell. 2018;172(5):1122-31.e9.
\newblock \url{https://doi.org/10.1016/j.cell.2018.02.010}.

\bibitem{zhou2018unetpp}
Zhou Z, Rahman~Siddiquee MM, Tajbakhsh N, Liang J.
\newblock {UNet++}: A Nested {U-Net} Architecture for Medical Image
  Segmentation.
\newblock In: Deep Learning in Medical Image Analysis and Multimodal Learning
  for Clinical Decision Support ({DLMIA} 2018). vol. 11045 of Lecture Notes in
  Computer Science. Springer; 2018. p. 3-11.
\newblock Available from: \url{https://doi.org/10.1007/978-3-030-00889-5_1}.

\bibitem{woo2018cbam}
Woo S, Park J, Lee JY, Kweon IS.
\newblock {CBAM}: Convolutional Block Attention Module.
\newblock In: Computer Vision -- {ECCV} 2018. Lecture Notes in Computer
  Science. Springer; 2018. p. 3-19.
\newblock Available from: \url{https://doi.org/10.1007/978-3-030-01234-2_1}.

\bibitem{deng2009imagenet}
Deng J, Dong W, Socher R, Li LJ, Li K, Fei-Fei L.
\newblock {ImageNet}: A large-scale hierarchical image database.
\newblock In: {IEEE} Conference on Computer Vision and Pattern Recognition
  ({CVPR}); 2009. p. 248-55.
\newblock \url{https://doi.org/10.1109/CVPR.2009.5206848}.

\bibitem{wightman2019timm}
Wightman R. {PyTorch} Image Models. GitHub; 2019.
\newblock \url{https://github.com/rwightman/pytorch-image-models}.
\newblock \url{https://doi.org/10.5281/zenodo.4414861}.

\bibitem{loshchilov2019adamw}
Loshchilov I, Hutter F.
\newblock Decoupled Weight Decay Regularization.
\newblock In: International Conference on Learning Representations ({ICLR});
  2019. Available from: \url{https://arxiv.org/abs/1711.05101}.

\bibitem{simeoni2025dinov3}
Sim{\'e}oni O, Vo HV, Seitzer M, Baldassarre F, Oquab M, Jose C, et~al.
  {DINOv3}; 2025.
\newblock arXiv preprint arXiv:2508.10104.
\newblock \url{https://doi.org/10.48550/arXiv.2508.10104}.

\bibitem{perezgarcia2025raddino}
P{\'e}rez-Garc{\'\i}a F, Sharma H, Bond-Taylor S, Bouzid K, Salvatelli V, Ilse
  M, et~al.
\newblock Exploring scalable medical image encoders beyond text supervision.
\newblock Nature Machine Intelligence. 2025;7(1):119-30.
\newblock \url{https://doi.org/10.1038/s42256-024-00965-w}.

\bibitem{zuiderveld1994clahe}
Zuiderveld K.
\newblock Contrast Limited Adaptive Histogram Equalization.
\newblock In: Heckbert PS, editor. Graphics Gems {IV}. Boston: AP Professional;
  1994. p. 474-85.
\newblock \url{https://doi.org/10.1016/B978-0-12-336156-1.50061-6}.

\bibitem{oppenheim1968homomorphic}
Oppenheim AV, Schafer RW, Stockham TG.
\newblock Nonlinear filtering of multiplied and convolved signals.
\newblock Proceedings of the {IEEE}. 1968;56(8):1264-91.
\newblock \url{https://doi.org/10.1109/PROC.1968.6570}.

\bibitem{cohen2022torchxrayvision}
Cohen JP, Viviano JD, Bertin P, Morrison P, Torabian P, Guarrera M, et~al.
\newblock {TorchXRayVision}: A library of chest {X}-ray datasets and models.
\newblock In: Konukoglu E, Menze B, Venkataraman A, Baumgartner C, Dou Q,
  Albarqouni S, editors. Proceedings of The 5th International Conference on
  Medical Imaging with Deep Learning. vol. 172 of Proceedings of Machine
  Learning Research. PMLR; 2022. p. 231-49.

\bibitem{lian2021chestxdet}
Lian J, Liu J, Zhang S, Gao K, Liu X, Zhang D, et~al.
\newblock A Structure-Aware Relation Network for Thoracic Diseases Detection
  and Segmentation.
\newblock {IEEE} Transactions on Medical Imaging. 2021;40(8):2042-52.
\newblock \url{https://doi.org/10.1109/TMI.2021.3070847}.

\bibitem{wang2004ssim}
Wang Z, Bovik AC, Sheikh HR, Simoncelli EP.
\newblock Image quality assessment: from error visibility to structural
  similarity.
\newblock {IEEE} Transactions on Image Processing. 2004;13(4):600-12.
\newblock \url{https://doi.org/10.1109/TIP.2003.819861}.

\bibitem{wang2003msssim}
Wang Z, Simoncelli EP, Bovik AC.
\newblock Multiscale structural similarity for image quality assessment.
\newblock In: The Thirty-Seventh Asilomar Conference on Signals, Systems \&
  Computers. vol.~2; 2003. p. 1398-402.
\newblock \url{https://doi.org/10.1109/ACSSC.2003.1292216}.

\bibitem{zhang2018lpips}
Zhang R, Isola P, Efros AA, Shechtman E, Wang O.
\newblock The Unreasonable Effectiveness of Deep Features as a Perceptual
  Metric.
\newblock In: Proceedings of the {IEEE/CVF} Conference on Computer Vision and
  Pattern Recognition ({CVPR}); 2018. p. 586-95.
\newblock \url{https://doi.org/10.1109/CVPR.2018.00068}.

\bibitem{huang2017densenet}
Huang G, Liu Z, van~der Maaten L, Weinberger KQ.
\newblock Densely Connected Convolutional Networks.
\newblock In: {IEEE} Conference on Computer Vision and Pattern Recognition
  ({CVPR}); 2017. p. 2261-9.
\newblock \url{https://doi.org/10.1109/CVPR.2017.243}.

\bibitem{ren2017fasterrcnn}
Ren S, He K, Girshick R, Sun J.
\newblock Faster {R-CNN}: Towards Real-Time Object Detection with Region
  Proposal Networks.
\newblock {IEEE} Transactions on Pattern Analysis and Machine Intelligence.
  2017;39(6):1137-49.
\newblock \url{https://doi.org/10.1109/TPAMI.2016.2577031}.

\bibitem{vanginneken2010anode09}
van Ginneken B, Armato SG, de~Hoop B, van Amelsvoort-van~de Vorst S, Duindam T,
  Niemeijer M, et~al.
\newblock Comparing and combining algorithms for computer-aided detection of
  pulmonary nodules in computed tomography scans: The {ANODE09} study.
\newblock Medical Image Analysis. 2010;14(6):707-22.
\newblock \url{https://doi.org/10.1016/j.media.2010.05.005}.

\bibitem{niemeijer2011cpm}
Niemeijer M, Loog M, Abr{\`a}moff MD, Viergever MA, Prokop M, van Ginneken B.
\newblock On Combining Computer-Aided Detection Systems.
\newblock {IEEE} Transactions on Medical Imaging. 2011;30(2):215-23.
\newblock \url{https://doi.org/10.1109/TMI.2010.2072789}.

\bibitem{efron1993bootstrap}
Efron B, Tibshirani RJ.
\newblock An Introduction to the Bootstrap.
\newblock Boston, MA: Springer US; 1993.
\newblock \url{https://doi.org/10.1007/978-1-4899-4541-9}.

\bibitem{wilcoxon1945individual}
Wilcoxon F.
\newblock Individual Comparisons by Ranking Methods.
\newblock Biometrics Bulletin. 1945;1(6):80-3.
\newblock \url{https://doi.org/10.2307/3001968}.

\bibitem{delong1988auc}
DeLong ER, DeLong DM, Clarke-Pearson DL.
\newblock Comparing the Areas under Two or More Correlated Receiver Operating
  Characteristic Curves: A Nonparametric Approach.
\newblock Biometrics. 1988;44(3):837-45.
\newblock \url{https://doi.org/10.2307/2531595}.

\bibitem{holm1979simple}
Holm S.
\newblock A Simple Sequentially Rejective Multiple Test Procedure.
\newblock Scandinavian Journal of Statistics. 1979;6(2):65-70.
\newblock Available from: \url{https://www.jstor.org/stable/4615733}.

\bibitem{icru1989report44}
{International Commission on Radiation Units and Measurements}.
\newblock Tissue Substitutes in Radiation Dosimetry and Measurement.
\newblock International Commission on Radiation Units and Measurements; 1989.
  44.
\newblock Journal of the ICRU, vol. os23, no. 1.
\newblock \url{https://doi.org/10.1093/jicru/os23.1.Report44}.

\bibitem{hubbell2004xraymass}
Hubbell JH, Seltzer SM. Tables of {X}-Ray Mass Attenuation Coefficients and
  Mass Energy-Absorption Coefficients (version 1.4); 2004.
\newblock NIST Standard Reference Database 126; originally published as NISTIR
  5632 (1995). Accessed 2026-11-03.
\newblock National Institute of Standards and Technology, Gaithersburg, MD.
  \url{http://physics.nist.gov/xaamdi}.
\newblock \url{https://doi.org/10.18434/T4D01F}.

\bibitem{poludniowski2021spekpy}
Poludniowski GG, Omar A, Bujila R, Andreo P.
\newblock Technical Note: {SpekPy} v2.0---a software toolkit for modeling x-ray
  tube spectra.
\newblock Medical Physics. 2021;48(7):3630-7.
\newblock \url{https://doi.org/10.1002/mp.14945}.

\bibitem{tejani2024claim}
Tejani AS, Klontzas ME, Gatti AA, Mongan JT, Moy L, Park SH, et~al.
\newblock Checklist for Artificial Intelligence in Medical Imaging ({CLAIM}):
  2024 Update.
\newblock Radiology: Artificial Intelligence. 2024;6(4):e240300.
\newblock \url{https://doi.org/10.1148/ryai.240300}.

\end{thebibliography}


\end{document}
